% !TEX program = xelatex
% 考研英语二 作文模板总纲（整合·改进版）
% 编译：xelatex 两遍
\documentclass[11pt,a4paper]{ctexart}

\usepackage{geometry}
\geometry{a4paper,top=2.1cm,bottom=2.0cm,left=2.0cm,right=2.0cm,headsep=6mm}

\usepackage{fontspec}
\setmainfont{Times New Roman}
\setsansfont{Arial}

\usepackage{xcolor}
\usepackage{amsmath,amssymb}
\usepackage{booktabs}
\usepackage{longtable}
\usepackage{array}
\usepackage{tabularx}
\usepackage{enumitem}
\usepackage{needspace}
\usepackage{fancyhdr}
\usepackage{tikz}
\usetikzlibrary{positioning,arrows.meta,calc}
\usepackage[hidelinks,bookmarksnumbered]{hyperref}
\hypersetup{
  pdftitle={考研英语二 作文模板总纲（整合·改进版）},
  pdfsubject={Part A 应用文（约 100 词 / 10 分）+ Part B 图表作文（约 150 词 / 15 分）},
  pdfkeywords={考研英语二, 作文模板, 图表作文, 应用文, 评分标准, 好词好句}
}

% ---------- 颜色（中性灰阶，点缀低饱和蓝灰） ----------
\definecolor{ink}{gray}{0.10}
\definecolor{mid}{gray}{0.42}
\definecolor{light}{gray}{0.72}
\definecolor{band}{gray}{0.955}
\definecolor{accent}{RGB}{42,74,105}

\color{ink}

% ---------- 版面 ----------
\linespread{1.22}
\setlength{\parindent}{0pt}
\setlength{\parskip}{4pt}

\ctexset{
  section/format = {\Large\bfseries\color{accent}},
  section/name   = {},
  section/number = \arabic{section},
  subsection/format = {\large\bfseries\color{ink}},
  subsubsection/format = {\normalsize\bfseries\color{ink}},
  section/beforeskip = 14pt plus 2pt minus 2pt,
  section/afterskip  = 8pt,
  subsection/beforeskip = 12pt plus 2pt minus 2pt,
  subsection/afterskip  = 6pt,
}

\pagestyle{fancy}
\fancyhf{}
\renewcommand{\headrulewidth}{0.3pt}
\renewcommand{\footrulewidth}{0pt}
\fancyhead[L]{\footnotesize\color{mid}考研英语二 作文模板总纲}
\fancyhead[R]{\footnotesize\color{mid}\leftmark}
\fancyfoot[C]{\footnotesize\color{mid}\thepage}

\renewcommand{\sectionmark}[1]{\markboth{#1}{}}

% ---------- 表格 ----------
\renewcommand{\arraystretch}{1.25}
\setlength{\tabcolsep}{5pt}

% ---------- 排版宏 ----------
% 填空占位
\newcommand{\fillin}[1]{{\bfseries【#1】}}
\newcommand{\en}[1]{\textit{#1}}
\newcommand{\kw}[1]{{\bfseries\color{accent}#1}}

% 雷区 / 对标评分 / 改进 标记
\newcommand{\zone}[1]{\par\addvspace{3pt}\noindent{\bfseries\color{accent}【雷区】}#1\par\addvspace{3pt}}
\newcommand{\rubric}[1]{\par\addvspace{3pt}\noindent{\bfseries【对标评分】}#1\par\addvspace{3pt}}
\newcommand{\impr}[1]{\par\addvspace{3pt}\noindent{\bfseries【本次改进】}#1\par\addvspace{3pt}}
\newcommand{\tip}[1]{\par\addvspace{3pt}\noindent{\bfseries【提示】}#1\par\addvspace{3pt}}

% 模板块：标题 + 上下细线，可跨页
\newenvironment{tpl}[1]{%
  \par\addvspace{7pt}\needspace{4\baselineskip}%
  \noindent{\bfseries\color{accent}#1}\par\nopagebreak
  \vspace{2pt}\noindent\rule{\linewidth}{0.7pt}\par\nopagebreak\vspace{3pt}%
  \begin{list}{}{\leftmargin=1.05em \rightmargin=0.25em \listparindent=0pt
    \itemindent=0pt \topsep=0pt \parsep=2.5pt \itemsep=2.5pt \partopsep=0pt}%
  \item\relax
}{%
  \end{list}\vspace{1pt}\noindent\rule{\linewidth}{0.7pt}\par\addvspace{8pt}%
}

% 说明块：左侧竖线缩进的小字
\newenvironment{note}{%
  \par\addvspace{5pt}%
  \begin{list}{}{\leftmargin=1.05em \rightmargin=0.25em \listparindent=0pt
    \itemindent=0pt \topsep=0pt \parsep=2pt \itemsep=2pt \partopsep=0pt}%
  \item\relax\small\color{ink}%
}{%
  \end{list}\addvspace{5pt}%
}

% 例句块
\newenvironment{eg}{%
  \par\addvspace{4pt}%
  \begin{list}{}{\leftmargin=1.4em \rightmargin=0.2em \listparindent=0pt
    \itemindent=0pt \topsep=0pt \parsep=2pt \itemsep=2pt \partopsep=0pt}%
  \item\relax
}{%
  \end{list}\addvspace{5pt}%
}

% 用于长表格的居中窄列
\newcolumntype{L}[1]{>{\raggedright\arraybackslash}p{#1}}
\newcolumntype{C}[1]{>{\centering\arraybackslash}p{#1}}

% 勾选框
\newcommand{\cb}{$\square$}

% 带圈数字 ①②③ 等：Times New Roman 无此字形，交给中文字体渲染
\xeCJKDeclareCharClass{CJK}{"2460 -> "24FF}
\xeCJKDeclareCharClass{CJK}{"25A0 -> "25FF}

\begin{document}

% ==================== 封面 ====================
\begin{titlepage}
\centering
\vspace*{2.2cm}

{\fontsize{30}{38}\selectfont\bfseries\color{accent} 考研英语二}\\[10pt]
{\fontsize{30}{38}\selectfont\bfseries\color{accent} 作文模板总纲}\\[18pt]

{\large\bfseries 整合 · 纠错 · 改进版}\\[6pt]
{\normalsize\color{mid} ——严格对标《英语二》阅卷评分档位与批改原则——}

\vspace{1.4cm}
\noindent\rule{\linewidth}{1pt}
\vspace{0.6cm}

\begin{minipage}{0.86\linewidth}
\centering
\large
\textbf{Part A 应用文} \quad 约 100 词 \quad 满分 10 分\\[4pt]
\textbf{Part B 图表作文} \quad 约 150 词 \quad 满分 15 分\\[10pt]
\normalsize\color{mid}
（英语二 Part B 自 2010 年分卷以来只考图表；本册主攻图表与应用文）
\end{minipage}

\vspace{1.0cm}
\noindent\rule{\linewidth}{1pt}
\vspace{0.8cm}

\begin{minipage}{0.88\linewidth}
\small
\textbf{本册材料来源}\\[4pt]
1.\ 《mmg 作文模板》（考研的秘密）——大作文静态/动态第一段、五类话题语料、小作文功能句库、话题词库；\\
2.\ 《英语作文模板》（御风 yufeng）——图表审题与写法、动态描述范式、议论句法、逻辑词汇总、真题范文；\\
3.\ 《考研英语批改指令》（英语二部分）——本册全部评分标准、档位描述与批改原则的唯一依据。\\[8pt]
\textbf{整合原则}\\[4pt]
源材料里的每一个好词、好句、好框架都已被收入；重复内容合并，错误内容更正，与评分规则冲突的内容改造；新增图表锚定写法、自查清单与真题批注。
\end{minipage}

\vfill
{\small\color{mid} 生成日期：\today}
\end{titlepage}

% ==================== 目录 ====================
\pagenumbering{roman}
\renewcommand{\contentsname}{目\quad 录}
\tableofcontents
\clearpage
\pagenumbering{arabic}

% ==================== 全书结构图 ====================
\section*{全书结构图}
\addcontentsline{toc}{section}{全书结构图}
\markboth{全书结构图}{}

\begin{center}
\begin{tikzpicture}[
  font=\small,
  box/.style={draw=ink,line width=0.6pt,rounded corners=1.5pt,align=left,
              inner sep=5pt,fill=white,text width=3.05cm},
  root/.style={draw=accent,line width=1pt,rounded corners=2pt,align=center,
               inner sep=7pt,fill=band,text width=3.3cm,font=\small\bfseries},
  leaf/.style={draw=light,line width=0.5pt,rounded corners=1.5pt,align=left,
               inner sep=4pt,fill=white,text width=4.35cm,font=\footnotesize},
  lbl/.style={font=\scriptsize\bfseries,color=accent,align=left,text width=2.5cm},
  ln/.style={draw=mid,line width=0.5pt},
]
\node[root] (R) at (0,0) {考研英语二写作\\25 分};

\node[box] (A) at (4.4,3.0) {\textbf{Part A 应用文}\\约 100 词 / 10 分\\书信、通知、告示、纪要};
\node[box] (B) at (4.4,-2.9) {\textbf{Part B 图表作文}\\约 150 词 / 15 分\\柱状、饼状、折线、表格、双图};

\node[leaf] (A1) at (9.9,4.6) {\textbf{第 1 段（2 句）}\\问候 + 点明写作目的};
\node[leaf] (A2) at (9.9,3.1) {\textbf{第 2 段（3--5 句）}\\事件 / 介绍 / 原因\\按题目要点逐条回应};
\node[leaf] (A3) at (9.9,1.6) {\textbf{第 3 段（1--2 句）}\\期待回应 / 期待结果};

\node[leaf] (B1) at (9.9,-1.4) {\textbf{第 1 段（2--3 句）}\\图表信息说明\\对象 + 时间 + 数值};
\node[leaf] (B2) at (9.9,-2.9) {\textbf{第 2 段（4--5 句）}\\描述与解读\\是什么 → 为什么 / 影响};
\node[leaf] (B3) at (9.9,-4.4) {\textbf{第 3 段（3 句）}\\总结 → 概括 → 建议\\不再出现数字};

\draw[ln] (R.east) -- ++(0.35,0) |- (A.west);
\draw[ln] (R.east) -- ++(0.35,0) |- (B.west);
\draw[ln] (A.east) -- ++(0.3,0) |- (A1.west);
\draw[ln] (A.east) -- ++(0.3,0) |- (A2.west);
\draw[ln] (A.east) -- ++(0.3,0) |- (A3.west);
\draw[ln] (B.east) -- ++(0.3,0) |- (B1.west);
\draw[ln] (B.east) -- ++(0.3,0) |- (B2.west);
\draw[ln] (B.east) -- ++(0.3,0) |- (B3.west);

\node[lbl,anchor=east,text width=1.9cm,align=right] (S) at (-1.55,-8.0) {评分只看\\四件事};
\node[box,text width=4.15cm,fill=band] (S1) at (2.05,-8.0)
  {\footnotesize （1）\textbf{回应任务}：要点全覆盖、不跑题\\
   （2）\textbf{用法正确}：拼写、语法、得体地道\\
   （3）\textbf{逻辑流畅}：段落与句间关系清楚\\
   （4）\textbf{丰富多样}：词汇句型不重复};
\node[box,text width=5.6cm,fill=band] (S2) at (9.15,-8.0)
  {\footnotesize【图表题四条红线】\\
   1. 不把「类别差异」写成「时间趋势」\\
   2. 不混淆「百分比」与「百分点」\\
   3. 不编造图中没有的数据\\
   4. 合理推测不能写成图表已证明的事实};
\draw[ln] (S.east) -- (S1.west);
\draw[ln] (S1.east) -- (S2.west);
\end{tikzpicture}
\end{center}

\begin{note}
看这张图只需记住一件事：\textbf{英语二的两篇文章，段落数量、段落功能、每段句数基本是固定的}。考场上的真正工作量，是把【方括号】里的内容按当次题目填对，而不是临场构思结构。
\end{note}

\clearpage

\section{先看规则：英语二评分标准与批改原则}
\markboth{先看规则：英语二评分标准与批改原则}{}

\begin{note}
本册所有模板，都是从下面的档位描述和批改原则**反推**出来的。先花 10 分钟看懂这一章，你对“什么该写、什么绝对不能写”的判断力，胜过背十篇模板。
\end{note}

\subsection{英语二写作的官方要求}

\begin{longtable}{L{2.6cm} L{3.2cm} L{2.6cm} L{2.6cm} L{3.4cm}}
\toprule
\textbf{部分} & \textbf{文体} & \textbf{字数} & \textbf{分值} & \textbf{提示信息形式}\\
\midrule
\endfirsthead
\toprule
\textbf{部分} & \textbf{文体} & \textbf{字数} & \textbf{分值} & \textbf{提示信息形式}\\
\midrule
\endhead
A 节（小作文） & 应用文：书信、通知、告示、纪要等 & 约 100 词 & 10 分 & 规定情景\\
B 节（大作文） & 说明文 / 议论文 & 约 150 词 & 15 分 & 文字、图画、图表等\\
\bottomrule
\end{longtable}

\begin{note}
\textbf{三个必须记住的细节：}
\begin{enumerate}[leftmargin=1.6em,itemsep=1pt,topsep=2pt]
\item 英语二写作合计\textbf{25 分}（英语一是 30 分：小作文 10 + 大作文 20）。全卷 100 分里，写作占四分之一，是投入产出比最高的板块。
\item 大纲写的是\textbf{“约 150 词”}。批改原则明确：“题目要求约 150 词时，不擅自改成必须不少于 150 词。”所以\textbf{不要为了凑字数把文章灌到 220 词}——本册把大作文的目标定在 \textbf{150--175 词}，小作文定在 \textbf{100--120 词}。
\item 提示信息形式虽然列了“文字、图画、图表”，但自 2010 年分卷以来，\textbf{英语二 Part B 只考过图表}（柱状、饼状、折线、表格、双图）。所以本册主体是图表作文。图画只作附录里的最低优先级保险（见附录 B）。
\end{enumerate}
\end{note}

\subsection{英语二大作文评分参考（满分 15 分）}

\begin{longtable}{C{1.5cm} L{6.2cm} L{7.4cm}}
\toprule
\textbf{档位} & \textbf{评分参考原文} & \textbf{这一档在阅卷现场的真正含义}\\
\midrule
\endfirsthead
\toprule
\textbf{档位} & \textbf{评分参考原文} & \textbf{这一档在阅卷现场的真正含义}\\
\midrule
\endhead
\textbf{13--15} & 内容切题，完整覆盖任务要点，图表描述准确，评论相关且充分；语言丰富准确，行文自然，层次清晰。 & 图表读对了（对象、时间、单位、数值、趋势都没错），任务要点一个不漏，评论与图表直接相关，通篇几乎无语言错误，段落层次一眼可辨。\\
\midrule
\textbf{10--12} & 较好完成任务，主要要点完整，允许遗漏少量次要信息；语言较丰富，偶有错误，整体连贯。 & 图表主要特征抓到了，可能漏了一个次要数据或次要评论角度；有个别语法/用词错误但不影响理解。\\
\midrule
\textbf{7--9} & 基本完成任务，涵盖多数要点，描述和评论基本合理；语言存在一些错误，但不影响整体理解。 & 写够了字，但图表描述比较笼统或数字用错，评论与图表关联较弱，语言错误开始变多。\\
\midrule
\textbf{4--6} & 遗漏重要内容或存在理解偏差；语言错误较多，影响理解，组织与衔接不足。 & 把静态图当动态图写、编造数据、或评论完全离题——这是最容易踩的两个坑。\\
\midrule
\textbf{1--3} & 严重偏离任务，遗漏主要内容，语言错误多，难以有效表达意思。 & 跑题，或只写了模板套话没有回应题目。\\
\midrule
\textbf{0} & 未作答、内容完全无关，或没有可评价的英文内容。 & ——\\
\bottomrule
\end{longtable}

\subsection{英语二小作文评分参考（满分 10 分）}

\begin{longtable}{C{1.5cm} L{6.2cm} L{7.4cm}}
\toprule
\textbf{档位} & \textbf{评分参考原文} & \textbf{这一档在阅卷现场的真正含义}\\
\midrule
\endfirsthead
\toprule
\textbf{档位} & \textbf{评分参考原文} & \textbf{这一档在阅卷现场的真正含义}\\
\midrule
\endhead
\textbf{8--9 的上限} & 原文把 10 分写在这一档里：“各方面表现充分、无明显不足时可给 10 分，\textbf{不人为封顶}”。\textbf{10 分不是一个独立档位，而是 8--9 档的满分说明。} & 任务要点全覆盖，语言准确自然且有变化，格式语气恰当。满分是真实存在的，不要自己先放弃。\\
\midrule
\textbf{8--9} & 很好地完成任务，要点完整，语言准确自然，表达有变化，行文连贯，格式和语气恰当。 & 该说的都说了，称呼/结束语/署名规范，语气与双方关系匹配。\\
\midrule
\textbf{6--7} & 较好完成任务，主要要点完整，允许遗漏少量次要信息；语言基本准确，偶有错误，组织清晰，格式和语气适当。 & 漏了一个次要要点，或称谓/落款有小瑕疵。\\
\midrule
\textbf{4--5} & 基本完成任务，涵盖多数要点；有一些语言错误，但不影响整体理解，结构和格式基本合理。 & 要点回应不全，或格式明显不对（如通知没写标题/落款）。\\
\midrule
\textbf{2--3} & 遗漏重要要点或夹杂无关内容；语言错误较多，影响理解，衔接不足，格式或语气不当。 & 写成大作文式的空泛议论，没有回应题干要点。\\
\midrule
\textbf{1} & 未完成任务，明显遗漏主要内容，语言错误严重，难以有效传达信息。 & ——\\
\midrule
\textbf{0} & 未作答、内容完全无关，或没有可评价的英文内容。 & ——\\
\bottomrule
\end{longtable}

\subsection{从批改原则反推的九条硬规则}

下面每一条都是阅卷指令里的原话（或原话的直接推论）。左边是规则，右边是它对写作的具体约束。

\subsubsection{大作文（图表题）的五条}

\begin{longtable}{C{0.9cm} L{7.0cm} L{7.6cm}}
\toprule
\textbf{序号} & \textbf{批改原则（原意）} & \textbf{写作时必须做到}\\
\midrule
\endfirsthead
\toprule
\textbf{序号} & \textbf{批改原则（原意）} & \textbf{写作时必须做到}\\
\midrule
\endhead
\textbf{1} & 先核对题目要求，再结合内容、语言和结构综合定档、给分，\textbf{不机械按错误数量扣分}。 & 不要因为怕写错就用最简单的句子；也不要以为堆砌复杂句就能加分。\\
\midrule
\textbf{2} & 图表题重点检查\textbf{对象、时间、单位、数值、比例、排序及变化趋势}。\textbf{区分静态对比与动态变化}，不把类别差异写成时间趋势，不混淆百分比与百分点，也\textbf{不编造数据}。 & 动笔前必须先判断“静态图 / 动态图”；每个数字都来自图表；“占 30\%”和“上升了 30 个百分点”是两回事。\\
\midrule
\textbf{3} & 描述应\textbf{突出主要特征和关键对比}，\textbf{不要求罗列全部数字}。评论须\textbf{与图表相关、逻辑合理}；\textbf{合理推测不能写成图表已经证明的事实}。不强求深入分析原因。 & 只挑 2--4 个最关键的数据；评论句必须有 “may / suggest / indicate” 之类的分寸感，不要写成 “The chart proves that...”。\\
\midrule
\textbf{4} & 统计英文词数，按题目要求评价长度；题目要求约 150 词时，\textbf{不擅自改成必须不少于 150 词}，也不编造固定扣分规则。 & 目标 150--175 词。写满 200+ 词不会加分，只会挤占阅读和检查时间。\\
\midrule
\textbf{5} & \textbf{正确、自然的简单表达可以保留}，不以生僻词、复杂句或连接词数量决定分数，\textbf{不因使用常见句式直接扣分}。 & 用你 100\% 写得对的句式。一个用对的一般句，胜过三个用错的“高级句”。\\
\bottomrule
\end{longtable}

\begin{tpl}{【雷区】图表作文的四条红线（踩一条，直接掉出 13--15 档）}
\textbf{红线一：把“类别差异”写成“时间趋势”。}\\
错：\en{The proportion of elderly people increased, while young people decreased.}（如果图里只有同一年的不同年龄段数据，就没有任何“increase/decrease”，只有“高低之差”。）\\
对：\en{The proportion of elderly respondents was markedly higher than that of younger ones.}\\[3pt]
\textbf{红线二：混淆“百分比”与“百分点”。}\\
错：\en{The figure rose by 33\% from 25\% to 58\%.}（25→58 是涨了 33 个\textbf{百分点}，不是 33\%。）\\
对：\en{The figure rose by 33 percentage points, from 25\% to 58\%.} / \en{...more than doubled, from 25\% to 58\%.}\\[3pt]
\textbf{红线三：编造图中没有的数据。}\\
把模板里的数字当成“待填的装饰”，随手写个 60\%，或者把两个类别的数字张冠李戴。阅卷老师手上有图，一眼就能看出来。\\
\textbf{红线四：把推测写成图表已证明的事实。}\\
错：\en{The chart proves that young people are lazy.}\\
对：\en{This may suggest that ... / One possible explanation is that ...}
\end{tpl}

\subsubsection{小作文（应用文）的四条}

\begin{longtable}{C{0.9cm} L{7.0cm} L{7.6cm}}
\toprule
\textbf{序号} & \textbf{批改原则（原意）} & \textbf{写作时必须做到}\\
\midrule
\endfirsthead
\toprule
\textbf{序号} & \textbf{批改原则（原意）} & \textbf{写作时必须做到}\\
\midrule
\endhead
\textbf{1} & 先核对写作身份、对象、目的及任务要点，再结合语言、结构和格式综合定档。 & 动笔前先把题干拆成“我是谁 / 写给谁 / 为什么写 / 要回应哪几个要点”。\\
\midrule
\textbf{2} & 按具体文体检查格式：\textbf{书信的称呼、结束语和署名；通知的标题和落款}。检查是否符合题目指定姓名，语气是否适合双方关系，\textbf{不把所有文章都改成正式公文}。 & 题干说 “use Li Ming instead”，署名就必须是 Li Ming；写给朋友和写给教授，语气必须不一样。\\
\midrule
\textbf{3} & 统计正文词数，称呼、结束语和署名单独说明，按题目要求评价长度，不编造固定扣分规则，不重复扣分。 & 约 100 词。称呼、落款不计入。\\
\midrule
\textbf{4} & 优先保证\textbf{信息完整、表达准确和简洁得体}。不以生僻词、复杂句或连接词数量决定分数，正确使用常见句式不扣分。 & 小作文的第一优先级是“该说的都说了”，第二优先级才是“说得漂亮”。\\
\bottomrule
\end{longtable}

\subsection{档位--行动对照表：想拿几分，就做这几件事}

\begin{longtable}{C{1.8cm} L{6.3cm} L{7.5cm}}
\toprule
\textbf{目标} & \textbf{大作文（15 分）要做的事} & \textbf{小作文（10 分）要做的事}\\
\midrule
\endfirsthead
\toprule
\textbf{目标} & \textbf{大作文（15 分）要做的事} & \textbf{小作文（10 分）要做的事}\\
\midrule
\endhead
\textbf{保底 10--12} & ① 图表读对，写清楚“图表的主题 + 1 个最突出的特征”；② 第一段总述句写对（图表类型 + 指标 + 对象 + 时间）；③ 第二、三段各写 2 句与图表相关的评论；④ 全文控制在 150--175 词。 & ① 称呼、正文、结束语、署名四件套齐全；② 题干给的写作要点逐条回应，一条不漏；③ 用自己最有把握的句式，不冒险。\\
\midrule
\textbf{冲 13--15} & 在上面基础上：① 第一段用“极值 + 对比”点出主要特征，而不是流水账；② 第二段用“是什么--为什么/影响”的论证结构，2--4 句成体系；③ 第三段回到图表本身收尾，\textbf{不出现数字}，也不过度升华；④ 每段之间有明确的逻辑衔接词；⑤ 至少用 2--3 个本册“弹药库”里的准确高级表达。 & ① 每个要点用两句来写（一句陈述、一句说明好处/影响），信息密度更高；② 语气与身份匹配（给朋友可以亲切，给教授要礼貌正式）；③ 第二段用具体的细节（时间/地点/活动内容）替代空泛的客套；④ 结尾句有个性、不套话。\\
\midrule
\textbf{必须避免} & 把静态图写成动态图；编造数据；混淆百分比与百分点；模板痕迹过重（通篇与题目无关的议论）。\textbf{另，本册建议：}第三段不再出现数字、全文控制在 190 词以内——这两条不是官方规则，是为了把时间留给回读和改错。 & 漏回应写作要点；格式残缺（通知没标题或落款）；署名写自己的真名；把小作文写成大作文式的空泛议论。\\
\bottomrule
\end{longtable}

\begin{note}
\textbf{一句话总结：}大作文的分水岭在“图表描述准不准 + 评论贴不贴图”；小作文的分水岭在“要点全不全 + 格式对不对”。这两件事都不需要天赋，只需要按本册的清单去核对。
\end{note}

\clearpage

\section{大作文总览：图表作文的三段结构与审题决策}
\markboth{大作文总览：图表作文的三段结构与审题决策}{}

\subsection{题干永远是同一个格式}

英语二大作文的 Directions 十年没变过，就是三句话加一个材料：

\begin{tpl}{Directions 模板（把年份换掉就是任何一年）}
\en{Directions:}\\
\en{Write an essay of about 150 words based on the chart below.}\\
\en{In your essay, you should}\\
\en{1) describe the chart briefly,}\\
\en{2) interpret the implied meaning, and}\\
\en{3) give your comments.}\\
\en{You should write neatly on the ANSWER SHEET. (15 points)}
\end{tpl}

这直接决定了段落分工：

\begin{longtable}{C{1.6cm} L{4.4cm} L{9.6cm}}
\toprule
\textbf{段落} & \textbf{对应题干要求} & \textbf{这一段的任务}\\
\midrule
\endfirsthead
\toprule
\textbf{段落} & \textbf{对应题干要求} & \textbf{这一段的任务}\\
\midrule
\endhead
\textbf{第 1 段} & \en{describe the chart briefly} & 说清这是**什么图、关于什么、什么时间范围**，并点出**最主要的特征**（极值 / 趋势 / 关键对比）。2--3 句。\\
\midrule
\textbf{第 2 段} & \en{interpret the implied meaning} & 解读：这些数据说明了什么？为什么会这样？会带来什么影响？4--5 句。\\
\midrule
\textbf{第 3 段} & \en{give your comments} & 评论：总结观点 + 建议/期待。3 句，**不再出现任何数字**。\\
\bottomrule
\end{longtable}

\begin{note}
题干固定，意味着\textbf{结构可以 100\% 预审题}——这正是本册模板的立足点。但要注意：从“套结构”到“套内容”只有一线之隔。阅卷指令里明确写了“模板文有明显的特征：内容优质但与题目无关”。所以本册的每个模板都带一个**锚定动作**：把自己写的每一句都指回图表里的某个具体信息。做不到这一点，模板就会从“高级感”变成“模板感”。
\end{note}

\subsection{审题七问：动笔前的 90 秒}

拿到图表，按顺序回答下面七个问题，答案决定了你用哪一套模板。

\begin{longtable}{C{0.9cm} L{3.6cm} L{11.0cm}}
\toprule
\textbf{序号} & \textbf{问题} & \textbf{怎么判断 / 为什么重要}\\
\midrule
\endfirsthead
\toprule
\textbf{序号} & \textbf{问题} & \textbf{怎么判断 / 为什么重要}\\
\midrule
\endhead
\textbf{1} & 图表类型是什么？ & 柱状图 \en{bar chart}、饼状图 \en{pie chart}、折线图 \en{line graph}、表格 \en{table}。第一句就要写对，写错了阅卷老师立刻察觉。\\
\midrule
\textbf{2} & \textbf{静态图还是动态图？} & \textbf{全篇最关键的一问。}横轴是“不同类别”→静态图；横轴是“不同时间”→动态图。静态图只比高低，动态图只讲变化。答案决定第 1 段和第 2 段的全部句式。\\
\midrule
\textbf{3} & 时态用什么？ & 描述图表本身用一般现在时（\en{The chart shows...}）。图内数据：指向过去的年份用一般过去时（\en{rose from 75.2 to 145.9}）；如果题目是一次“调查”且没有强调时间，用一般现在时陈述调查结果也完全可以（\en{The survey shows that 91.3\% of students mention...}）。有“预计/到 2030 年”用 \en{is projected to}。\textbf{唯一不能做的是同一张图的数据时态来回跳。}\\
\midrule
\textbf{4} & 主要信息是什么？ & 静态图找\textbf{极值、排序、聚类}；动态图找\textbf{总体趋势、极值、首末数据、交叉点}。这是第 1 段要写的东西。\\
\midrule
\textbf{5} & 描述对象和指标各是什么？ & 对象＝图里被比较的那些东西（空调/冰箱/洗衣机）；指标＝衡量的量（每百户拥有量、比例、人数）。第一段的骨架就是“指标 + 对象 + 范围限定”。\\
\midrule
\textbf{6} & 单位、口径是什么？ & \% 还是人数？“元”还是“万元”？“每百户”还是“每千人”？写错单位比写错数字更致命。\\
\midrule
\textbf{7} & 描述与评论怎么分配篇幅？ & 英语二可以五五开。基础一般：描述多写、评论写 3--4 句稳妥的；基础较好：描述力求精简（2 句点出关键对比），把火力放到第 2 段。无论哪种，\textbf{第 3 段都要回到图表}。\\
\bottomrule
\end{longtable}

\subsection{静态图 / 动态图判断树}

\begin{center}
\begin{tikzpicture}[
  font=\small,
  q/.style={draw=accent,line width=0.8pt,rounded corners=2pt,align=center,
            inner sep=6pt,fill=band,text width=4.6cm},
  a/.style={draw=ink,line width=0.6pt,rounded corners=1.5pt,align=left,
            inner sep=6pt,fill=white,text width=5.6cm,font=\footnotesize},
  e/.style={draw=mid,line width=0.5pt,rounded corners=1.5pt,align=center,
            inner sep=4pt,fill=white,text width=3.4cm,font=\scriptsize},
  ln/.style={draw=mid,line width=0.6pt,-{Stealth[length=4pt]}},
  lb/.style={font=\scriptsize\bfseries,color=accent,inner sep=2pt}
]
\node[q] (Q) at (0,0) {横轴 / 主维度是\textbf{时间}（年份、月份、时期）吗？};

\node[e] (Y) at (-4.0,-3.0) {\textbf{是} \(\rightarrow\) \textbf{动态图}};
\node[e] (N) at (4.2,-3.0) {\textbf{否} \(\rightarrow\) \textbf{静态图}};

\node[a] (D) at (-4.0,-7.0)
  {\textbf{只讲“变化”，不比“大小”}\\[2pt]
   第 1 段：总述 + 总体趋势 + 1--2 个最突出的对象\\[2pt]
   第 2 段：变化说明了什么（发展 / 进步 / 隐忧）\\[2pt]
   常用：\en{increase / decline / steadily / surge /} \en{from ... to ... / by ... / peak at}\\[3pt]
   \textbf{忌}：把不同对象的高低说成“变化”};

\node[a] (S) at (4.2,-7.0)
  {\textbf{只讲“高低”，不谈“变化”}\\[2pt]
   第 1 段：总述 + 极值 + 排序 + 聚类对比\\[2pt]
   第 2 段：这种结构说明了什么（观念 / 现状 / 隐忧）\\[2pt]
   常用：\en{account for / rank first / stand at /} \en{the highest / in contrast / respectively}\\[3pt]
   \textbf{忌}：给不同类别编出“上升/下降”};

\draw[ln] (Q.south) -- ++(0,-0.9) -| node[lb,pos=0.25]{是} (Y.north);
\draw[ln] (Q.south) -- ++(0,-0.9) -| node[lb,pos=0.25]{否} (N.north);
\draw[ln] (Y.south) -- (D.north);
\draw[ln] (N.south) -- (S.north);

\node[draw=light,line width=0.5pt,rounded corners=1.5pt,fill=band,align=left,
      inner sep=6pt,text width=15.2cm,font=\footnotesize] (H) at (0.1,-10.6)
 {\textbf{第三种：半动态图（横轴是年龄段 / 阶段，不是时间）}\ ——比如“不同年龄员工的工作满意度”。
  它两边都沾一点：可以按\textbf{动态}写（\en{People are more likely to study out of interest as they grow older.}），
  也可以按\textbf{静态}写（\en{Older people are more likely to study out of interest than younger ones.}）。
  \textbf{选定一个视角，通篇不要来回跳。}详见第 5 章。};

\draw[ln] (D.south) -- ++(0,-0.6) -| (H.west);
\draw[ln] (S.south) -- ++(0,-0.6) -| (H.east);
\end{tikzpicture}
\end{center}

\zone{判断错了会怎样？如果图里只有“2014 年三类电器的拥有量”，你写成 \en{air conditioners increased steadily}，就是\textbf{凭空造出一个时间趋势}——按批改原则第 2 条，这属于“把类别差异写成时间趋势”，是从 13--15 档直接掉到 4--6 档的错误。反过来，图里有 2014--2023 四年数据，你只写 \en{air conditioners ranked first}，只是浪费了信息，不至于跑题，但也拿不到“图表描述准确、评论充分”。}

\subsection{字数预算：150--175 词怎么分}

\begin{longtable}{C{1.6cm} C{1.9cm} C{1.9cm} L{9.4cm}}
\toprule
\textbf{段落} & \textbf{句数} & \textbf{词数} & \textbf{内容分配}\\
\midrule
\endfirsthead
\toprule
\textbf{段落} & \textbf{句数} & \textbf{词数} & \textbf{内容分配}\\
\midrule
\endhead
\textbf{第 1 段} & 2--3 句 & 40--55 词 & 第 1 句总述（图表类型 + 指标 + 对象 + 时间/范围）；第 2 句主要特征（极值或趋势 + 2--4 个关键数字）；第 3 句（可选）一句过渡到评论。\\
\midrule
\textbf{第 2 段} & 4--5 句 & 60--75 词 & 总领 1 句 + 分论点 3--4 句。每句都要能指回图表里的某个信息。\\
\midrule
\textbf{第 3 段} & 3 句 & 35--45 词 & 总结 1 句 + 概括/建议 1--2 句。不出现数字。\\
\bottomrule
\end{longtable}

\begin{note}
\textbf{为什么不是越长越好？}第一，批改原则第 4 条明确说了“约 150 词”不擅自改成“不少于 150 词”，多写不加分。第二，考场上大作文的合理时间预算是 \textbf{25--30 分钟}，写 220 词意味着没有时间回读和改错，而“语法错误极少”恰恰是 13--15 档的要求。第三，答案卡的行数是有限的。
\end{note}

\subsection{九句骨架：全篇的“承重墙”}

把全篇拆成九个功能位，你就拥有了一张不会跑偏的施工图。第 3--8 章做的是同一件事：给每一个功能位配一套可替换的句子库。

\begin{longtable}{C{1.0cm} C{1.5cm} L{3.5cm} L{9.0cm}}
\toprule
\textbf{位} & \textbf{段} & \textbf{功能} & \textbf{句子任务}\\
\midrule
\endfirsthead
\toprule
\textbf{位} & \textbf{段} & \textbf{功能} & \textbf{句子任务}\\
\midrule
\endhead
1 & 第 1 段 & 总述 & \en{The bar chart illustrates ...}（图表类型 + 指标 + 对象 + 范围）\\
\midrule
2 & 第 1 段 & 主要特征 & 极值 / 趋势 / 关键对比 + 2--4 个数字\\
\midrule
3 & 第 1 段 & 过渡 & 一句把描述引向评论（可省）\\
\midrule
4 & 第 2 段 & 总领 & 数据说明了什么（一句话概括你的解读方向）\\
\midrule
5 & 第 2 段 & 分论点一 & 是什么 / 为什么（图表锚定）\\
\midrule
6 & 第 2 段 & 分论点二 & 为什么 / 影响（图表锚定）\\
\midrule
7 & 第 2 段 & 分论点三 & 深入一层，或指出隐忧（图表锚定）\\
\midrule
8 & 第 3 段 & 总结 & \en{Overall, ...}（回到图表的核心矛盾）\\
\midrule
9 & 第 3 段 & 建议 / 期待 & 政府 / 学校 / 媒体 / 个人应该做什么，或一句有分寸的展望\\
\bottomrule
\end{longtable}

\begin{note}
第 5--7 位是“变量位”，必须按题目现场生成；1、2、8、9 位是“常量位”，可以提前背到条件反射。这就是为什么本册反复强调：\textbf{结构预审题，内容现场填}。
\end{note}

\clearpage

\section{第 1 段（静态图）：图表信息说明}
\markboth{第 1 段（静态图）：图表信息说明}{}

\begin{note}
\textbf{适用前提：}横轴是不同\textbf{类别}（国家、品牌、活动、年龄段），没有时间维度，或者虽有时间但题目只让你比“同一时点的不同对象”。此时你的任务只有一件事：\textbf{把高低关系写清楚}。任何“上升/下降/增长”都不该出现在这一段的句子里。
\end{note}

\subsection{总述句：静态图的第一句}

总述句的骨架是 \textbf{图表类型 + 指标 + 对象 + 范围限定}。这四件套来自御风的“标准范式”，本册把它固定下来：

\begin{tpl}{总述句公式}
\en{The} \fillin{图表类型（bar chart / pie chart / table）} \en{shows / illustrates / presents} \fillin{指标（the proportion / the number / the percentage）} \en{of} \fillin{对象} \en{in / among} \fillin{范围限定（人群、地点、年份）}.
\end{tpl}

\begin{eg}
\en{The bar chart illustrates the primary leisure activities of the elderly in a community.}\\
中文：该柱状图展示了某社区老年人的主要休闲活动。\\[3pt]
\en{The pie chart presents the factors that consumers take into account when choosing a restaurant.}\\
中文：该饼状图展示了消费者选择餐厅时所考虑的因素。\\[3pt]
\en{The table shows the results of a job satisfaction survey among employees of different age groups in a certain company.}\\
中文：该表格展示了某公司不同年龄段员工工作满意度调查的结果。
\end{eg}

\begin{note}
\textbf{为什么第一句这么重要？}它是阅卷老师读到的第一行。图表类型写对、指标和对象写准，老师立刻知道“这个学生读懂了图”；反之若开头就写错图表类型，后面的印象分会一路走低。
\end{note}

\subsubsection{总述句的可替换零件}

\begin{longtable}{L{3.2cm} L{5.4cm} L{6.9cm}}
\toprule
\textbf{零件} & \textbf{可选词} & \textbf{说明}\\
\midrule
\endfirsthead
\toprule
\textbf{零件} & \textbf{可选词} & \textbf{说明}\\
\midrule
\endhead
图表类型 & \en{table}（表格）、\en{bar chart}（柱状图）、\en{line graph}（折线图）、\en{pie chart}（饼状图） & 双图写 \en{The two charts} 或 \en{The chart and the table}。表格不要写成 \en{table chart}。\\
\midrule
谓语动词 & \en{shows}、\en{illustrates}、\en{presents}、\en{reveals}、\en{depicts}、\en{provides a snapshot of} & \en{show} 本身已经地道且万能，不必刻意替换。\en{compare} 用于“对比”类题目：\en{The chart compares ...}。\\
\midrule
指标词 & \en{the proportion}、\en{the percentage}、\en{the share}、\en{the number}、\en{the amount}、\en{the figure}、\en{the index}、\en{the rate} & \en{number} 接可数（人数），\en{amount} 接不可数（时间、金钱），\en{proportion/percentage/share} 接比例。\\
\midrule
范围限定 & \en{among college students}、\en{in a certain city}、\en{across the country}、\en{at the end of 2023}、\en{aged 15 to 69}、\en{per 100 households} & 时间、地点、人群、单位口径都属于“范围限定”，\textbf{宁可多写一个介词短语，也不要漏掉单位}。\\
\bottomrule
\end{longtable}

\subsection{主要特征句：静态图的第二、三句}

静态图的信息密度通常比动态图低，所以第 1 段要写得更“满”一点。用下面两套写法中的任意一套。

\subsubsection{写法 A：逐项写法（信息量大、稳，适合基础一般）}

每一项给一句，句式变化靠“换个连接词 + 换个动词”实现。这是 mmg 原书的思路，本册保留了它的**高级表达表**并补齐了全部例句。

\begin{tpl}{逐项写法公式}
\fillin{关于} \fillin{什么}，\fillin{指标} \fillin{怎么样}.
\end{tpl}

\begin{longtable}{C{3.6cm} C{3.2cm} C{3.2cm} L{5.5cm}}
\toprule
\textbf{关于（前置）} & \textbf{指标} & \textbf{怎么样（动词）} & \textbf{备注}\\
\midrule
\endfirsthead
\toprule
\textbf{关于（前置）} & \textbf{指标} & \textbf{怎么样（动词）} & \textbf{备注}\\
\midrule
\endhead
\en{Concerning} & \en{the percentage} & \en{accounts for} & 最通用，占……\\
\en{Regarding} & \en{the share} & \en{makes up} & \en{make up} 后接比例，合起来是“构成”\\
\en{As for} & \en{the quantity} & \en{hits} / \en{reaches} & 口语感稍强，适合过渡\\
\en{In the case of} & \en{the amount} & \en{stands at} & 强调“停在一个数值上”\\
\en{With regard to} & \en{the proportion} & \en{stays at} & 强调“保持不变”\\
\en{In terms of} & \en{the figure} & \en{is estimated at} & 用于估算值\\
\en{When it comes to} & \en{the index} & \en{ranks first / last} & 用于排序\\
\bottomrule
\end{longtable}

\begin{note}
这些词不是越“高级”越好，而是\textbf{越准越好}。用 \en{hit} 去描述 3\% 会显得夸张；用 \en{stay at} 去描述一个正在变的数字则自相矛盾。批改原则第 5 条说得很清楚：不以生僻词、复杂句决定分数。选词的唯一标准是“\textbf{换成你自己的话，意思还是不是一样}”。
\end{note}

\begin{tpl}{逐项写法例句（全部可替换）}
1. \en{Regarding the proportion of watching TV, the figure stood at 90.8\%, representing the largest share among all categories.}\\
2. \en{As for household expenditure on healthcare, spending reached 8{,}000 yuan per capita, reflecting rising health awareness.}\\
3. \en{With respect to the number of museum visitors, the figure surged to 1.2 billion, demonstrating the rapid spread of cultural tourism.}\\
4. \en{In terms of public transportation, the percentage of daily commuters remained at approximately 55\%, suggesting stable demand.}\\
5. \en{When it comes to gardening, the participation rate stayed at \fillin{数值}, far below that of walking.}\\
6. \en{Concerning those who study out of interest, the share rose with age, from 12\% among the youngest group to 41\% among the oldest.}
\end{tpl}

\subsubsection{写法 B：聚类写法（精简、有高级感，适合冲高分）}

聚类写法把若干个同类项打包，用一句话说完，是御风原书最推荐的手法。它需要三个东西：\textbf{极值判断} + \textbf{并列/对比衔接词} + \textbf{概括性抽象名词}。

\begin{tpl}{聚类写法公式}
\en{Overall,} \fillin{事物 A} \en{ranks first, accounting for} \fillin{数值 A}\en{, followed by} \fillin{事物 B}\en{, which makes up} \fillin{数值 B}\en{.} \fillin{事物 C 和 D} \en{come next, accounting for} \fillin{数值 C} \en{and} \fillin{数值 D} \en{respectively, while} \fillin{事物 E} \en{ranks last at} \fillin{数值 E}\en{.}
\end{tpl}

\begin{eg}
\en{Overall, watching TV ranks first, accounting for 90.8\%, followed by walking, which makes up 68.3\%. Gardening, reading and playing chess come next, accounting for \fillin{数值} each, while the remaining activities are mentioned by fewer than one in ten respondents.}
\end{eg}

\begin{tpl}{把聚类写法的两头写扎实}
\textbf{开头的概括词}：\en{Overall,} / \en{Broadly speaking,} / \en{As can be seen from the chart,}\\[3pt]
\textbf{表示“并列/接近”}：\en{similarly}、\en{likewise}、\en{in much the same way}、\en{at a similar level}、\en{the same is true of ...}\\[3pt]
\textbf{表示“反差/落差”}：\en{in contrast}、\en{on the other hand}、\en{by comparison}、\en{distinct from}、\en{the gap between X and Y}、\en{markedly lower / higher than}\\[3pt]
\textbf{表示“被压制/微不足道”}：\en{are overshadowed by}、\en{are far behind}、\en{account for only a small share}、\en{none of them exceeds} \fillin{数值}\\[3pt]
\textbf{总结“整体结构”}：\en{The overall structure reveals ...} / \en{Such a distribution suggests that ...}
\end{tpl}

\subsection{静态图描述功能句十条（可直接抄）}

这十条来自 mmg 原书，本册逐条补齐了中文、修正了错误，并标注了使用场景。

\begin{longtable}{C{0.9cm} L{7.3cm} L{6.9cm}}
\toprule
\textbf{编号} & \textbf{英文功能句} & \textbf{中文 / 使用场景}\\
\midrule
\endfirsthead
\toprule
\textbf{编号} & \textbf{英文功能句} & \textbf{中文 / 使用场景}\\
\midrule
\endhead
1 & \en{As for \fillin{事物 1}, it accounts for \fillin{百分比 1}.} & 就事物 1 而言，占比为百分比 1。最通用的一条。\\
\midrule
2 & \en{Turning to \fillin{事物 1}, it makes up \fillin{百分比 1} of the total.} & 转而看事物 1，其占总体的百分比 1。用于“换个对象继续说”。\\
\midrule
3 & \en{Regarding \fillin{事物 1}, the share stands at \fillin{百分比 1}.} & 关于事物 1，占比为百分比 1。\en{stand at} 强调“停在某数值”。\\
\midrule
4 & \en{For \fillin{事物 1}, the proportion is \fillin{百分比 1}.} & 事物 1 的比例为百分比 1。最简句式，用于不想重复花样时。\\
\midrule
5 & \en{In the case of \fillin{事物 1}, the slice reaches \fillin{百分比 1}.} & 事物 1 的扇形占比达到百分比 1。\textbf{只用于饼图}（\en{slice} 是饼图的扇形）。\\
\midrule
6 & \en{As for \fillin{事物 1} and \fillin{事物 2}, they register \fillin{百分比 1} and \fillin{百分比 2} respectively.} & 就事物 1 与事物 2 而言，它们分别为百分比 1 和百分比 2。\en{respectively} 是配对神器，务必用对顺序。\\
\midrule
7 & \en{Compared with \fillin{事物 2} at \fillin{百分比 2}, \fillin{事物 1} comes in at \fillin{百分比 1}.} & 与处于百分比 2 的事物 2 相比，事物 1 为百分比 1。用于制造对比。\\
\midrule
8 & \en{While \fillin{事物 2} holds \fillin{百分比 2}, \fillin{事物 1} takes \fillin{百分比 1}.} & 事物 2 占百分比 2，而事物 1 为百分比 1。\en{while} 表对比，非常地道。\\
\midrule
9 & \en{\fillin{事物 1} accounts for \fillin{百分比 1}, whereas \fillin{事物 3} posts \fillin{百分比 3}.} & 事物 1 占百分比 1，而事物 3 为百分比 3。\en{whereas} 比 \en{but} 更书面。\\
\midrule
10 & \en{\fillin{事物 1} and \fillin{事物 2} together make up \fillin{百分比合计}, with individual shares of \fillin{百分比 1} and \fillin{百分比 2}.} & 事物 1 与事物 2 合计占百分比合计，单项占比分别为百分比 1 与百分比 2。用于“打包两个主要项”。\\
\bottomrule
\end{longtable}

\subsection{把前三句连起来：静态图第 1 段的三种成品}

\subsubsection{成品一：极值型（最常用）}

\begin{tpl}{静态图第 1 段 · 极值型}
\en{The bar chart illustrates the primary leisure activities of the elderly in a community. Overall, watching TV and walking stand out as the two most common activities, accounting for 90.8\% and 68.3\% respectively. By contrast, other activities, including gardening, reading and playing chess, are much less popular, none of them exceeding 35\%.}
\end{tpl}

\begin{note}
第 1 句总述（图表类型 + 指标 + 对象 + 范围）；第 2 句用 \en{Overall} 亮出\textbf{两个极值}并用 \en{respectively} 配对；第 3 句用 \en{By contrast} 一次性打包剩下的所有项，用“封顶值”（\en{none of them exceeding 35\%}）代替罗列。全程只有 3 个数字，但高低关系完全清楚。
\end{note}

\subsubsection{成品二：排序型（适合饼图和表格）}

\begin{tpl}{静态图第 1 段 · 排序型}
\en{The pie chart presents the factors that consumers take into account when choosing a restaurant. Taste ranks first, accounting for 36.3\%, followed by environment and service, which make up 24.1\% and 19.5\% respectively. Price and other factors come next, while reputation ranks last at less than 10\%.}
\end{tpl}

\subsubsection{成品三：带过渡句型（想在第 1 段就点出观点）}

\begin{tpl}{静态图第 1 段 · 带过渡句}
\en{The bar chart illustrates a survey on students' gains from a labor-practice class at a university. The two most frequently cited benefits are the acquisition of relevant knowledge and the improvement of practical skills, accounting for 84.5\% and 91.3\% respectively, whereas the enjoyment factor represents a mere 54.5\%. Such a distribution is revealing, and the reasons behind it deserve closer examination.}
\end{tpl}

\begin{note}
最后一句是御风原书的过渡句 \en{Succinct as it is, the chart conveys abundant insights, and below are my perspectives.} 的改写版。原句本身没问题（倒装结构漂亮），但被用滥后会让阅卷老师一眼识别出模板。本册建议：\textbf{过渡句要么自己写具体的，要么干脆不写}，把省下的字数还给第 2 段。
\end{note}

\subsection{静态图第 1 段自查}

\begin{longtable}{C{0.8cm} L{13.6cm}}
\toprule
& \textbf{检查项}\\
\midrule
\endfirsthead
\toprule
& \textbf{检查项}\\
\midrule
\endhead
\cb & 第一句写对了图表类型（\en{bar chart / pie chart / table}），没有把表格写成 \en{table chart}。\\
\cb & 第一句包含“指标 + 对象 + 范围限定（时间/人群/单位）”。\\
\cb & 全文\textbf{没有出现任何表示时间变化的动词}（\en{increase / decrease / rise / fall / grow / decline}）。\\
\cb & 用的每一个数字都能在图上指出来，没有“约等于”或“编一个”。\\
\cb & 单位写对了（\%、人数、元、每百户……），没有把 \% 和百分点混用。\\
\cb & 描写了\textbf{至少一个对比}（谁高于谁 / 谁只有谁的几分之一），而不是干巴巴列数字。\\
\cb & 数字总数控制在 3--5 个以内（批改原则第 3 条：不要求罗列全部数字）。\\
\cb & 全段 2--3 句，词数 40--55 之间。\\
\bottomrule
\end{longtable}

\clearpage

\section{第 1 段（动态图）：图表信息说明}
\markboth{第 1 段（动态图）：图表信息说明}{}

\begin{note}
\textbf{适用前提：}横轴是不同\textbf{时间}（年份、月份、时期、阶段），你要讲的是“变化”。这一段的句子里\en{increase / decrease / rise / fall / grow / decline / remain stable} 之类的动词是必需品的，反而\textbf{不能}出现“某类占多少比例”这种静态比较句去替代趋势——除非你要补一个时点数据（如“到 2023 年达到 145.9”）。
\end{note}

\subsection{总述句：动态图的第一句}

\begin{tpl}{总述句公式（动态）}
\en{The} \fillin{图表类型} \en{illustrates the changes in} \fillin{指标 + 对象} \en{from} \fillin{起始时间} \en{to} \fillin{结束时间}\en{, as well as projected changes to} \fillin{预测年份（若有）}\en{.}
\end{tpl}

\begin{eg}
\en{The line graph illustrates the changes in the number of museums and museum visits in China from 2013 to 2017.}\\
中文：该折线图展示了 2013 至 2017 年中国博物馆数量与参观人次的变化。\\[3pt]
\en{The line graph shows the amount of energy consumed in the US from 1980 to 2015, as well as projected consumption to 2030.}\\
中文：该折线图展示了 1980 至 2015 年美国能源消耗量，以及到 2030 年的预测消耗量。\\[3pt]
\en{The bar chart traces the change in the volume of express delivery business in China between 2018 and 2022.}\\
中文：该柱状图呈现了 2018 至 2022 年中国快递业务量的变化。
\end{eg}

\begin{note}
\textbf{两个容易丢分的细节：}
一是\textbf{“预测”要用对}。图里有虚线延伸到未来年份，就必须用 \en{projected} / \en{is expected to} / \en{is predicted to}，不能直接说 \en{will}，更不能用过去时。
二是\textbf{年份范围要写全}。\en{from 1980 to 2015} 里的两个数字都属于“范围限定”，漏一个就等于漏了任务信息。
\end{note}

\subsection{动态描述的标准范式（全册最重要的一个公式）}

御风原书给出了这个公式，本册把它规范化：

\begin{tpl}{动态图描述六件套}
（可选：逻辑衔接词）+ 【时间（点 / 时期）】+ 【事物】+ 【趋势动词】+ 【数值】+ 【意义】
\end{tpl}

\begin{note}
六件套里的“\textbf{意义}”是很多人忽略的一环，但它正是拉开差距的地方：\en{, indicating that ...} / \en{, reflecting ...} / \en{, underscoring ...} 这样一个分词短语，既做了数据解读，又把描述和评论缝在了一起，正好命中评分档位里“描述准确、评论充分”的要求。\textbf{建议每一段描述里至少有两句带“意义”尾。}
\end{note}

\begin{tpl}{六件套的六个成品（全部可直接替换）}
1. \en{Between 2015 and 2024, the proportion of online learners underwent a steady increase, from 25\% to 58\%, indicating the growing popularity of digital education.}\\[2pt]
2. \en{During the decade, the number of electric vehicles surged from 3 million to 18 million, reflecting the accelerating transition toward green transportation.}\\[2pt]
3. \en{By 2022, the figure had climbed to a peak of 72\%, underscoring the widespread acceptance of the policy.}\\[2pt]
4. \en{During the period concerned, the figure climbed steadily to 72\%, highlighting the growing acceptance of digital technology.}\\[2pt]
5. \en{Between 2018 and 2024, the percentage rose by a modest 5 percentage points, suggesting gradual progress.}\\[2pt]
6. \en{The figure was twice as high as that recorded in 2015, underscoring impressive growth.}
\end{tpl}

\begin{note}
注意第 5 句里 \en{by a modest 5 percentage points} 的用法——这是“百分点”最正确的写法，也是本册在第 1 章红线二里强调的那一点。反过来，如果原文是 25\% 涨到 50\%，写 \en{rose by 25 percentage points} 对，写 \en{rose by 25\%} 就错了（那等于说从 25\% 涨到 31.25\%）。
\end{note}

\subsection{趋势与程度的速查词（完整版见第 13 章）}

\begin{longtable}{L{2.8cm} L{5.8cm} L{6.6cm}}
\toprule
\textbf{方向} & \textbf{常规词（稳）} & \textbf{加重词（要慎用）}\\
\midrule
\endfirsthead
\toprule
\textbf{方向} & \textbf{常规词（稳）} & \textbf{加重词（要慎用）}\\
\midrule
\endhead
上升 & \en{climb}、\en{increase}、\en{rise}、\en{grow}、\en{go up} & \en{surge}、\en{soar}、\en{rocket}、\en{shoot up}、\en{leap}、\en{spiral}（程度重）\\
\midrule
下降 & \en{dip}、\en{slide}、\en{fall}、\en{decline}、\en{drop}、\en{decrease}、\en{diminish} & \en{plunge}、\en{plummet}、\en{nosedive}、\en{tumble}、\en{slump}、\en{plumb}（程度重）\\
\midrule
平稳 & \en{remain stable}、\en{level off}、\en{plateau at}、\en{stay roughly constant}、\en{hold steady at} & ——\\
\midrule
波动 & \en{fluctuate}、\en{show slight fluctuations}、\en{swing between X and Y} & \en{volatility increased}（用于波动加大时）\\
\midrule
程度（轻） & \en{slightly}、\en{modestly}、\en{moderately}、\en{marginally} & ——\\
\midrule
程度（重） & \en{considerably}、\en{remarkably}、\en{notably}、\en{markedly}、\en{substantially}、\en{significantly}、\en{dramatically} & \en{radically}、\en{exponentially}（用于成倍增长时）\\
\bottomrule
\end{longtable}

\zone{加重词用错了会显得“读不懂图”。图里从 80\% 涨到 85\%，你写 \en{soared}，阅卷老师会认为你不懂 \en{soar} 的含义；图里从 3 涨到 18（6 倍），你写 \en{increased slightly} 同样是失真。判断标准只有一个：\textbf{变化幅度和图里看起来的感觉是否一致}。拿不准就用 \en{rose from X to Y}，永不出错。}

\subsection{动态图描述功能句（可直接抄）}

\begin{longtable}{C{0.9cm} L{7.3cm} L{6.9cm}}
\toprule
\textbf{编号} & \textbf{英文功能句} & \textbf{中文 / 使用场景}\\
\midrule
\endfirsthead
\toprule
\textbf{编号} & \textbf{英文功能句} & \textbf{中文 / 使用场景}\\
\midrule
\endhead
1 & \en{From \fillin{时间 1} to \fillin{时间 2}, \fillin{事物 1} rose from \fillin{数值 1} to \fillin{数值 2}, showing a steady upward trend.} & 自时间 1 至时间 2，事物 1 由数值 1 升至数值 2，呈现稳定上升趋势。最常用的一条。\\
\midrule
2 & \en{Over the same span, \fillin{事物 2} fell from \fillin{数值 3} to \fillin{数值 4}, indicating a clear decline.} & 同一时期，事物 2 由数值 3 降至数值 4，表明明显下滑。承接上句做对比。\\
\midrule
3 & \en{By \fillin{时间 2}, \fillin{事物 3} plateaued around \fillin{数值 5}, and its growth slowed and stabilized.} & 到时间 2 时，事物 3 在数值 5 附近进入平台期，增速放缓并趋于稳定。用于“先涨后平”的形状。\\
\midrule
4 & \en{After \fillin{时间 1}, \fillin{事物 1} accelerated, and the increase from \fillin{数值 1} to \fillin{数值 2} outpaced earlier gains.} & 时间 1 之后，事物 1 出现加速，从数值 1 到数值 2 的增幅超过此前增幅。用于“后期加速”。\\
\midrule
5 & \en{Volatility increased for \fillin{事物 3}, with wider swings around \fillin{数值区间}, and short-term fluctuations became more frequent.} & 事物 3 的波动性上升，在数值区间附近的振幅加大，短期起伏更为频繁。用于锯齿状折线。\\
\midrule
6 & \en{During this period, \fillin{事物 A} witnessed a sharp increase, rising from \fillin{数值 1} to \fillin{数值 2}.} & 在此期间，事物 A 出现了明显上升，从数值 1 上升至数值 2。\en{witness} 的主语是“事物”，非常地道。\\
\midrule
7 & \en{Meanwhile, the number of \fillin{事物 B} also experienced a remarkable increase, eventually reaching approximately \fillin{数值}.} & 与此同时，事物 B 的数量也出现了显著增长，最终达到约数值。用于“第二个对象”。\\
\midrule
8 & \en{In contrast, \fillin{事物 C} dropped slightly, from \fillin{数值} to \fillin{数值}.} & 相较之下，事物 C 略有下降，从数值降至数值。用于“逆势”对象。\\
\bottomrule
\end{longtable}

\subsection{动态图第 1 段的成品与批注}

\subsubsection{成品一：一句话总趋势 + 两个对象（最稳）}

\begin{tpl}{动态图第 1 段 · 成品一}
\en{The chart clearly illustrates the noteworthy change in the health literacy level of residents in our country. As is shown in the line graph, a clear upward trend can be observed over the past decade: the proportion of people aged between 15 and 69 possessing basic health literacy rose sharply from 8.8\% to 25.4\%, almost tripling in a mere ten years.}
\end{tpl}

\begin{note}
\textbf{得分点解析：}① 第一句写全了“指标 + 对象 + 范围”；② \en{As is shown in the line graph} 明确点出图表类型；③ \en{almost tripling in a mere ten years} 用一个分词短语把“变化幅度”具象化——比单说“上升”有力得多；④ 只有两个数字，但趋势和幅度都说清楚了。这正是批改原则第 3 条欣赏的“突出主要特征，不罗列全部数字”。
\end{note}

\subsubsection{成品二：多对象对比型}

\begin{tpl}{动态图第 1 段 · 成品二}
\en{The bar chart traces the change in the volume of express delivery business in China between 2018 and 2022. During this period, the figure witnessed a dramatic increase, rising from 50.7 billion pieces to 110.6 billion pieces, which means it more than doubled within only five years.}
\end{tpl}

\subsubsection{成品三：静态起步 + 动态推进（表格题常用）}

\begin{tpl}{动态图第 1 段 · 成品三}
\en{The table illustrates changes in the average ownership of three major durable consumer goods per 100 Chinese households at the end of 2014, 2017, 2020 and 2023. Overall, all three categories registered growth, though at noticeably different rates. Air conditioners showed the most striking increase, rising from 75.2 units in 2014 to 145.9 in 2023, whereas washing machines grew more moderately, from 83.7 to 98.2.}
\end{tpl}

\begin{note}
\textbf{得分点解析：}① 表格题的“范围限定”最容易漏——\en{per 100 Chinese households} 和四个年份都是必写信息；② \en{Overall} 先给总趋势，再给极值对象，最后一句用 \en{whereas} 做对比收尾，三句话完成了“总—分—对比”的完整结构；③ \en{though at noticeably different rates} 是一个非常高级的插入语，一句话就为下文的分类讨论铺好了路。
\end{note}

\subsection{动态图第 1 段自查}

\begin{longtable}{C{0.8cm} L{13.6cm}}
\toprule
& \textbf{检查项}\\
\midrule
\endfirsthead
\toprule
& \textbf{检查项}\\
\midrule
\endhead
\cb & 第一句写对了图表类型，且写全了时间范围（起止都要有）。\\
\cb & 图里若有“预测年份/虚线”，用了 \en{projected} / \en{is expected to}，没有误用过去时或 \en{will}。\\
\cb & 至少有一句点明了\textbf{总体趋势}（\en{Overall, all three categories registered growth...}）。\\
\cb & 至少有一个\textbf{极值}（最高的那个对象 / 最大的那个变化）被单独点出。\\
\cb & 每个趋势动词的“力度”与图中幅度匹配（没有把小涨幅写成 \en{soared}）。\\
\cb & 数值写法正确：用 \en{from X to Y} 表达区间，用 \en{by X} 表达幅度，\en{by 5 percentage points} 没有写成 \en{by 5\%}。\\
\cb & 至少有两句带了“意义尾”（\en{, indicating / reflecting / underscoring ...}）。\\
\cb & 数字总数控制在 4--6 个以内。\\
\cb & 全段 2--3 句，词数 40--55 之间。\\
\bottomrule
\end{longtable}

\clearpage

\section{第 1 段的三种特殊情况：半动态图、双图、表格}
\markboth{第 1 段的三种特殊情况}{}

\subsection{半动态图：横轴是“年龄段 / 阶段”}

\begin{note}
\textbf{什么是半动态图？}横轴不是年份，而是“不同年龄段”“不同年级”“不同阶段”（如 2012 年英语二：不同年龄段员工的工作满意度）。它同时具备动态与静态的特点，所以叫半动态。
\end{note}

这类题最大的好处是\textbf{切入角度多}，你可以二选一：

\begin{longtable}{C{2.6cm} L{5.9cm} L{5.9cm}}
\toprule
\textbf{视角} & \textbf{写法} & \textbf{例句}\\
\midrule
\endfirsthead
\toprule
\textbf{视角} & \textbf{写法} & \textbf{例句}\\
\midrule
\endhead
\textbf{动态视角} & 把“年龄”当作“时间”来写，讲“随着……而变化” & \en{With age, people are more likely to study out of interest.}\newline \en{The proportion of people who study out of interest increases with age.}\\
\midrule
\textbf{静态视角} & 把不同年龄组当作“不同类别”，横向比高下 & \en{Older people are more likely to study out of interest than younger people.}\\
\bottomrule
\end{longtable}

\zone{\textbf{选定一个视角，全篇不要来回跳。}最典型的错误是：第 1 段用静态视角写“中年的不满意率最高”，第 2 段却写“随着年龄增长，满意度先降后升”。后者其实是在叙述图里\textbf{并不存在}的一条时间轨迹，属于批改原则第 2 条明令禁止的“把类别差异写成时间趋势”。}

\subsubsection{年龄相关的表达（务必背准）}

\begin{longtable}{C{4.4cm} L{10.0cm}}
\toprule
\textbf{表达方式} & \textbf{例句}\\
\midrule
\endfirsthead
\toprule
\textbf{表达方式} & \textbf{例句}\\
\midrule
\endhead
\en{People aged + 年龄段} & \en{People aged under 24}；\en{people aged between 30 and 39}；\en{people aged over 50}\\
\midrule
\en{in + 年龄层} & \en{People in their twenties}；\en{people in their fifties}（注意：指整十年段，用复数）\\
\midrule
\en{Age group} & \en{the under-24 age group}；\en{the 41--50 age group}\\
\midrule
\en{-year-old} 复合形容词 & \en{under-24-year-old people}（作定语时 \en{year} 不加 \en{s}）\\
\bottomrule
\end{longtable}

\begin{tpl}{半动态图第 1 段通用句式}
\en{The table demonstrates the results of a survey conducted among people of different age groups. While \fillin{某一组} show clear \fillin{态度 A}, their counterparts \fillin{另外两组} maintain a comparatively \fillin{态度 B} attitude. Notably, \fillin{某一组} also stand out as the only group that \fillin{独特表现}, with a figure of \fillin{数值}. Such a contrast suggests that \fillin{你的初步解读}.}
\end{tpl}

\begin{eg}
\en{The table demonstrates the results of a job satisfaction survey conducted in a certain company. While employees aged 41 to 50 show clear discontent, their counterparts aged under 41 and over 50 maintain a comparatively favorable attitude towards their jobs. The dissatisfaction rate of the 41--50 age group reaches 64\%, exceeding the combined rates of the other two groups. Such a contrast suggests that employees at different stages of life hold markedly different views on their work.}
\end{eg}

\subsection{双图表}

\begin{note}
2024 年英语一出现过“漫画 + 图表”的混合题，2026 年英语一出现了“饼图 + 柱状图”的双图题。英语二目前还没考过双图，但大纲允许，值得准备。\textbf{好消息是：双图在写法上没有任何新东西}——无非把总述句写成 \en{X and Y}，然后分别按“静态”或“动态”各自写一句。
\end{note}

\begin{tpl}{双图第 1 段公式}
\en{The} \fillin{图 A 类型} \en{shows} \fillin{图 A 的指标与对象}\en{, while the} \fillin{图 B 类型} \en{reveals} \fillin{图 B 的指标与对象}\en{.} \fillin{用一句话各自点出最突出的特征}\en{.}
\end{tpl}

\begin{eg}
\en{Two charts present consumers' acceptance of elderly-care robots and their primary concerns. The pie chart shows that 72.1\% of respondents express at least partial acceptance, whereas the remaining 27.9\% convey outright rejection. According to the bar chart, safety is regarded as the dominant concern, accounting for 46.3\%, followed by price and convenience, at approximately 25\% and 10\% respectively.}
\end{eg}

\begin{note}
\textbf{双图的三个注意事项：}
① 总述句一定要让读者知道“有两张图”，不要只写一张；
② 两张图的顺序要和题干或图表标注一致，不要自作主张颠倒；
③ 字数紧张时，每张图只写一句、只挑一个最突出的特征，把篇幅留给第 2 段。
\end{note}

\subsection{表格题}

表格题本质上可以是静态、动态或半动态，\textbf{它的特殊性在排版，不在写法}。需要注意三点：

\begin{longtable}{C{0.9cm} L{13.4cm}}
\toprule
& \textbf{表格题专项提醒}\\
\midrule
\endfirsthead
\toprule
& \textbf{表格题专项提醒}\\
\midrule
\endhead
1 & \textbf{表格的“表头”就是你的范围限定。}单位（\en{per 100 households}、\en{yuan per capita}、\en{\%}）、年份、人群都要在第一句里交代清楚。2025 年英语一表格题的关键信息就是 \en{per 100 Chinese households}——漏掉它，所有数字都失去意义。\\
\midrule
2 & \textbf{别把表格写成流水账。}表格数据最全，最容易写成“X 是 1，Y 是 2，Z 是 3”。正确做法是先给\textbf{总体判断}（\en{Overall, all three categories registered growth, though at noticeably different rates.}），再挑 2--3 个最有说头的对象。\\
\midrule
3 & \textbf{注意行列交叉的“组合信息”。}例如“2014 年最低、2023 年反而最高”这类\textbf{排名变化}，是表格特有的观察点，写出来会显得你读得很细。\\
\bottomrule
\end{longtable}

\subsection{第 1 段三种情况的共同自查}

无论静态、动态还是半动态，第 1 段都要过这四关：

\begin{longtable}{C{0.8cm} L{13.6cm}}
\toprule
& \textbf{四关}\\
\midrule
\endfirsthead
\toprule
& \textbf{四关}\\
\midrule
\endhead
\textbf{一} & \textbf{类型关}：图表类型写对，且与图一致。\\
\midrule
\textbf{二} & \textbf{口径关}：指标、对象、时间、单位（\%/人数/元/每百户）一个不漏。\\
\midrule
\textbf{三} & \textbf{匹配关}：句子的“动/静”性质与图的类型匹配（静态不谈变化，动态不做纯比较）。\\
\midrule
\textbf{四} & \textbf{取舍关}：数字不超过 6 个，“最突出的特征”必须被单独点出。\\
\bottomrule
\end{longtable}

\clearpage

\section{第 2 段：描述与解读（全篇得分的重心）}
\markboth{第 2 段：描述与解读}{}

\begin{note}
第 2 段对应题干里的 \en{interpret the implied meaning}。判分档位里“评论相关且充分”这五个字，说的就是这一段。它也是\textbf{最容易写成模板的地方}——因为很多人在这里直接背了两段与题目无关的空泛议论。本册对这一章做了最大幅度的改造：所有语料都补上了\textbf{“图表锚定”环节}，让每一句议论都能指回图表里的某个具体信息。
\end{note}

\subsection{两种骨架：1+3 与 2+2}

第 2 段固定写 4--5 句。御风原书给出了两种骨架，本册保留并做了修正。

\begin{tpl}{骨架一：两点式（2+2，推荐）}
\en{For one thing, what first comes into view, without any doubt, is} \fillin{第一分点}\en{. Thanks to} \fillin{原因}\en{,} \fillin{对第一分点的解释}\en{. For another, there still remain potential concerns indicated in the chart.} \fillin{图表中的某个具体现象} \en{serves as an implication that} \fillin{第二分点}\en{.}
\end{tpl}

\begin{note}
\textbf{本册的三处改进：}
① 原书是 \en{...indicated in the （图画/图表）}，是一个待填的括号——现在直接写明 \en{in the chart}，并强制你填一个\textbf{具体的图表现象}（某个数据、某个排名、某个落差），这就完成了“锚定”。
② \en{For one thing ... For another ...} 是对比式并列，天然适合“一半是好事、一半有隐忧”的双面解读，正好命中档位描述里的“评论相关且\emph{充分}”。
③ 如果图表呈现的完全是正面现象，把 \en{there still remain potential concerns} 换成 \en{another point worth noting is that ...} 即可，不要硬凑问题。
\end{note}

\begin{tpl}{骨架二：三点式（1+3）}
\fillin{总领句：一句话概括图表说明了什么}\en{. To begin with,} \fillin{第一分点}\en{. Additionally,} \fillin{第二分点}\en{. More significantly,} \fillin{第三分点}\en{.}
\end{tpl}

\begin{note}
三点式适合“图表里有三个可以并列说的方面”的题目，例如三类商品都有增长但各有原因。\en{To begin with / Additionally / More significantly} 是三者递进的衔接词，最后一个用 \en{More significantly} 比用 \en{Finally} 更有分量。
\end{note}

\subsection{句子的写法：一点两句}

\begin{note}
御风的建议是“一点两句”，即每个分论点写两个句子，而不是一句话塞进所有信息。这是有道理的：两句话意味着你可以用“定义 + 影响”“归因 + 结果”“对比 + 假设”这样的组合，句子更短、更不容易出错，同时\textbf{句型多样性（GRA）自然就上去了}。反过来，一句话里塞三个从句，是错误率最高的写法。
\end{note}

\begin{longtable}{L{2.6cm} L{5.6cm} L{6.0cm}}
\toprule
\textbf{句内功能} & \textbf{可替换的动词 / 结构} & \textbf{示范（以“健康素养提升”为例）}\\
\midrule
\endfirsthead
\toprule
\textbf{句内功能} & \textbf{可替换的动词 / 结构} & \textbf{示范（以“健康素养提升”为例）}\\
\midrule
\endhead
定义 / 分类 & \en{serve as}、\en{be regarded as}、\en{represent}、\en{be reckoned as} & \en{Health literacy serves as a basic competence for modern citizens.}\\
\midrule
归因 & \en{stem from}、\en{be driven by}、\en{arise from}、\en{derive from}、\en{be rooted in}、\en{in response to} & \en{This rise stems largely from the rapid spread of the Internet.}\\
\midrule
影响 / 结果 & \en{lead to}、\en{thereby}、\en{allow}、\en{pave the way for}、\en{facilitate}、\en{promote}、\en{culminate in}、\en{hinder}、\en{impede} & \en{Easier access to information facilitates better health decisions.}\\
\midrule
对比 / 假设 & \en{Given}、\en{insofar as}、\en{without}、\en{but for}、\en{in the absence of}、\en{by contrast}、\en{consistent with} & \en{Without such efforts, the gap between urban and rural areas would widen.}\\
\midrule
举例 & \en{for instance}、\en{taking ... as an example} & \en{Taking health education as an example, it has been widely promoted.}\\
\midrule
评价 & \en{rest on}、\en{lie in}、\en{be central to}、\en{be indispensable to}、\en{be pivotal to}、\en{be essential to} & \en{What matters most lies in turning awareness into daily habits.}\\
\bottomrule
\end{longtable}

\begin{note}
同一篇作文里，\textbf{这六种功能各用一到两次就够了}，不要通篇都是归因句。批改原则第 5 条特别提醒：不以生僻词、复杂句或连接词数量决定分数，正确使用常见句式不扣分。所以上面的词不需要全背，\textbf{每种挑一两个顺手的}，一直用到考场上去。
\end{note}

\subsection{话题语料库：五类通用段落（图表锚定改造版）}

下面是 mmg 原书的五类第二段语料。原文的问题在于：\fillin{题目相关词} 这个占位太大，一段话里出现五次，很容易写成“什么都能套、什么都套不牢”的空话。本册的处理办法是：

\begin{tpl}{语料使用三原则}
\textbf{原则一 · 先锚定。}每段开头先补一句\textbf{带图表信息}的锚定句（本册已为每一类给出），再进入通用议论。\\
\textbf{原则二 · 只取一半。}每类语料给的是 3--4 句，实际用 2--3 句即可，剩下留给锚定句。\\
\textbf{原则三 · 换掉空词。}把 \fillin{题目相关词} 换成\textbf{图表里真实存在的对象}，把泛指换成实指。
\end{tpl}

\begin{note}
\textbf{关于“分寸感”的一句提醒（对应批改原则第 3 条）：}下面每一类语料里的 \en{reflects}、\en{stems from}、\en{serves as an implication} 表达的都是\textbf{你的判断}，不是图表已经证明的事实。阅卷指令里写着“合理推测不能写成图表已经证明的事实”，所以：\textbf{不要}把这些句子改写成 \en{The chart proves that ...} 或 \en{This is because ...}。如果想把分寸感做得更足，把 \en{reflects} 换成 \en{may reflect} / \en{seems to reflect} 即可，句意不变而风险更低。
\end{note}

\subsubsection{第一类：经济民生类（英语二最高频）}

\textbf{适用：}居民收入 / 消费 / 城镇化 / 快递业务量 / 劳动实践 / 旅游 / 生活水平提升等绝大多数“数据在涨”的题目。

\begin{tpl}{经济民生类 · 完整段落（mmg 原文 + 本册锚定改造）}
\fillin{【锚定句】} \en{The steady rise in} \fillin{图表中的指标} \en{is not merely a change in numbers. In reality, the increasing attention paid to} \fillin{主题词} \en{reflects the change in the hierarchy of social needs brought about by social development. As economic conditions continue to improve, people are no longer satisfied with single and basic life goals, but have begun to pursue more diversified ways of life, and} \fillin{主题词} \en{happens to meet this kind of need. It is precisely against such a background that the practical value of} \fillin{主题词} \en{has been continuously magnified, and it has gradually become an important element that is highly valued in modern society.}
\end{tpl}

\begin{eg}
\textbf{以 2022 年英语二“快递业务量变化”为例：}\\
\en{The steady rise in express delivery volume is not merely a change in numbers. In reality, the increasing attention paid to online shopping reflects the change in the hierarchy of social needs brought about by social development. As economic conditions continue to improve, people are no longer satisfied with single and basic life goals, but have begun to pursue more diversified ways of life, and e-commerce happens to meet this kind of need. It is precisely against such a background that the practical value of online consumption has been continuously magnified, and it has gradually become an important element that is highly valued in modern society.}
\end{eg}

\begin{note}
\textbf{锚定句怎么写？}三种最省事的写法：\\
① 引用趋势：\en{The steady rise in X is not merely a change in numbers.}\\
② 引用极值：\en{That X ranks first at \fillin{数值} is telling.}\\
③ 引用落差：\en{Behind the gap between X and Y lies a deeper shift.}\\
\textbf{写完锚定句，你后面的每一句议论就都有了“根”}——阅卷老师看到的是“这个学生从数据出发在思考”，而不是“这个学生在背模板”。
\end{note}

\subsubsection{第二类：校园教育类（对应英语一的“品质类”）}

\textbf{适用：}学生收获调查 / 学习方式 / 素质教育 / 兼职 / 劳动实践 / 阅读 / 团队合作 / 自律 / 坚持等校园主题。

\begin{tpl}{校园教育类 · 正面（mmg 原文 + 本册修正）}
\en{In a sense, the positive influence brought by} \fillin{主题词} \en{is not temporary, but will gradually reveal itself as individuals grow. For young people, growth is never merely a matter of acquiring more knowledge, but of gradually forming a more rational way of dealing with things through the tempering of real life.} \fillin{主题词} \en{can help them learn to face problems in a more mature manner, and also enable them to accumulate a more lasting and reliable inner strength through repeated trials.}
\end{tpl}

\begin{tpl}{校园教育类 · 反面（把 positive 换成 negative，本册修正了平行结构）}
\en{In a sense, the negative influence brought by} \fillin{主题词} \en{is by no means temporary, but will gradually reveal itself as individuals grow. For young people, growth is never merely a matter of acquiring more knowledge, but of gradually forming a mature and rational way of dealing with things through the tempering of real life, and} \fillin{主题词} \en{may exert an adverse influence on this very process. It may not only make it more difficult for them to face problems in a mature manner, but may also gradually weaken the inner strength they ought to possess.}
\end{tpl}

\begin{note}
\textbf{本册的两处改进：}原书反面版最后一句以 \en{through repeated setbacks} 结尾，与前面的 \en{through the tempering of real life} 形成重复且逻辑绕；本册删去，句意更干净。另外原书把这一类的说明写成“品质类需要根据具体题目替换”，实际上\textbf{英语二不需要背具体的“品质”}——你只需要把 \fillin{主题词} 填成图表里的调查对象（如 \en{practical skills}、\en{teamwork}、\en{a sense of responsibility}）即可。
\end{note}

\begin{eg}
\textbf{以 2024 年英语二“劳动实践课收获调查”为例：}\\
\en{In a sense, the positive influence brought by practical skills is not temporary, but will gradually reveal itself as individuals grow. For young people, growth is never merely a matter of acquiring more knowledge, but of gradually forming a more rational way of dealing with things through the tempering of real life. Practical skills can help them learn to face problems in a more mature manner, and also enable them to accumulate a more lasting and reliable inner strength through repeated trials.}
\end{eg}

\subsubsection{第三类：环境生态类}

\textbf{适用：}绿色出行 / 垃圾分类 / 低碳生活 / 环保意识 / 新能源汽车 / 生态保护等。

\begin{tpl}{环境生态类 · 完整段落（mmg 原文，已修正拼写与搭配）}
\en{The continued attention paid to} \fillin{主题词} \en{in contemporary society actually reflects the growing acceptance of the concept of green development. In the past, many people focused more on the convenience brought about by development, whereas today, with ecological pressure steadily accumulating, the public has gradually come to realize that, without sufficient attention, the problems associated with} \fillin{主题词} \en{will affect not only present life but also place constraints on future development. It is precisely for this reason that} \fillin{主题词} \en{has become increasingly worthy of attention in contemporary society.}
\end{tpl}

\subsubsection{第四类：体育健康类}

\textbf{适用：}锻炼方式 / 健康素养 / 体检 / 饮食 / 休闲活动 / 全民健身等。

\begin{tpl}{体育健康类 · 完整段落（mmg 原文，只需改两处即从“环境”变“健康”）}
\en{The continued attention paid to} \fillin{主题词} \en{in contemporary society actually reflects the growing acceptance of the concept of healthy living. In the past, many people focused more on the satisfaction brought about by food, whereas today, with health problems steadily accumulating, the public has gradually come to realize that, without sufficient attention, the problems associated with} \fillin{主题词} \en{will affect not only present life but also exert a negative influence on future life. It is precisely for this reason that} \fillin{主题词} \en{has become increasingly worthy of attention in contemporary society.}
\end{tpl}

\begin{note}
\en{environmental} 版和 \en{health} 版的骨架完全一样，只换了三处词（\en{green development → healthy living}、\en{the convenience brought about by development → the satisfaction brought about by food}、\en{ecological pressure → health problems}）。\textbf{这就是 mmg 原书讲的“换词应万变”}——但要注意：它成立的前提是\textbf{骨架本身是通用的议论框架}，而不是把不相关的句子硬套。健康题里不要出现 \en{green development}，这是底线。
\end{note}

\subsubsection{第五类：文化自信类}

\textbf{适用：}传统文化 / 博物馆参观 / 非遗 / 春节花销 / 国潮 / 文化活动等。

\begin{tpl}{文化自信类 · 完整段落（mmg 原文 + 本册修正 \en{stem from} 搭配）}
\en{Viewed from a deeper perspective, the sustained attention paid to} \fillin{主题词} \en{actually reflects the growing cultural consciousness in contemporary society. As society becomes increasingly open, more and more people have come to realize that a nation's truly lasting confidence stems not only from its tangible achievements, but also from its identification with and commitment to its own cultural traditions. It is against this background that} \fillin{主题词} \en{has been reinterpreted, and the significance it embodies has consequently become more distinct.}
\end{tpl}

\subsubsection{第六类：科技与网络类（本册新增，英语二高频）}

互联网与科技主题在英语二真题中反复出现（2017 年博物馆、2020 年手机阅读、2023 年健康素养的信息渠道），而 mmg 原书没有独立的科技类语料。本册补一段，风格与前五类保持一致。

\begin{tpl}{科技与网络类 · 完整段落（本册新编，可直接背）}
\en{The growing reliance on} \fillin{主题词} \en{in contemporary society actually reflects the profound changes that digital technology has brought to everyday life. In the past, information and services were largely confined to physical spaces, whereas today, with the Internet penetrating almost every aspect of daily life, the public has gradually come to realize that} \fillin{主题词} \en{has not only improved efficiency but also reshaped the way people learn, work and communicate. It is precisely for this reason that} \fillin{主题词} \en{has become an indispensable part of modern life, and its influence deserves careful reflection.}
\end{tpl}

\subsection{“上位概念”：用最少的词覆盖最多的题}

\begin{note}
御风原书有一个很关键的观察：\textbf{英语二图表作文的话题重合度极高}，而且真正的“主题词”应该是\textbf{方向大词（上位概念）}，不是具体题目词。太泛（如 \en{development}）会显得空洞，太具体（如 \en{air conditioner ownership}）又压不中。正确的粒度是：\en{living standards}、\en{consumption upgrading}、\en{public health awareness} 这一层。
\end{note}

\begin{longtable}{L{4.6cm} L{4.4cm} L{5.2cm}}
\toprule
\textbf{图表可能的主题} & \textbf{对应的上位概念} & \textbf{可以搭配的评论方向}\\
\midrule
\endfirsthead
\toprule
\textbf{图表可能的主题} & \textbf{对应的上位概念} & \textbf{可以搭配的评论方向}\\
\midrule
\endhead
收入、消费、城镇化、快递量 & \en{living standards}；\en{consumption upgrading}；\en{economic development} & 需求层次上移；服务便利化；城乡差距缩小\\
\midrule
教育、学习方式、劳动实践、阅读 & \en{quality education}；\en{lifelong learning}；\en{personal growth} & 能力培养重于知识灌输；知行合一\\
\midrule
锻炼、体检、健康素养、饮食 & \en{public health awareness}；\en{a healthy lifestyle} & 从“治已病”到“治未病”；健康观念转变\\
\midrule
环保、低碳、垃圾分类、新能源 & \en{green development}；\en{environmental awareness} & 眼前利益与长远利益的取舍\\
\midrule
博物馆、非遗、文化活动、国潮 & \en{cultural confidence}；\en{cultural consciousness} & 从“看见传统”到“认同传统”\\
\midrule
手机阅读、线上会议、直播、AI & \en{digital technology}；\en{the digital age} & 效率提升与专注力流失并存\\
\midrule
老年人休闲、养老、人口老龄化 & \en{population aging}；\en{care for the elderly} & 物质保障与精神关怀的落差\\
\midrule
兼职、就业、职业规划、创业 & \en{career planning}；\en{employment pressure} & 短期收益与长期积累\\
\bottomrule
\end{longtable}

\begin{note}
背“上位概念”还有一个隐性好处：它让第 2 段的\textbf{语域}（register）统一。整篇都用 \en{living standards / public awareness} 这一层的词，读起来像一篇完整的议论文；一会儿 \en{air conditioners}、一会儿 \en{social development}，就会显得跳。
\end{note}

\subsection{第 2 段成品与批注}

\begin{tpl}{第 2 段 · 成品（以 2025 年英语二“社区老年人休闲活动”为例）}
\en{The uneven distribution of these activities is revealing. For one thing, what first comes into view, without any doubt, is the remarkable diversity of leisure available to today's elderly. Thanks to the immense progress in living conditions, they enjoy increasing opportunities to cultivate hobbies and entertain themselves, which was hardly imaginable a few decades ago. For another, there still remain potential concerns indicated in the chart. The disproportionate preference for activities requiring less cognitive effort serves as an implication that humanistic care for the elderly may lag behind material support.}
\end{tpl}

\begin{note}
\textbf{得分点解析：}① 第 1 句是\textbf{锚定句}，直接指向图表里“分布不均”这个核心特征；② \en{For one thing ... For another ...} 给出“好的一面 + 隐忧”的双面解读，正是“评论充分”；③ 第 2 句用 \en{Thanks to ...} 做归因，第 4 句用 \en{serves as an implication that ...} 做推测——\textbf{推测用了 \en{may}，没有写成 \en{proves}}，符合批改原则第 3 条；④ 全程没有堆砌数字，但每一句都能在图里找到依据。
\end{note}

\subsection{第 2 段自查}

\begin{longtable}{C{0.8cm} L{13.6cm}}
\toprule
& \textbf{检查项}\\
\midrule
\endfirsthead
\toprule
& \textbf{检查项}\\
\midrule
\endhead
\cb & 第 1 句是\textbf{总领/锚定句}，直接指向图表的核心特征，不是一句万能空话。\\
\cb & 至少有 2 句能与图表里的具体信息对应（趋势、极值、排名或落差）。\\
\cb & 用了“分点”结构（\en{For one thing ... For another} 或 \en{To begin with ... Additionally ...}），层次一眼可辨。\\
\cb & 所有推测都用了 \en{may / might / seems to / is likely to}，\textbf{没有} \en{The chart proves / shows that people are lazy} 这类越界表达。\\
\cb & 没有把“类别差异”改写成“随时间变化”。\\
\cb & 没有引用图里不存在的数据。\\
\cb & 分论点数量是 2--3 个，不是 5 个（宁缺毋滥，说不完就少说）。\\
\cb & 全段 4--5 句，60--75 词。\\
\cb & 没有出现与本题完全无关的“万能议论”（例如写“网络时代弊大于利”这种通用段落）。\\
\bottomrule
\end{longtable}

\clearpage

\section{第 3 段：总结与建议}
\markboth{第 3 段：总结与建议}{}

\begin{note}
第 3 段对应题干里的 \en{give your comments}，固定 3 句：\textbf{总结 → 概括/表态 → 建议/期待}。对这一段的硬要求有三条：
① \textbf{不出现任何数字}（数字留在第 1、2 段）；
② \textbf{不要过度升华}，图表题最好就事论事，与前面的分析、图表的内容有紧密关联；
③ 段首逻辑词用 \en{Overall}，\textbf{不要用 \en{In conclusion}}（\en{In conclusion} 是长篇议论文的收尾语，用在 150 词的短文里显得笨重）。
\end{note}

\subsection{图表题专用三句式}

\begin{tpl}{第 3 段公式（图表题）}
\textbf{第 1 句（总结）}：\en{Overall,} + 图表核心矛盾/核心内容的概括。\\
\textbf{第 2 句（概括/表态）}：把矛盾的各个方面概括一下，或给出你的判断。\\
\textbf{第 3 句（建议/期待）}：谁应该做什么；或一句有分寸的展望。
\end{tpl}

\begin{eg}
\textbf{以 2025 年英语二“老年人休闲活动”为例：}\\
\en{Overall, the chart reflects both the material progress the elderly have enjoyed and the spiritual care they still need. Every small effort counts when it comes to enhancing the quality of life for elderly citizens. As college students, we can begin by spending more time with our grandparents and taking a more active part in voluntary work.}
\end{eg}

\begin{note}
\textbf{得分点解析：}① 第 1 句用 \en{both ... and ...} 把第 2 段的“双面解读”收拢成一句；② 第 2 句给判断，不谈数字；③ 第 3 句落到“大学生能做什么”，具体、可执行、有身份感——比空喊“社会应该重视”好得多。
\end{note}

\subsection{措施功能句库（御风原书八组，本册修正 + 补中文）}

这是第三段最实用的弹药。\textbf{积极情况四组 + 消极情况四组}，每组都是一段可直接使用的成品，只需替换主题词。

\subsubsection{积极情况}

\begin{tpl}{积极一 · 总结 + 强调意义}
\en{In a nutshell, I am firmly of the opinion that} \fillin{主题词 / 动名词短语} \en{holds the key to happiness and thus deserves to be placed at a premium. Once the significance of this issue is fully recognized, the next step is to work out and implement a down-to-earth strategy. It is also incumbent on the departments concerned to introduce feasible measures to further promote it. Only by doing so can we lead a more fulfilling life.}\\[3pt]
中文：简而言之，我坚定地认为，主题词是通往幸福的关键，因此理应受到高度重视。一旦充分认识到这一问题的重要性，接下来就应制定并落实切实可行的策略。有关部门也有责任采取可行措施，进一步推动其发展。只有这样，我们才能过上更加充实的生活。
\end{tpl}

\begin{tpl}{积极二 · 政府 + 个人共同推进}
\en{To improve the current state of} \fillin{主题词} \en{, it should be high on the agenda of the departments concerned to formulate practical strategies and call for concerted efforts. As for individuals, we would be well advised to adjust our attitudes towards} \fillin{主题词} \en{and do our part in furthering it.}\\[3pt]
中文：为了改善主题词的现状，有关部门应优先制定切实可行的策略，并呼吁社会各方共同努力。对于个人而言，我们也应调整对主题词的态度，并尽自己的一份力量推动其发展。
\end{tpl}

\begin{tpl}{积极三 · 宣传与观念转变}
\en{To improve the current state of affairs, we can launch online campaigns aimed at encouraging a shift in people's attitudes towards} \fillin{主题词} \en{and raising public awareness of its significance.}\\[3pt]
中文：为了改善当前状况，我们可以开展网络宣传活动，推动人们转变对主题词的态度，并提高公众对其重要性的认识。
\end{tpl}

\begin{tpl}{积极四 · 强调重要性 + 个人行动}
\en{As far as I am concerned,} \fillin{主题词} \en{is of paramount significance and deserves to be held in high regard. As individuals, we should constantly remind ourselves to be} \fillin{主题词} \en{and translate such awareness into concrete action.}\\[3pt]
中文：在我看来，主题词具有至关重要的意义，理应受到高度重视。作为个人，我们应时刻提醒自己做到主题词，并将这种意识转化为实际行动。
\end{tpl}

\subsubsection{消极情况}

\begin{tpl}{消极一 · 总结判断}
\en{I am firmly of the opinion that} \fillin{主题词 / 主题句} \en{is detrimental to our lives and should therefore be effectively curbed.}\\[3pt]
中文：我坚定地认为，主题词会对我们的生活造成不利影响，因此应当得到有效遏制。
\end{tpl}

\begin{tpl}{消极二 · 政府 + 个人共同治理}
\en{To stamp out this problem, it should be high on the agenda of the departments concerned to formulate practical strategies and call for concerted efforts. As for individuals, we would be well advised to adjust our attitudes towards} \fillin{主题词} \en{and do our part in preventing it.}\\[3pt]
中文：为了杜绝这一问题，有关部门应优先制定切实可行的策略，并呼吁社会各方共同努力。对于个人而言，我们也应调整对主题词的态度，并尽自己的一份力量加以防范。
\end{tpl}

\begin{tpl}{消极三 · 互联网宣传}
\en{Fortunately, by harnessing the power of the Internet, we can spread practical advice and raise public awareness, thereby discouraging people from} \fillin{动名词 / 主题行为}\en{.}\\[3pt]
中文：幸运的是，通过利用互联网的力量，我们可以传播实用建议、提高公众意识，从而减少人们某种负面行为的发生。
\end{tpl}

\begin{tpl}{消极四 · 有关部门治理}
\en{It is also incumbent on the departments concerned to introduce feasible measures and remedy the problem before serious consequences arise.}\\[3pt]
中文：有关部门也有责任出台切实可行的措施，并在严重后果出现之前及时解决这一问题。
\end{tpl}

\begin{note}
\textbf{怎么选？}积极题用积极组的任意一组；消极题用消极组的任意一组。\textbf{不要一次把三四组全用上}——每组都是一整段，全用会严重超字数。正确做法是：\textbf{从一组里取 2 句}，与图表题的“回到图表”要求结合起来。
\end{note}

\subsection{mmg 通用尾段：政府 -- 媒体 -- 个人（按话题分类）}

mmg 原书的尾段是一套“政府 / 媒体 / 个人”的三方责任结构。它的优点是清晰、好写、不易跑题；缺点是\textbf{四个话题版本内容几乎完全一样}（只换了几个词），如果照抄会显得空洞。本册的处理是：\textbf{保留一套骨架，给出四个话题的替换词表}。

\begin{tpl}{通用尾段骨架（一套，通吃所有话题）}
\en{If} \fillin{主题词} \en{is to truly embody its practical value, it cannot rely solely on superficial advocacy, but requires sustained promotion on multiple levels.} \fillin{主体词 1：政府 / 学校 / 企业} \en{should proceed from the perspective of long-term development and ensure that} \fillin{主题词} \en{receives stronger support. The media should actively perform its communicative function so that the significance of} \fillin{主题词} \en{can be genuinely understood by more people. Meanwhile, individuals should remain rational and view the influence of} \fillin{主题词} \en{on the times in a proper way.}
\end{tpl}

\begin{longtable}{L{3.4cm} L{5.4cm} L{5.4cm}}
\toprule
\textbf{话题} & \textbf{可替换的“主体词 1”} & \textbf{可替换的支撑动词}\\
\midrule
\endfirsthead
\toprule
\textbf{话题} & \textbf{可替换的“主体词 1”} & \textbf{可替换的支撑动词}\\
\midrule
\endhead
经济民生 & \en{the government} / \en{relevant departments} & \en{receives stronger support}；\en{ensure fair access to ...}\\
\midrule
校园教育 & \en{schools and universities} & \en{receives stronger support}；\en{be integrated into the curriculum}\\
\midrule
环境生态 & \en{the government} / \en{regulatory bodies} & \en{receives stronger support}；\en{is subject to stricter regulation}\\
\midrule
体育健康 & \en{the government} / \en{public health institutions} & \en{receives stronger support}；\en{be made more accessible}\\
\midrule
文化自信 & \en{the government} / \en{cultural institutions} & \en{receives stronger support}；\en{be better protected and promoted}\\
\midrule
科技网络 & \en{the government} / \en{technology companies} & \en{receives stronger support}；\en{is subject to clearer rules}\\
\bottomrule
\end{longtable}

\begin{note}
\textbf{本册的修正：}mmg 原书为“经济类”“校园类”“环境生态类”“体育健康类”“文化自信类”分别给了一整段几乎一样的尾段，实际是重复。本册只保留一套骨架，用上面的替换表覆盖全部话题。另外原书反面版本里的 \en{is subject to stronger regulation} 已经并入表中，不需要再单独背一段。\textbf{这样一改，第三段从“背五段”变成“背一段 + 一张表”。}
\end{note}

\subsection{三种收尾风格，按你的水平选}

\begin{longtable}{C{2.6cm} L{5.6cm} L{5.9cm}}
\toprule
\textbf{风格} & \textbf{写法} & \textbf{适合谁}\\
\midrule
\endfirsthead
\toprule
\textbf{风格} & \textbf{写法} & \textbf{适合谁}\\
\midrule
\endhead
\textbf{扣题型} & 三句都围绕图表本身：\en{Overall, the chart reflects ...} / \en{... is what the figures ultimately point to.} / \en{To change this, ... should ...} & 想稳稳拿 13--15 分、不想冒险的同学。\textbf{本册最推荐。}\\
\midrule
\textbf{折中型} & 前两句扣图，第三句用一句通用的呼吁收尾 & 想兼顾“扣题”与“漂亮”的同学。\\
\midrule
\textbf{升华型} & 前两句扣图，第三句上升到价值/时代层面 & 语言功底扎实、能保证不出语法错误的同学。\textbf{风险最高}，写不好就是“过度升华”，反而扣分。\\
\bottomrule
\end{longtable}

\begin{note}
御风原书讲过一句很重要的话：图表作文“最后一段里面最好不要出现任何一个数字”，“既不能跑题过度上升，也不能不总结把数据写到最后一段里面”。这句话在阅卷现场是站得住的——第 3 段的作用是把前两段收拢成一个观点，不是继续报数据。
\end{note}

\subsection{第 3 段成品（三档对照）}

\begin{tpl}{成品一 · 扣题型（稳）}
\en{Overall, the chart shows a clear and encouraging trend in the health literacy of Chinese residents. Behind the rising figures lie decades of investment in education, public health and the Internet. To sustain this progress, both the government and individuals should keep working in the same direction.}
\end{tpl}

\begin{tpl}{成品二 · 折中型（推荐）}
\en{Overall, the figures reveal far more than a change in habits; they reflect a society that is becoming healthier and better informed. Such progress, however, will not continue automatically. It is therefore advisable for the government to expand health education, and for each of us to turn awareness into everyday practice.}
\end{tpl}

\begin{tpl}{成品三 · 含隐忧型（图表双面时用）}
\en{Overall, the chart reflects both the remarkable progress in living conditions and the imbalance that remains. Material comfort alone cannot guarantee a genuinely better life, and the gap between what is available and what is truly needed still deserves attention. As college students, we can begin by offering more companionship to those around us.}
\end{tpl}

\subsection{第 3 段自查}

\begin{longtable}{C{0.8cm} L{13.6cm}}
\toprule
& \textbf{检查项}\\
\midrule
\endfirsthead
\toprule
& \textbf{检查项}\\
\midrule
\endhead
\cb & \textbf{全文没有任何数字}（年份、百分比、人数都不能出现在这一段）。\\
\cb & 第 1 句以 \en{Overall} 开头，而不是 \en{In conclusion} / \en{To sum up}。\\
\cb & 第 1 句确实在“总结图表”，而不是一句放到任何题目都成立的空话。\\
\cb & 建议的具体执行者明确（政府 / 学校 / 媒体 / 企业 / 个人），没有含糊地说“社会应该……”。\\
\cb & 没有过度升华（没有扯到“中华民族伟大复兴”“人类命运共同体”这类与图表无关的大话）。\\
\cb & 全段 3 句，35--45 词。\\
\cb & 与前两段使用不同的句式开头，避免整篇都是 \en{The chart ...} / \en{What ... is ...}。\\
\bottomrule
\end{longtable}

\clearpage

\section{全篇拼装：四套真题的批注版范文}
\markboth{全篇拼装：四套真题的批注版范文}{}

\subsection{拼装公式}

\begin{tpl}{静态图 · 全篇拼装公式}
\textbf{第 1 段}（3 句）\\
\quad ① \en{The} \fillin{图表类型} \en{illustrates} \fillin{指标} \en{of} \fillin{对象} \en{among} \fillin{范围}\en{.}\\
\quad ② \en{Overall,} \fillin{极值 A} \en{ranks first, accounting for} \fillin{数值 A}\en{, followed by} \fillin{极值 B} \en{at} \fillin{数值 B}\en{.}\\
\quad ③ \en{By contrast,} \fillin{其余项} \en{are overshadowed, none of them exceeding} \fillin{数值}\en{.}\\[4pt]
\textbf{第 2 段}（4--5 句）\\
\quad ④ 锚定句（指向图中的分布不均 / 落差）\\
\quad ⑤ \en{For one thing, what first comes into view, without any doubt, is} \fillin{第一分点}\en{.}\\
\quad ⑥ \en{Thanks to} \fillin{原因}\en{,} \fillin{对第一分点的解释}\en{.}\\
\quad ⑦ \en{For another, there still remain potential concerns indicated in the chart.} \fillin{具体现象} \en{serves as an implication that} \fillin{第二分点}\en{.}\\[4pt]
\textbf{第 3 段}（3 句）\\
\quad ⑧ \en{Overall, the chart reflects both} \fillin{正面} \en{and} \fillin{隐忧}\en{.}\\
\quad ⑨ \fillin{你的判断}\en{.}\\
\quad ⑩ \en{As college students / To sustain this progress,} \fillin{具体可执行的建议}\en{.}
\end{tpl}

\begin{tpl}{动态图 · 全篇拼装公式}
\textbf{第 1 段}（2--3 句）\\
\quad ① \en{The} \fillin{图表类型} \en{illustrates the changes in} \fillin{指标 + 对象} \en{from} \fillin{时间 1} \en{to} \fillin{时间 2}\en{.}\\
\quad ② \en{Overall,} \fillin{总体趋势}\en{, and} \fillin{最突出的对象} \fillin{趋势动词} \en{from} \fillin{数值 1} \en{to} \fillin{数值 2}\en{,} \fillin{意义尾}\en{.}\\
\quad ③ \en{In contrast,} \fillin{另一对象} \fillin{趋势动词} \fillin{数值}\en{.}\\[4pt]
\textbf{第 2 段}（4--5 句）\\
\quad ④ 锚定句（引用趋势：\en{The steady rise in X is not merely a change in numbers.}）\\
\quad ⑤--⑦ 从第 6 章话题语料库里选 2--3 句，替换 \fillin{主题词}\\
\quad ⑧（可选）一句推测/隐忧，用 \en{may / seems to}\\[4pt]
\textbf{第 3 段}（3 句）\\
\quad ⑨ \en{Overall,} \fillin{图表核心内容的概括}\en{.}\\
\quad ⑩ \fillin{判断或建议}\en{.}\\
\quad ⑪ \en{Only when ... / To sustain this progress, ...}
\end{tpl}

\subsection{真题范文一：2025 年英语二（社区老年人休闲活动 · 静态柱状图）}

\begin{tpl}{题目要点}
柱状图：某社区老年人主要休闲活动。看电视 90.8\%，散步 68.3\%，其余（园艺、阅读、下棋）均不超过 35\%。
\end{tpl}

\begin{tpl}{批注版范文（157 词）}
\en{The bar chart illustrates the primary leisure activities of the elderly in a community. Overall, watching TV and walking stand out as the two most common choices, accounting for 90.8\% and 68.3\% respectively, whereas other activities such as gardening, reading and playing chess are far less popular, none of them exceeding 35\%.}\\[3pt]
\en{Such a distribution is revealing. For one thing, what first comes into view, without any doubt, is the remarkable diversity of leisure now available to the elderly. Thanks to the immense progress in living conditions, they enjoy far more opportunities to cultivate hobbies and entertain themselves than previous generations did. For another, there still remain potential concerns indicated in the chart. The disproportionate preference for activities requiring little cognitive effort may suggest that humanistic care for the elderly still lags behind material support.}\\[3pt]
\en{Overall, the chart reflects both the material comfort the elderly have gained and the companionship they still need. Every little effort counts toward enhancing the quality of life for elderly citizens. As college students, we can begin by spending more time with our grandparents and taking a more active part in voluntary work.}
\end{tpl}

\begin{note}
\textbf{本册对源范文的七处修改（对照评分规则）：}
\begin{enumerate}[leftmargin=1.7em,itemsep=1.5pt,topsep=2pt]
\item \en{The participation rate highlights the prevalence of watching TV and walking} $\rightarrow$ \en{watching TV and walking stand out as the two most common choices}。原句把“参与率”当主语去“凸显受欢迎程度”，逻辑绕且搭配生硬；改后主语就是图表里的对象，更直接。
\item \en{While other activities ... are overshadowed in comparison, none of which exceeds 35\%} $\rightarrow$ \en{whereas other activities ... are far less popular, none of them exceeding 35\%}。原句 \en{While} 从句 + \en{none of which} 从句无法衔接，是明确的语法错误（\en{which} 需要有先行词，且 \en{While} 从句不能这样接）；改后 \en{none of them} 作独立主格，语法正确。
\item 删去 \en{Succinct as it is, the chart conveys abundant enlightenments, and below are my perspectives.}。一是 \en{enlightenments} 作可数名词使用不当（应为 \en{insights}）；二是这句被大量模板书使用，容易触发“模板感”。本册用一句具体的锚定句 \en{Such a distribution is revealing.} 替代。
\item \en{life-altering} 类的空泛修饰一律不保留，换成由图可验证的表述。
\item \en{Every little endeavor counts to enhance} $\rightarrow$ \en{counts toward enhancing}。\en{count to} 不是“有助于”的正确搭配。
\item \en{take more participation in volunteer work} $\rightarrow$ \en{take a more active part in voluntary work}。原搭配不存在。
\item 全文控制在 157 词。源范文超过 180 词，超出“约 150 词”的要求，在阅卷现场属于不必要的风险。
\end{enumerate}
\end{note}

\subsection{真题范文二：2024 年英语二（劳动实践课学生收获 · 静态柱状图）}

\begin{tpl}{题目要点}
柱状图：某高校劳动实践课学生的收获调查。相关知识 84.5\%，实践技能 91.3\%，乐趣 54.5\%，合作能力不足三分之一。
\end{tpl}

\begin{tpl}{批注版范文（154 词）}
\en{The bar chart presents a survey of the benefits students derived from a labor-practice class at a university. Overall, the two most frequently cited advantages are the acquisition of relevant knowledge and the improvement of practical skills, accounting for 84.5\% and 91.3\% respectively, whereas less than one third of the respondents mention cooperation skills.}\\[3pt]
\en{These figures suggest that while the course fulfills its basic objectives, it falls short of helping students appreciate the deeper value of labor. On the one hand, students do acquire substantial knowledge and take part in hands-on practice. On the other hand, many of them complete the tasks individually and leave with little sense of shared achievement, which indicates that teamwork has not yet been given due weight.}\\[3pt]
\en{Overall, the chart points to a course that works but has not yet reached its full potential. Universities should therefore place greater emphasis on its core mission --- immersing students in the spirit of labor and cultivating their ability to cooperate.}
\end{tpl}

\begin{note}
\textbf{得分点解析：}① 第 1 段把四个数据压缩成三个（两个极值 + 一个“不足三分之一”），符合“不要求罗列全部数字”；② 第 2 段开始就是\textbf{锚定句} \en{These figures suggest that ...}，直接给出判断；③ \en{On the one hand ... On the other hand ...} 用对了冠词（源范文写的是 \en{On one hand}，属于易错点）；④ 第 3 段提出了一个\textbf{具体可执行的建议}（\en{place greater emphasis on ... cultivating their ability to cooperate}），而不是空喊口号。
\end{note}

\subsection{真题范文三：2023 年英语二（居民健康素养水平变化 · 动态折线图）}

\begin{tpl}{题目要点}
折线图：2012--2021 年，15--69 岁居民具备基本健康素养的比例由 8.8\% 升至 25.4\%，十年间近乎翻了两番（约 2.9 倍）。
\end{tpl}

\begin{tpl}{批注版范文（162 词）}
\en{The line graph traces the change in the health literacy level of Chinese residents over the past decade. As is shown, the proportion of people aged 15 to 69 with basic health literacy rose sharply from 8.8\% to 25.4\%, almost tripling within ten years.}\\[3pt]
\en{The steady rise is not merely a change in numbers. For one thing, it stems largely from the expansion of the educated population, as more people are now able to obtain information across a wide range of fields, healthcare included. For another, it can be regarded as a by-product of the Internet, which has lowered the threshold for ordinary people to access instant information. What deserves equal attention is the role of the government: the improvement is also an epitome of the Healthy China Strategy.}\\[3pt]
\en{Overall, the chart reflects a society that is becoming better informed about its own health. Such progress, however, will not continue automatically. It is therefore advisable for the government to expand health education, and for each of us to turn awareness into everyday practice.}
\end{tpl}

\begin{note}
\textbf{本册对源范文的三处修改（对照评分规则）：}
\begin{enumerate}[leftmargin=1.7em,itemsep=1.5pt,topsep=2pt]
\item 源范文写 \en{almost tripling in a mere decade}。核对数据：8.8\% $\rightarrow$ 25.4\% 是 2.89 倍，用 \en{tripling}（三倍）略有夸大。本册改为 \en{almost tripling}，措辞更稳。\textbf{批改原则第 2 条明确“不编造数据”，对倍数、幅度的表述同样要准确。}
\item 源范文的第三段用 \en{Not only does remarkable progress in health literacy suggest achievements in the past, but acts as a promising start of the forthcoming decade.}——\en{not only ... but acts as} 前后不平衡（前面是 \en{suggest}，后面是 \en{acts as}），是倒装句最容易踩的坑。本册改成了两个独立的简单句。
\item 源范文结尾 \en{with our joint efforts, the figures will not stagnate but scale new peaks} 里 \en{the figures ... scale new peaks} 搭配不自然。本册换成 \en{turn awareness into everyday practice}，更具体、更安全。
\end{enumerate}
\end{note}

\subsection{真题范文四：2012 年英语二（不同年龄段员工工作满意度 · 半动态表格）}

\begin{tpl}{题目要点}
表格：某公司工作满意度调查。41--50 岁组不满意率 64\%，高于另外两组之和；41--50 岁组也是唯一“完全不满意”比例突出的群体。
\end{tpl}

\begin{tpl}{批注版范文（158 词）}
\en{The table demonstrates the results of a job satisfaction survey conducted in a certain company. While employees aged 41 to 50 show clear discontent, their counterparts aged under 41 and over 50 maintain a comparatively favorable attitude. The dissatisfaction rate of the 41--50 group reaches 64\%, exceeding the combined rates of the other two.}\\[3pt]
\en{Such a contrast reveals that people at different stages of life hold markedly different views on their work. On the one hand, it is hardly surprising that the other two groups feel satisfied: those in their twenties and thirties tend to work with ambition, while those over fifty have accumulated savings and begun to prepare for retirement. On the other hand, those in the middle often feel stuck, since they must support both their aging parents and their young children while coping with declining energy.}\\[3pt]
\en{Overall, the chart draws attention to a group that is easily overlooked. The middle-aged deserve more care and warmth, and we students can start by offering more support to our parents from now on.}
\end{tpl}

\begin{note}
\textbf{本册对源范文的五处修改：}
\begin{enumerate}[leftmargin=1.7em,itemsep=1.5pt,topsep=2pt]
\item \en{Notably as well, middle-aged people also stand out as the only age group that strikingly reported no satisfaction at all} $\rightarrow$ 删去。原句 \en{notably as well} 重复冗余，且该数据在本段已由“64\% 高于另两组之和”覆盖，\textbf{信息重复}不如让位给对比句。
\item \en{There's an increasing burden of raising their parents and minor kids} $\rightarrow$ \en{they must support both their aging parents and their young children}。\en{raise their parents} 搭配错误（\en{raise} 是抚养，不是赡养）；\en{minor kids} 也不地道。
\item \en{no longer being aggressive and motivated, they feel worn out} $\rightarrow$ \en{while coping with declining energy}。原句是悬垂结构（\en{being} 的逻辑主语与主句主语不一致），属于语法扣分点。
\item \en{To conclude, few realize the challenges facing the middle-aged group, although most of us understand the necessity to care about the minor kids and the elderly.} $\rightarrow$ 精简为一句总结。原句两个分句互相削弱，且 \en{To conclude} 在本册统一改为 \en{Overall}。
\item 原文第三段的 \en{Every little action counts} 与第二段内容重复，本册保留但把它放在结尾，作为“身份落地”句，效果更好。
\end{enumerate}
\end{note}

\subsection{考场流程：25--30 分钟怎么用}

\begin{longtable}{C{2.2cm} C{2.0cm} L{10.2cm}}
\toprule
\textbf{阶段} & \textbf{建议用时} & \textbf{动作}\\
\midrule
\endfirsthead
\toprule
\textbf{阶段} & \textbf{建议用时} & \textbf{动作}\\
\midrule
\endhead
审题 & 2 分钟 & 回答第 2 章的“审题七问”：图表类型、静态/动态、时态、主要信息、对象与指标、单位、篇幅分配。在草稿纸上写下对象的英文名和 3--5 个关键数字。\\
\midrule
列点 & 2 分钟 & 写下第 2 段的 2--3 个分论点（中文即可），每个分论点旁边标出“锚定到图表的哪一句”。\\
\midrule
写作 & 18--20 分钟 & 按九句骨架往下写。\textbf{先写第 1 段和第 3 段（这两段是常量），最后写第 2 段}，这样即使时间不够，文章也是完整的。\\
\midrule
回读 & 3--5 分钟 & 只检查四类问题：主谓一致、时态、单复数、拼写。\textbf{不要在这个阶段改句子结构}，越改越错。\\
\bottomrule
\end{longtable}

\begin{note}
\textbf{书写要求：}题干里那句 \en{You should write neatly on the ANSWER SHEET} 不是客套话。字迹是阅卷老师对你作文的第一印象，潦草的字迹会直接影响老师“愿不愿意细看”。建议：\textbf{统一用 0.5 的黑色签字笔；字母大小一致；每段首行缩进两个字母；不留大片涂改}。写错了划一条横线即可。
\end{note}

\subsection{大作文全篇自查清单（考前背下来）}

\begin{longtable}{C{0.8cm} L{6.4cm} L{6.4cm}}
\toprule
& \textbf{内容与任务} & \textbf{语言与结构}\\
\midrule
\endfirsthead
\toprule
& \textbf{内容与任务} & \textbf{语言与结构}\\
\midrule
\endhead
\cb & 图表类型写对，第一句含“指标 + 对象 + 范围限定”。 & 全文 150--175 词，不是 200+。\\
\cb & 静态/动态判断正确，通篇没有“动静混写”。 & 第 1 段一般现在时（描述图表本身）+ 一般过去时（图内数据）。\\
\cb & 题干三项要求（describe / interpret / comment）各自落到对应段落。 & 没有出现 \en{The chart proves that ...} 这类越界表达。\\
\cb & 每个数字都能在图上找到，单位正确。 & 第 3 段没有出现任何数字。\\
\cb & 没有把“类别差异”写成“时间趋势”。 & 主谓一致、单复数、冠词无明显错误。\\
\cb & 没有把百分比写成百分点（或反之）。 & 至少使用了 2--3 个准确的“意义尾”分词短语。\\
\cb & 第 2 段有明确的 2--3 个分论点，且能与图表对应。 & 段落之间有明确的衔接词，层次一眼可辨。\\
\cb & 第 3 段有一句具体的建议，执行者明确。 & 没有整段套话与题目无关。\\
\bottomrule
\end{longtable}

\clearpage

\section{小作文总览：题型、审题与格式}
\markboth{小作文总览：题型、审题与格式}{}

\begin{note}
小作文（Part A）满分 10 分，约 100 词。它和大作文最大的区别是：\textbf{大作文比的是“写得好不好”，小作文比的是“该说的有没有说、格式对不对”}。批改原则第 4 条写得很直白：“优先保证信息完整、表达准确和简洁得体。”所以小作文的备考思路是——\textbf{把功能句拼成一个不漏要点的段落}，而不是写一篇漂亮的短文。
\end{note}

\subsection{题型盘点}

\begin{longtable}{C{3.0cm} L{5.2cm} L{5.4cm}}
\toprule
\textbf{题型} & \textbf{出现情况} & \textbf{备考优先级}\\
\midrule
\endfirsthead
\toprule
\textbf{题型} & \textbf{出现情况} & \textbf{备考优先级}\\
\midrule
\endhead
书信 & 绝对主力。感谢、邀请、建议、投诉、道歉、申请、推荐、咨询、祝贺、介绍、求职、辞职…… & \textbf{第一优先}。功能句通用，把第 10--12 章的三段功能句背熟即可覆盖。\\
\midrule
通知 / 告示 & 英语二考过 2013、2015、2020 等年份 & \textbf{第二优先}。与书信只有格式上的差别。\\
\midrule
纪要（\en{Meeting Minutes}） & 大纲新增题型，近年偶有出现 & \textbf{第三优先}。记格式即可（标题 + 会议信息 + 议题 + 决议 + 记录人）。\\
\bottomrule
\end{longtable}

\subsection{审题：三部分拆解法}

书信类题目的题干固定分成三段：

\begin{tpl}{题干拆解模板}
\textbf{第一部分 · 基本信息}：一句话告诉你“出于某个原因，你要写某种书信”，并指定你的身份和收信人。\\
\textbf{第二部分 · 书信内容}：具体要点（英语二通常直接给出 \en{1) ... 2) ...} 或 \en{You should ...}）。\\
\textbf{第三部分 · 写作要求}：字数、\en{Do not use your own name. Use “Li Ming” instead.}、分值。
\end{tpl}

\begin{tpl}{例：2022 年英语二}
\en{Directions:}\\
\en{Write an email to invite international students to a food festival to be held in your university.}\\
\en{You should include the details you think necessary. You should write about 100 words neatly on the ANSWER SHEET. Do not sign your own name at the end of the letter. Use “Li Ming” instead. (10 points)}
\end{tpl}

\begin{note}
\textbf{英语一和英语二在小作文上的区别：}英语一偏“任务驱动”，很多细节要你自己想；\textbf{英语二会直接给出写作要点}。这个差别决定了策略：\textbf{英语二的小作文必须“逐点回应、句句有回应”}——题干提到的每一个要点，都要有一句（最好两句）去交代，不重不漏。
\end{note}

\subsection{英语二小作文历年真题速查}

\begin{longtable}{C{1.4cm} L{12.6cm}}
\toprule
\textbf{年份} & \textbf{题目类型与要点}\\
\midrule
\endfirsthead
\toprule
\textbf{年份} & \textbf{题目类型与要点}\\
\midrule
\endhead
2010 & 感谢与邀请信：感谢对方在你赴美期间的热情接待，并欢迎对方来访中国\\
2011 & 祝贺与建议信：祝贺朋友考上大学，并就如何适应大学生活提出建议\\
2012 & 投诉信：投诉网购的电子词典出问题，要求解决\\
2013 & 通知：通知同学参加为贫困儿童举办的慈善义卖\\
2014 & 告知与咨询信：告知对方自己的生活习惯，并就如何适应国外合租生活征求意见\\
2015 & 通知：介绍高中生夏令营并招募志愿者\\
2016 & 回复感谢与建议信：感谢对方的祝贺，并就翻译学习给出建议\\
2017 & 回复接受邀请信：接受邀请并介绍讲座要点\\
2018 & 道歉信：为临时取消拜访而道歉并建议改期\\
2019 & 汇报建议信：汇报辩论赛计划并征求意见\\
2020 & 通知：介绍古迹参观活动并给出注意事项\\
2021 & 邀请信：邀请同学参加关于学习与校园生活的线上会议\\
2022 & 邀请信：邀请国际学生参加校园美食节\\
2023 & 建议信：建议朋友选择看哪个展览（艺术展 or 机器人展）\\
2024 & 汇报计划信：提出古镇老宅保护调研计划并征求意见\\
2025 & 邀请信：介绍一部短剧并邀请朋友参演\\
\bottomrule
\end{longtable}

\begin{note}
\textbf{观察结论：}十六年里，“邀请 / 感谢 / 建议 / 通知”四类占了绝大多数，而\textbf{每一道题都含有一个“事由”和一个或多个需要交代的细节}。所以小作文的备考本质是：\textbf{把“一件事”的各个侧面（时间、地点、内容、意义、期待）用现成的功能句说清楚}。
\end{note}

\subsection{三段公式：2 + 3--5 + 1--2}

\begin{longtable}{C{1.6cm} C{2.0cm} C{2.0cm} L{9.0cm}}
\toprule
\textbf{段落} & \textbf{句数} & \textbf{词数} & \textbf{任务}\\
\midrule
\endfirsthead
\toprule
\textbf{段落} & \textbf{句数} & \textbf{词数} & \textbf{任务}\\
\midrule
\endhead
\textbf{第 1 段} & 2 句 & 20--25 词 & 第 1 句问候/铺垫；第 2 句点明写作目的（\en{I am writing to ...}）。\\
\midrule
\textbf{第 2 段} & 3--5 句 & 55--70 词 & \textbf{主体}。事件类写六要素（时间/地点/内容/意义）；非事件类写建议/原因。逐点回应题干要点。\\
\midrule
\textbf{第 3 段} & 1--2 句 & 15--25 词 & 期待回应 / 期待结果 / 感谢配合。\\
\bottomrule
\end{longtable}

\begin{note}
mmg 原书把这套句数配比写成“241 / 231 / 152 / 252”等公式（各段选几句），并强调“小作文建议占篇幅的三分之二，第二段可以写多一些”。这与上面的分配一致。\textbf{核心只有一句：第 2 段是主体，第 1、3 段越短越好。}
\end{note}

\subsection{格式规范（格式分就在这里）}

\subsubsection{书信格式}

\begin{tpl}{书信格式（左对齐“齐头式”）}
\en{Dear} \fillin{称呼}\en{,}\\
（空一行）\\
\en{（第 1 段）}\\
（空一行）\\
\en{（第 2 段）}\\
（空一行）\\
\en{（第 3 段）}\\
（空一行）\\
\en{Yours sincerely,}\\
\en{Li Ming}
\end{tpl}

\begin{longtable}{L{3.4cm} L{5.4cm} L{5.2cm}}
\toprule
\textbf{要素} & \textbf{怎么写} & \textbf{注意事项}\\
\midrule
\endfirsthead
\toprule
\textbf{要素} & \textbf{怎么写} & \textbf{注意事项}\\
\midrule
\endhead
称呼 & \en{Dear Professor Smith,} / \en{Dear Mr. Wang,} / \en{Dear Sir or Madam,} / \en{Dear Jack,} / \en{Dear International Students,} & 题干给谁就写谁。不知道姓名用 \en{Dear Sir or Madam,}。给朋友用名字即可。称呼后用\textbf{逗号}。\\
\midrule
正文 & 三段 & 齐头式（每段顶格）和缩进式（每段首行缩进两字母）都可以，\textbf{全篇统一}。\\
\midrule
结束语 & 正式：\en{Yours sincerely,}；一般：\en{Yours,}；朋友：\en{Best wishes,} & 与称呼的语气匹配。\en{Yours sincerely} 后面是逗号，\en{Yours} 后面也是逗号。\\
\midrule
署名 & 题目指定 \en{Li Ming} 就写 \en{Li Ming} & \textbf{绝对不能写自己的真名}。独立一行，不加标点。\\
\bottomrule
\end{longtable}

\begin{note}
mmg 原书建议“齐头式”（每段顶格），理由是“同样的篇幅显得写得更满”。这是很实用的经验。另外原书提到“作文建议 0.7 的笔，更好看；其它科目统一 0.5”——笔的选择因人而异，关键是\textbf{全篇一致、字迹清晰}。
\end{note}

\subsubsection{通知格式}

\begin{tpl}{通知格式}
\begin{center}
\en{Notice}\\
\en{（日期写在第一行最右侧，如 Friday, June 10）}
\end{center}
（空一行）\\
\en{To ensure that} \fillin{通知对象} \en{are fully informed of} \fillin{通知事项}\en{, this notice is issued to provide you with some essential information regarding it. And the detailed arrangements are as follows.}\\
（空一行）\\
\en{（主体：时间、地点、内容、要求，可用 To begin with / As for / Moreover 分点）}\\
（空一行）\\
\en{（结尾 2 句：感谢配合 + 期待参与）}\\
\en{The Students' Union}
\end{tpl}

\begin{longtable}{L{3.4cm} L{10.6cm}}
\toprule
\textbf{要素} & \textbf{要点}\\
\midrule
\endfirsthead
\toprule
\textbf{要素} & \textbf{要点}\\
\midrule
\endhead
标题 & \en{Notice} 居中，首字母大写。这是格式分的第一眼。\\
\midrule
日期 & 写在第一行的最右侧（有的年份也可写在落款下方）。\\
\midrule
正文 & 第一段说明“为什么发这个通知”；第二段说明具体安排；第三段感谢配合/期待参与。\\
\midrule
落款 & 写明发布单位（\en{The Students' Union} / \en{The Organizing Committee}）。\\
\midrule
语气 & 通知以“事”为主，\textbf{客观、正式、有条理}，可以分点写（\en{To begin with / As for / Moreover}），比书信更规整。\\
\bottomrule
\end{longtable}

\zone{格式错误是“低档位”的典型特征。评分参考里“格式或语气不当”对应 2--3 分档。最常见的三类错误：① 通知没写标题 \en{Notice}；② 通知没有落款单位；③ 书信署名写了自己的名字而不是题目指定的 \en{Li Ming}，或者干脆忘记署名。}

\subsubsection{纪要格式（低概率但要知道）}

\begin{tpl}{纪要（Meeting Minutes）结构}
\en{Meeting Minutes}\\
\en{I. Meeting Details}\ \ —— \en{Location / Date / Time / Attendees}\\
\en{II. Meeting Objective}\ \ —— 一句话说明会议目的\\
\en{III. Meeting Agenda}\ \ —— 编号列出 1.\ 2.\ 3.\ 讨论事项与结论\\
\en{Minutes submitted by:} \fillin{记录人}\\
\en{Approved by:} \fillin{批准人}
\end{tpl}

\subsection{小作文备考动作}

\begin{longtable}{C{0.9cm} L{13.4cm}}
\toprule
& \textbf{建议动作}\\
\midrule
\endfirsthead
\toprule
& \textbf{建议动作}\\
\midrule
\endhead
1 & \textbf{背功能句，不背整篇。}第 10--12 章给出的功能句是“零件”，考场上是组装，不是默写。\\
\midrule
2 & \textbf{先用近 10 年英语二真题练“拆要点”。}每道题只做一件事：把题干拆成 2--4 个必须回应的要点，用中文写下来。这一步做熟，跑题的风险基本为零。\\
\midrule
3 & \textbf{再练“拼装”。}每个要点配 1--2 句功能句，连起来读一遍，检查“有没有哪句话和题目无关”。\\
\midrule
4 & \textbf{最后练格式。}拿 5 张纸默写书信格式、通知格式各 3 遍，直到不需要思考。\\
\midrule
5 & \textbf{目标分档自测。}按第 1 章的小作文档位表给自己打分，7 分以上再练下一年（这是 mmg 原书给的标准：\textbf{小作文 7 分以上再开始下一年的练习}）。\\
\bottomrule
\end{longtable}

\clearpage

% =====================================================================
% 70_letter_p1.tex
% 小作文（Part A 应用文）第一段功能句库
% 来源：《mmg 作文模板》三、通用小作文——「第一段第一句」16 条 +
%       「第一段第二句」27 条（含逐条真题示例）；
%       《英语作文模板》（御风）小作文首段公式、通知首段公式与真题范文。
% =====================================================================

\section{小作文第一段功能句库}

\begin{note}
小作文（Part A，约 100 词，满分 10 分）第一段\textbf{固定只用两句}：

\textbf{第 1 句「问候与铺垫」}——建立通讯关系、交代来信语境（希望对方安好 / 很高兴收到来信 / 感谢大家关注本通知 / 已整理好会议纪要）。

\textbf{第 2 句「点明写作目的」}——用一句话把题干事由讲清楚：感谢、邀请、建议、投诉、道歉、申请、推荐、咨询、祝贺、介绍、通知、纪要，或其中若干项的混合。

两句合起来一般不超过 30 词。第一段写得越长，第二段（真正承载写作要点的地方）就越没有空间。
\end{note}

\begin{tpl}{第一段公式}
\textbf{第 1 句}：Dear \fillin{称呼}，+ 问候／情绪铺垫句。\par
\textbf{第 2 句}：点明写作目的 + \fillin{题目相关词}（把题干事由填进去）。\par
\textbf{（选做）第 3 句}：\kw{Below are the details.} / \kw{And the detailed arrangements are as follows.}——只作过渡，引出第二段，不含新信息。
\end{tpl}

\rubric{第一段对应评分档位里的「\textbf{格式和语气恰当}」与「\textbf{回应任务}（点题）」两项。称呼、首句语气、目的句对上了，格式分与语气分就拿到了；第一段本身不承担内容要点的展开。}

\tip{称呼与署名一律按题目要求写。题目说「Do not use your own name, use “Li Ming” instead.」时，称呼用 Dear Sir or Madam / Dear Jack / Dear Professor Smith，署名只写 \kw{Li Ming}。题干给了具体人名（如 Jack、David、John）就必须用题目给的名字，不要自己编。}


% =====================================================================
\subsection{第 1 句：问候与铺垫}
% =====================================================================

\begin{note}
下面 25 条按使用场景分成四组。带「首推」标记的 7 条覆盖英语二历年绝大多数真题场景，优先背这 7 条；其余条目作同义替换，进考场前扫一眼即可。所有句子都在句末按需加逗号或句号，接到第 2 句的「目的句」上。
\end{note}

\subsubsection{（1）通用问候（任何书信都能用）}

\begin{eg}
1.\ I hope this message finds you well.\\
希望你一切安好。\kw{（首推·万能）}

2.\ I hope everything is going smoothly with your study and life.\\
希望你的学习和生活一切顺利。

3.\ I hope you are doing well these days.\\
希望你最近一切顺利。\kw{（首推）}

4.\ I hope everything is going well on your side.\\
希望你那边一切顺利。

5.\ I hope you are enjoying good health and a pleasant life.\\
希望你身体健康、生活愉快。

6.\ I hope all is well with you.\\
希望你一切都好。

7.\ It is a great pleasure to write to you.\\
能给你写这封信，我感到十分高兴。\kw{（首推·正式）}

8.\ It is a great pleasure for me to write this letter to you.\\
能给你写这封信，我感到十分高兴。（原书完整版，可用于正式书信）

9.\ I feel honored to be able to write to you today.\\
今天能够给你写信，我感到很荣幸。\kw{（首推·对师长/机构）}

10.\ I am truly glad to get in touch with you.\\
能够和你取得联系，我由衷感到高兴。
\end{eg}

\subsubsection{（2）收到来信 / 再次联系}

\begin{eg}
1.\ I am so delighted to receive your letter.\\
非常高兴收到你的来信。

2.\ I am delighted to have received your letter.\\
收到你的来信，我非常高兴。（书面化改写，语气更稳）

3.\ It is really nice to hear from you again.\\
很高兴再次收到你的消息。\kw{（首推·回信类）}

4.\ Words cannot express how delighted I was to learn that \fillin{事由}.\\
得知\fillin{事由}，我高兴得难以言表。\kw{（首推·情绪铺垫）}

5.\ Words failed to express my delight when I received your email.\\
收到你的邮件时，我高兴得难以言表。（御风范文句，注意时态用 failed 更自然）

6.\ Words cannot express my excitement when I decide to tell you about \fillin{事由}.\\
一想到要告诉你\fillin{事由}，我就激动得难以言表。（御风 2025 英二范文起手式）
\end{eg}

\tip{「Words fail to express my ...」这一结构的用法见本节末【本次改进】：原书写成 \textit{words fail to express my excitement to inform you about ...} 是不地道的搭配，已按上表第 4、6 条改写。}

\subsubsection{（3）通知、告示类开场（Notice / Announcement）}

\begin{note}
通知、告示的第一句不问候个人，而是对全体读者致谢或致意。主语用 we（学生会、组委会、学院），或改用被动式 this notice is issued to 更正式。中文里「本通知」对应 Notice 居中大写标题，时间写在第一行最右。
\end{note}

\begin{eg}
1.\ We greatly appreciate all students for sparing your precious time to pay attention to this notice.\\
我们十分感谢全体同学抽出宝贵时间关注本通知。

2.\ Our sincere gratitude goes to you for your willingness to read and comply with the arrangements in this notice.\\
诚挚感谢各位愿意阅读并遵照本通知中的相关安排。

3.\ We would like to extend our thanks to every student for sparing your time amid busy schedules to pay attention to this notice.\\
谨向每位同学致谢，感谢大家百忙之中关注本通知。

4.\ We are pleased to inform you of \fillin{通知事项}.\\
我们很高兴通知大家\fillin{通知事项}。（更亲切的语气）

5.\ This notice is issued to provide you with some essential information regarding \fillin{通知要点}.\\
本通知旨在就\fillin{通知要点}向大家提供必要信息。（更正式的语气）
\end{eg}

\subsubsection{（4）会议纪要类开场（Meeting Minutes）}

\begin{note}
纪要（Meeting Minutes）第一句用「已记录／已整理」的完成时，直接引出下文的分条要点。纪要的正题通常按 I. Meeting Details（Date / Time / Location / Attendees）、II. Meeting Objective、III. Meeting Agenda 的固定框架排布，第一段的功能句只负责把它带出来。
\end{note}

\begin{eg}
1.\ I have documented the content of the previous seminar, and the key points of the meeting minutes are listed as follows.\\
我已记录上次课程研讨会的内容，本次会议纪要的核心要点整理如下。\kw{（首推）}

2.\ The content of the last graduation project team meeting has been recorded, and the detailed meeting minutes are presented below.\\
我已记录上次毕设项目组会议的内容，具体会议纪要内容如下所示。

3.\ I have compiled the content of the previous academic exchange meeting, and the structured meeting minutes are provided as follows.\\
我已整理上次学术交流会的内容，条理化的会议纪要内容详见下文。

4.\ I have documented the discussion at the \fillin{会议名称}, and the minutes are presented below.\\
我已记录\fillin{会议名称}的讨论内容，纪要如下。
\end{eg}

\rubric{第 1 句只需「让人知道你在跟谁说话、以什么身份说话」，不需要任何题目信息。它决定阅卷人对「格式与语气」的第一印象，因此宁可短、宁可平实，不要为了华丽而写长。}

\zone{不要在第 1 句里塞题目要点。要点全部放在第二段；第 1 句一旦提前把要点说完，第二段就只能重复，阅卷人会直接判定「内容单薄」。}


% =====================================================================
\subsection{第 2 句：点明写作目的（按信件类型）}
% =====================================================================

\begin{note}
下表把原书「第一段第二句」27 条全部收入，并按信件类型重新归位（原书是平铺编号，类型混杂，同一类型散落在 1、14、15、18、19、23、24 等多处）。每条都已还原 OCR 错误（\textit{thc / mattcr / sinccre / duc coursc / 題目相关词} 等）并配中文；表中方括号处填写题干事由。
\end{note}

\begin{longtable}{L{2.0cm}L{9.2cm}L{4.6cm}}
\toprule
\textbf{信件类型} & \textbf{功能句（第 2 句）} & \textbf{中文} \\
\midrule
\endfirsthead
\toprule
\textbf{信件类型} & \textbf{功能句（第 2 句）} & \textbf{中文} \\
\midrule
\endhead
\bottomrule
\endlastfoot

% ---- 混合类 ----
混合类 & I am writing to you regarding the matter of \fillin{题目相关词}. & 我写信是为了\fillin{题目相关词}这一事项。 \\
混合类 & I am writing on behalf of \fillin{代表对象} to discuss \fillin{题目相关词}. & 我代表\fillin{代表对象}给你写信，讨论\fillin{题目相关词}。 \\
混合类 & The purpose of this letter is to discuss \fillin{题目相关词} with you in detail. & 这封信的目的是想与你详细讨论\fillin{题目相关词}。 \\
混合类 & I am writing to inform you of \fillin{事由}. & 我写信是想告知你\fillin{事由}。 \\
\midrule

% ---- 感谢 ----
感谢信 & I am writing to express my sincere gratitude for \fillin{帮助内容}. & 我写信是为了对\fillin{帮助内容}表达诚挚的感谢。 \\
感谢信 & This letter is meant to convey my heartfelt thanks for your support with \fillin{题目相关词}. & 这封信旨在对你在\fillin{题目相关词}方面给予的支持表达由衷感谢。 \\
感谢信 & The purpose of this letter is to extend my sincere gratitude for your help with \fillin{题目相关词}. & 本信的目的是为你在\fillin{题目相关词}上给予的帮助致以诚挚谢意。 \\
\midrule

% ---- 邀请 ----
邀请信 & I am writing to invite you to take part in \fillin{题目相关词}. & 我写信是想邀请你参加\fillin{题目相关词}。 \\
邀请信 & The purpose of this letter is to invite you to take an active part in \fillin{题目相关词}. & 本信的目的是邀请你积极参加\fillin{题目相关词}。 \\
邀请信 & Through this letter, I sincerely invite you to join us in \fillin{题目相关词}. & 借此来信，我诚挚邀请你参与我们组织的\fillin{题目相关词}。 \\
\midrule

% ---- 建议 ----
建议信 & With this letter, I would like to share my views on \fillin{题目相关词}. & 借此来信，我想和你分享我对\fillin{题目相关词}的看法。 \\
建议信 & The primary aim of this letter is to offer some suggestions on \fillin{题目相关词}. & 本信的首要目的是就\fillin{题目相关词}提出若干建议。 \\
建议信 & I am writing to share with you some suggestions concerning \fillin{题目相关词}. & 我写信是想就\fillin{题目相关词}和你分享一些建议。 \\
\midrule

% ---- 投诉 ----
投诉信 & This letter serves as a formal complaint about the recent problems with \fillin{题目相关词}. & 这封信是一份正式投诉，反映近期与\fillin{题目相关词}有关的问题。 \\
投诉信 & Through this letter, I would like to call your attention to \fillin{题目相关词}. & 借此来信，我想引起你对\fillin{题目相关词}的重视。 \\
\midrule

% ---- 道歉 ----
道歉信 & This letter is intended to offer a sincere apology for the problems caused by \fillin{题目相关词}. & 这封信旨在就\fillin{题目相关词}所带来的问题向你致以诚挚的歉意。 \\
道歉信 & In this letter, I would like to explain the reasons behind \fillin{题目相关词} and ask for your understanding. & 在信中，我想说明\fillin{题目相关词}背后的原因，并请求你的理解。 \\
道歉信 & I am writing to express my sincere apology for \fillin{过失行为}. & 我写信是想就\fillin{过失行为}表达我诚挚的歉意。 \\
\midrule

% ---- 申请 / 求职 ----
申请、求职信 & I am writing to apply for the position of \fillin{职位} advertised in \fillin{信息来源}. & 我写信是想应聘在\fillin{信息来源}上发布的\fillin{职位}一职。 \\
申请、求职信 & I am writing to express my strong interest in \fillin{项目或岗位} and to apply for it. & 我写信是想表达我对\fillin{项目或岗位}的浓厚兴趣，并申请加入。 \\
\midrule

% ---- 推荐 ----
推荐信 & The purpose of this letter is to recommend \fillin{推荐对象} as a valuable opportunity for you. & 本信旨在向你推荐\fillin{推荐对象}，这对你而言是一次宝贵的机会。 \\
推荐信 & I am writing to recommend \fillin{推荐对象} to you, which I believe is fully qualified for \fillin{目标用途}. & 我写信是想向你推荐\fillin{推荐对象}，我认为它完全适合\fillin{目标用途}。 \\
\midrule

% ---- 咨询 / 求助 ----
咨询、求助信 & The central purpose of this letter is to seek your advice on \fillin{题目相关词}. & 此信的核心目的是就\fillin{题目相关词}向你征求意见。 \\
咨询、求助信 & I am writing to seek your guidance on \fillin{题目相关词}. & 我写信是想就\fillin{题目相关词}向你寻求指导。 \\
\midrule

% ---- 祝贺 ----
祝贺信 & Please allow me to extend my warmest congratulations to you. & 请允许我向你致以最诚挚的祝贺。 \\
祝贺信 & This letter is intended to extend my warm congratulations on your achievement in \fillin{题目相关词}. & 这封信旨在对你在\fillin{题目相关词}方面取得的成就表示热烈祝贺。 \\
祝贺信 & The purpose of this letter is to congratulate you sincerely on your remarkable progress in \fillin{题目相关词}. & 本信旨在对你在\fillin{题目相关词}上取得的显著进步致以诚挚祝贺。 \\
\midrule

% ---- 介绍 ----
介绍信 & In this letter, I would like to introduce \fillin{介绍对象} to you as an excellent option. & 在这封信中，我想向你介绍\fillin{介绍对象}，把它作为一个极佳的选择推荐给你。 \\
介绍信 & I am writing to give you a brief introduction to \fillin{介绍对象}. & 我写信是想向你简要介绍\fillin{介绍对象}。 \\
\midrule

% ---- 通知 / 告示 ----
通知、告示 & In order to further promote \fillin{题目相关词}, our school has decided to issue the following notice. & 为了进一步推进\fillin{题目相关词}，学校决定发布如下通知。 \\
通知、告示 & To ensure the smooth implementation of \fillin{题目相关词}, the following arrangements are hereby announced. & 为确保\fillin{题目相关词}顺利开展，现将相关安排通知如下。 \\
通知、告示 & To raise students' awareness of \fillin{题目相关词}, the school will carry out the following activities. & 为提升学生对\fillin{题目相关词}的认识，学校将开展以下活动。 \\
通知、告示 & This announcement aims to provide detailed information about \fillin{题目相关词}. & 本公告旨在提供关于\fillin{题目相关词}的详细信息。 \\
通知、告示 & The purpose of this notice is to remind everyone of the upcoming activities related to \fillin{题目相关词}. & 本通知旨在提醒大家即将开展的、与\fillin{题目相关词}相关的活动。 \\
通知、告示 & This notice is intended to clarify the requirements for participating in \fillin{题目相关词}. & 本通知旨在说明参与\fillin{题目相关词}的具体要求。 \\
\midrule

% ---- 纪要 ----
纪要 & I am writing to submit the minutes of \fillin{会议名称} held on \fillin{日期}. & 现提交\fillin{日期}召开的\fillin{会议名称}会议纪要。 \\
纪要 & The following minutes record the key points of \fillin{会议名称}. & 以下纪要记录了\fillin{会议名称}的核心要点。 \\

\end{longtable}

\begin{note}
\textbf{通知、告示的第一段标准写法}（御风，可直接套用）：\par
To ensure that \fillin{通知对象} are fully informed of \fillin{通知事项}, this notice is issued to provide you with some essential information regarding its \fillin{通知要点}. And the detailed arrangements are as follows.\par
为确保\fillin{通知对象}充分了解\fillin{通知事项}，本通知旨在就\fillin{通知要点}向大家提供必要信息，具体安排如下。\par
语气微调：\kw{this notice is issued to}（正式）／\kw{we are pleased to}（亲切）／\kw{I am writing to}（个人身份发通知，如班长）。
\end{note}


% =====================================================================
\begin{note}
下面示范里的收信人姓名（Bob、Jack、David、Professor Smith 等）\textbf{只是示例}。考场上一切以题干为准：题目给了姓名就用题目的姓名，题目只说 “your friend” 就用一个常见的英文名，题目要求 “Do not use your own name. Use 'Li Ming' instead.” 时，\textbf{只有署名受此限制，称呼仍按题干给的人写}。
\end{note}

\subsection{真题示范（第一段整段）}
% =====================================================================

\begin{note}
以下 8 段全部按「第 1 句问候 + 第 2 句目的」写成，称呼按真题要求统一处理，署名均为 Li Ming。英文原句来自 mmg 的真题示例与御风范文，已逐句纠正 OCR 错误并统一文体。
\end{note}

\begin{tpl}{2010 英语二 · 感谢 + 邀请（混合类）}
Dear Bob,\\
I hope this message finds you well. I am writing to express my sincere thanks for your warm reception during my stay in the United States, and to invite you to visit China in due course.\\[2pt]
希望你一切安好。我写信是想表达两件事：一是对你在我赴美期间热情接待的诚挚感谢，二是诚挚邀请你在方便的时候来中国访问。
\end{tpl}

\begin{tpl}{2011 英语二 · 祝贺 + 建议（混合类）}
Dear Li Ming,\\
It is a great pleasure to write to you. Please allow me to extend my warmest congratulations on your admission to the university, and I would also like to offer you some suggestions on how to get prepared for university life.\\[2pt]
能给你写信，我感到十分高兴。请允许我对你考入大学致以最诚挚的祝贺，同时我也想就如何为大学生活做准备给你一些建议。
\end{tpl}

\begin{tpl}{2012 英语二 · 投诉}
Dear Sir or Madam,\\
I feel honored to be able to write to you today. This letter serves as a formal complaint about the recent problems with the electronic dictionary I bought from your online store.\\[2pt]
今天能够给您写信，我感到很荣幸。这封信是一份正式投诉，反映我最近在贵网店购买的电子词典所出现的问题。
\end{tpl}

\begin{tpl}{2014 英语二 · 咨询}
Dear Sir or Madam,\\
I hope everything is going smoothly with your study and life. I am writing to tell you about my daily habits and to seek your advice on how I can better adapt to our shared apartment life abroad.\\[2pt]
希望你的学习和生活一切顺利。此信的核心目的是就我如何更好地适应今后在国外的合租生活向你征求意见。
\end{tpl}

\begin{tpl}{2018 英语二 · 道歉}
Dear Professor Smith,\\
I am truly glad to get in touch with you. This letter is intended to offer a sincere apology for the problems caused by my having to cancel my visit to you at the last minute.\\[2pt]
能够和你取得联系，我由衷感到高兴。这封信旨在就我在最后时刻不得不取消拜访而给你带来的种种麻烦，致以诚挚的歉意。
\end{tpl}

\begin{tpl}{2021 英语二 · 邀请}
Dear Jack,\\
I hope you are doing well these days. I am writing to invite you to take part in our upcoming online meeting about study and campus life at our university.\\[2pt]
希望你最近一切顺利。我写信是想邀请你参加我们即将举行的、以本校学习和校园生活为主题的线上会议。
\end{tpl}

\begin{tpl}{2023 英语二 · 建议}
Dear David,\\
It is really nice to hear from you again. With this letter, I would like to share my views on whether you should choose the art exhibition or the robot show this Sunday.\\[2pt]
很高兴再次收到你的消息。借此来信，我想和你分享我对本周日该选择参加艺术展还是机器人展的看法。
\end{tpl}

\begin{tpl}{2024 英语二 · 汇报计划（讨论类）}
Dear John,\\
It is a great pleasure to write to you. The purpose of this letter is to discuss our survey plan on the protection of old houses in the ancient town with you in detail.\\[2pt]
能给你写信，我感到十分高兴。这封信的目的是想和你详细讨论我们关于古镇老宅保护调研计划的具体安排。
\end{tpl}

\begin{note}
\textbf{两处需要更正的原书标注}：

（1）原书把「The purpose of this letter is to congratulate you sincerely on your remarkable progress in English-Chinese translation.」标注为\textbf{2016 英语二}示例。但 2016 英语二的真题指令是：朋友 Jack 来信祝贺你获奖并询问翻译建议，你回信须 1）感谢他 2）给出建议——属于\textbf{感谢 + 建议}类，不是祝贺类。该句更贴合 2011 英语二（祝贺考入大学）或通用祝贺语境，本册已按类归位。

（2）原书把 2022 英语二的示例误抄成了 2020 英语一的句子（「To raise students' awareness of ...」条下的示例沿用了上一页的校园歌唱比赛句）。2022 英语二真题为校园美食节（International Food Festival），正确示例见第二段章节。

\tip{第一段占比：100 词里第一段 25～30 词，第二段 55～65 词，第三段 15～20 词。这是英语二应用文「内容为主、格式为辅」的正常配比。}
\end{note}

\impr{对第一段所做的改动与理由，附改进前后对照：}

\begin{longtable}{L{6.4cm}L{4.6cm}L{4.8cm}}
\toprule
\textbf{原书 / OCR 原句} & \textbf{问题} & \textbf{改进后} \\
\midrule
\endfirsthead
\toprule
\textbf{原书 / OCR 原句} & \textbf{问题} & \textbf{改进后} \\
\midrule
\endhead
\bottomrule
\endlastfoot

Words fail to express my excitement to inform you about our approaching innovation contest. & express 后面直接跟 to inform 目的状语，语义变成「表达我的兴奋是为了通知你」，搭配不成立；excitement 也不宜作 express 的宾语来引出消息。 & Words cannot express how excited I am to tell you about our upcoming innovation contest. \\
\midrule
Words failed to express my delight when I received your email. & 原书是 words fail（现在时），事件已发生，时态不一致。 & Words failed to express my delight when I received your email.（御风范文原句，已保留并标注） \\
\midrule
It is a great pleasure for me to write this letter to you. & 语法正确，但作为第 1 句与第 2 句连用会吃掉太多字数。 & It is a great pleasure to write to you. \\
\midrule
I am writing to you regarding the matter of my sincere thanks for your warm reception during my stay in the United States and my invitation for you to visit China in due course. & regarding the matter of 后接两个抽象名词并列，句子过长；且 my invitation 与 my thanks 并列时谓语语义重叠，读者要读两遍才懂是在说「感谢 + 邀请」两件事。 & I am writing to express my sincere thanks for your warm reception during my stay in the United States, and to invite you to visit China in due course. \\
\midrule
I hope you are enjoyinggood health and a pleasant life.（OCR 粘连） & 识别错误；且表述略中式。 & I hope you are enjoying good health and a pleasant life.（可用）／更简：I hope all is well with you. \\
\midrule
I have documented the content of the previous seminar, and the key points of the meeting minutes are listed as tollows.（OCR） & tollows $\rightarrow$ as follows；且 meeting minutes 与 key points 语义重复。 & I have documented the content of the previous seminar, and the key points of the meeting minutes are listed as follows.（保留原句，仅纠错） \\
\midrule
1am writing to you regarding thc mattcr of ...（OCR） & I/1 混淆、thc/mattcr 等系统性识别错误。 & I am writing to you regarding the matter of ... \\
\midrule
Iam writing on behalf of our class to discuss ...（OCR） & Iam 粘连。 & I am writing on behalf of our class to discuss ... \\
\midrule
To raise students' awareness of 题目相关词, thc school will carry out the following activities.（OCR） & thc $\rightarrow$ the。 & To raise students' awareness of \fillin{题目相关词}, the school will carry out the following activities. \\

\end{longtable}

\begin{note}
\textbf{未做合并、全部保留的条目}：原书「第一段第一句」16 条与「第一段第二句」27 条已全部收入本节的四组问候句、目的句长表，无一条被删除。只做了三件事——按类型重排、纠正 OCR 错误、为每条补写中文（原书中文对应的是「题目相关词」式占位说明，本册全部改写成可读译文）。
\end{note}

% =====================================================================
% 75_letter_p2.tex
% 小作文（Part A 应用文）第二段功能句库
% 来源：《mmg 作文模板》三、通用小作文——第二段第一句（16 条）、
%       建议类 / 邀请类 / 道歉类 / 祝贺类 / 介绍类 / 通知类
%       各含「第二句 / 第三句 / 第四句功能句挑选」；
%       《英语作文模板》（御风）事件六要素表达、各类信件「结果/效果」句、
%       通知中间段写法与真题范文第二段。
% =====================================================================

\section{小作文第二段功能句库}

\begin{note}
第二段是\textbf{得分主体}，写 3～5 句，占全文 55～65 词。它与第一段的区别在于：第一段是套话，第二段\textbf{必须由你按题目填充}——把题干的写作要点逐条回应，再挂上功能句的骨架。
\end{note}

\begin{tpl}{第二段公式}
\textbf{第 1 句（统领句）}：\kw{To be more specific, the main points can be arranged as follows.}——把话说下去，不承载信息。\par
\textbf{第 2 句}：第一层实质内容（第一类建议 / 邀请 + 收获 / 承认过失 + 影响 / 夸成就 + 品质 / 给出定义 / 交代时间地点）。\par
\textbf{第 3 句}：第二层实质内容（追加建议 / 参加的价值 / 解释原因 / 归因 / 展开特点 / 内容与环节）。\par
\textbf{第 4 句}：收束（\kw{Last but not least ...} + 好处 / 补救措施 / 未来展望 / 注意事项）。\par
\textbf{第 5 句（可选）}：如果题干有三个要点，用 \kw{Taken together / Handled in this way} 再补一句。
\end{tpl}

\rubric{第二段对应评分档位里的「\textbf{内容要点完整}」（英语二题目通常直接给出 1、2 两条写作要点，漏一条即扣档）与「\textbf{格式和语气恰当}」。阅卷时先看要点有没有逐条落地，再看语言是否准确、得体。}

\zone{第二段的\textbf{头号雷区}是把它写成大作文式的空泛议论。大作文第一段讲完现象就可以往上位概念套，小作文不行——英语二会直接给出写作要点，你\textbf{必须逐点回应，不重、不漏、不跑题}。凡是题目提到的要点，一定要全部回应。}

\zone{第二号雷区是\textbf{分点罗列}。书信是偏生活化的文风，不要写成 \kw{1. ... 2. ... 3. ...} 的公文条目，就用 First and foremost / Moreover / Last but not least 这类连接词把要点串成一段连贯的话。只有通知（Notice）允许写得规整并列。}

\tip{「第二段第二句」不一定要写在第二句。原书特别提示：把 \kw{First}、\kw{What's more} 这类前缀换一换，功能句可以放在第 2、3、4 句的任意位置。真正固定的是\textbf{三类内容各出现一次}：核心内容 → 追加内容 → 收束。}


% =====================================================================
\subsection{第二段通用统领句（第 1 句）}
% =====================================================================

\begin{note}
下表完整收入原书「第二段第一句功能句挑选」16 条，一条未删。其中 \textbf{8 条标了「首推」}，是按文体首选；其余 8 条同样可用，只是与前面语义较为接近，适合用来换词、避免多篇作文重复。选句原则：\textbf{按第二段要写的内容选}——写建议选 measures / suggestions 类，写活动选 arrangements / planned event 类，写介绍选 specific details 类，写通知选 notice / schedule 类。
\end{note}

\begin{longtable}{L{0.9cm}L{10.4cm}L{4.5cm}}
\toprule
\textbf{序} & \textbf{统领句（英文）} & \textbf{中文} \\
\midrule
\endfirsthead
\toprule
\textbf{序} & \textbf{统领句（英文）} & \textbf{中文} \\
\midrule
\endhead
\bottomrule
\endlastfoot

1 & To be more specific, the main points can be arranged as follows. \kw{（首推·万能）} & 更具体地说，主要内容可以概括为以下几点。 \\
2 & In terms of concrete arrangements, the details are as follows. & 就具体安排而言，细节如下。 \\
3 & With regard to the practical steps, the key information is listed below. & 关于实际步骤，关键信息如下所示。 \\
4 & As for the specific arrangements, they can be outlined in the following way. \kw{（首推·活动）} & 至于具体安排，可以概括为以下几点。 \\
5 & When it comes to the concrete details, they may be presented as follows. \kw{（首推·介绍）} & 说到具体细节，可以作如下说明。 \\
6 & As for the practical measures, my suggestions are listed as follows. \kw{（首推·建议）} & 至于具体措施，我的建议如下。 \\
7 & Concerning the next steps, several suggestions can be made as follows. & 关于下一步的做法，可以提出如下建议。 \\
8 & In light of the above, the following suggestions may be of some help. & 鉴于以上情况，以下几点建议或许有所帮助。 \\
9 & To address this issue more effectively, the following actions are recommended. \kw{（首推·建议）} & 为更有效地解决这一问题，建议采取以下行动。 \\
10 & As for the underlying reasons, they can be explained as follows. \kw{（首推·原因）} & 至于背后的原因，可以作如下解释。 \\
11 & To justify my opinion, several points need to be mentioned as follows. & 为说明我的观点，有必要提出以下几点。 \\
12 & With respect to the causes of this phenomenon, they can be summarized as follows. & 至于这一现象的成因，可以归纳为以下几点。 \\
13 & As for the planned event, the main arrangements are as follows. \kw{（首推·活动）} & 至于本次活动的具体安排，主要如下。 \\
14 & With respect to the upcoming activity, the relevant information is given below. & 关于即将举行的活动，相关信息如下。 \\
15 & Regarding the details of this notice, please pay attention to the following points. \kw{（首推·通知）} & 关于本通知的具体内容，请留意以下几点。 \\
16 & As far as the schedule is concerned, the key items are listed as follows. & 就时间安排而言，主要事项如下。 \\

\end{longtable}


% =====================================================================
\subsection{建议类}
% =====================================================================

\begin{tpl}{建议类 · 第二段骨架}
\textbf{第 2 句}：提出第一条也是最核心的建议，并说明它带来的好处。\par
\textbf{第 3 句}：追加第二条建议，或强调「你必须做某事」的必要性。\par
\textbf{第 4 句}：最后一条建议 + \kw{Only in this way} 收尾，落到好处上。\par
（题目要两点就写 2、4 两句；要三点就三句都写。建议类最常见的落点是「怎么做 + 为什么有用」。）
\end{tpl}

\subsubsection{第 2 句功能句（7 条）}

\begin{eg}
1.\ First and foremost, putting \fillin{题目相关词} into consistent practice is strongly recommended, as it can help you have a clearer insight into \fillin{题目相关短语或词}.\\
首先，强烈建议你把\fillin{题目相关词}持续付诸实践，因为这有助于你更清楚地认识\fillin{题目相关短语或词}。

2.\ First and foremost, cultivating \fillin{题目相关词} in your daily routine is strongly recommended, as it can help you have a clearer insight into \fillin{题目相关短语或词}.\\
首先，强烈建议你在日常作息中不断培养\fillin{题目相关词}，因为这能帮助你更清楚地认识\fillin{题目相关短语或词}。

3.\ At this stage, it would be wise for you to \fillin{建议行动，动词短语} in terms of \fillin{题目相关词}, as this will help you better understand \fillin{题目相关短语或词}.\\
在当前阶段，建议你在\fillin{题目相关词}方面\fillin{建议行动，动词短语}，因为这将帮助你更好地理解\fillin{题目相关短语或词}。

4.\ In particular, you are advised to \fillin{建议行动，动词短语} with a focus on \fillin{题目相关词} so that you can gradually improve \fillin{题目相关短语或词}.\\
具体而言，建议你围绕\fillin{题目相关词}\fillin{建议行动，动词短语}，以便逐步提升\fillin{题目相关短语或词}。

5.\ A sensible step would be to \fillin{建议行动，动词短语} with a clear focus on \fillin{题目相关词}, which is essential for \fillin{题目相关短语或词}.\\
明智之举是\fillin{建议行动，动词短语}，并把重点放在\fillin{题目相关词}上，这对\fillin{题目相关短语或词}至关重要。

6.\ To move from intention to action, you might \fillin{建议行动，动词短语} with a focus on \fillin{题目相关词}, thereby strengthening \fillin{题目相关短语或词}.\\
为把想法化为行动，你不妨围绕\fillin{题目相关词}\fillin{建议行动，动词短语}，从而强化\fillin{题目相关短语或词}。

7.\ To make your efforts more effective, it is preferable to \fillin{建议行动，动词短语} with a focus on \fillin{题目相关词}, which will sharpen \fillin{题目相关短语或词}.\\
为使你的努力更为有效，最好能结合\fillin{题目相关词}\fillin{建议行动，动词短语}，这将有助于提升\fillin{题目相关短语或词}。
\end{eg}

\subsubsection{第 3 句功能句（7 条）}

\begin{eg}
1.\ Furthermore, there is no denying that it is of great necessity for you to \fillin{要做的事}.\\
此外，不可否认的是，你有必要\fillin{要做的事}。

2.\ Moreover, it is widely recognized that it is of vital importance for you to \fillin{要做的事}.\\
此外，普遍认为你\fillin{要做的事}具有至关重要的意义。

3.\ In addition, it is strongly recommended that you should \fillin{要做的事} so as to improve the situation.\\
另外，强烈建议你\fillin{要做的事}，以便改善当前状况。

4.\ What is more, it is highly advisable for you to \fillin{要做的事} in order to achieve better results.\\
更重要的是，为了取得更好的效果，非常建议你\fillin{要做的事}。

5.\ Besides, it is of great significance for you to \fillin{要做的事}, which will benefit you in the long run.\\
此外，你\fillin{要做的事}具有重要意义，这将在长远来看对你大有裨益。

6.\ Furthermore, you are expected to \fillin{要做的事}, which can help you handle \fillin{题目相关词} more effectively.\\
此外，希望你能够\fillin{要做的事}，这将帮助你更有效地应对\fillin{题目相关词}。

7.\ What is more, it is essential for you to \fillin{要做的事}, because this is closely related to \fillin{题目相关词}.\\
更重要的是，你必须\fillin{要做的事}，因为这与\fillin{题目相关词}密切相关。
\end{eg}

\subsubsection{第 4 句功能句（7 条）}

\begin{eg}
1.\ Last but not least, it is advisable for you to \fillin{要做的事}. Only in this way can you truly \fillin{好处是什么}.\\
最后，你最好\fillin{要做的事}。唯有如此，你才能真正\fillin{好处是什么}。

2.\ Last but not least, you are strongly encouraged to \fillin{要做的事}. Only in this way will you be able to \fillin{好处是什么}.\\
最后，强烈建议你\fillin{要做的事}。只有这样，你才有可能\fillin{好处是什么}。

3.\ Last but not least, it would be a wise choice for you to \fillin{要做的事}. Only in this way can you better \fillin{好处是什么}.\\
最后，你\fillin{要做的事}不失为明智之举。唯有如此，你才能更好地\fillin{好处是什么}。

4.\ Last but not least, it is of great necessity for you to \fillin{要做的事}. Only in this way can you gradually \fillin{好处是什么}.\\
最后，你\fillin{要做的事}是十分必要的。只有这样，你才能逐步\fillin{好处是什么}。

5.\ Last but not least, you should make a conscious effort to \fillin{要做的事}. Only in this way can you effectively \fillin{好处是什么}.\\
最后，你应当有意识地\fillin{要做的事}。唯有如此，你才能有效地\fillin{好处是什么}。

6.\ Last but not least, \fillin{要做的事} will enable you to \fillin{好处是什么} in terms of \fillin{题目相关词}.\\
最后，\fillin{要做的事}将使你在\fillin{题目相关词}方面得以\fillin{好处是什么}。

7.\ Only by \fillin{要做的事} can you truly \fillin{好处是什么} when dealing with \fillin{题目相关词}.\\
唯有\fillin{要做的事}，你在面对\fillin{题目相关词}时才能真正\fillin{好处是什么}。
\end{eg}


% =====================================================================
\subsection{邀请类}
% =====================================================================

\begin{tpl}{邀请类 · 第二段骨架}
\textbf{第 2 句}：正式发出邀请，交代活动（时间 / 地点）+ 参加能收获什么。\par
\textbf{第 3 句}：说明参加的价值——平台、机会、能力或品质的提升。\par
\textbf{第 4 句}：呼吁出席，强调「不应错过 / 有了你活动更有意义」。\par
（邀请类的核心是「给对方一个来的理由」：收获、平台、不可错过。）
\end{tpl}

\subsubsection{第 2 句功能句（7 条）}

\begin{eg}
1.\ In particular, I would like to invite you to \fillin{受邀行动，如 take part in} \fillin{题目相关词}, which will offer you \fillin{参加收获}.\\
具体而言，我诚挚邀请你\fillin{受邀行动，如 take part in}\fillin{题目相关词}，这将为你带来\fillin{参加收获}。

2.\ More specifically, your participation in \fillin{题目相关词} at \fillin{时间地点信息} would be greatly appreciated, as it will enable you to \fillin{参加收获}.\\
更具体地说，如果你能于\fillin{时间地点信息}参与\fillin{题目相关词}，我们将不胜感激，因为这将使你得以\fillin{参加收获}。

3.\ As for the concrete arrangement, you are warmly invited to \fillin{受邀行动，动词短语} in connection with \fillin{题目相关词} so that you can \fillin{参加收获}.\\
至于具体安排，特此诚挚邀请你围绕\fillin{题目相关词}\fillin{受邀行动，动词短语}，以便你能够\fillin{参加收获}。

4.\ At this point, it would be a great honor if you could \fillin{受邀行动，动词短语} for \fillin{题目相关词}, which serves as an ideal chance to \fillin{参加收获}.\\
在此时刻，若你能为\fillin{题目相关词}\fillin{受邀行动，动词短语}，将是我们的莫大荣幸，这也是\fillin{参加收获}的绝佳机会。

5.\ In relation to \fillin{题目相关词}, we sincerely invite you to \fillin{受邀行动，动词短语}, allowing you to \fillin{参加收获}.\\
就\fillin{题目相关词}而言，我们诚挚邀请你\fillin{受邀行动，动词短语}，从而使你得以\fillin{参加收获}。

6.\ With regard to \fillin{题目相关词}, it is my pleasure to invite you to \fillin{受邀行动，动词短语}, which can greatly broaden \fillin{参加相关短语}.\\
关于\fillin{题目相关词}，我很高兴邀请你\fillin{受邀行动，动词短语}，这将极大拓展\fillin{参加相关短语}。

7.\ To make full use of \fillin{题目相关词}, you are cordially invited to \fillin{受邀行动，动词短语}, through which you may \fillin{参加收获}.\\
为更好地发挥\fillin{题目相关词}的价值，特此诚挚邀请你\fillin{受邀行动，动词短语}，借此你可以\fillin{参加收获}。
\end{eg}

\subsubsection{第 3 句功能句（7 条）}

\begin{eg}
1.\ Furthermore, this will be a valuable opportunity for you to \fillin{要做的事}, especially in the context of \fillin{题目相关词}.\\
此外，借此你将获得一次\fillin{要做的事}的宝贵机会，尤其是在\fillin{题目相关词}这一背景下。

2.\ Moreover, it is widely believed that \fillin{要做的事} at \fillin{题目相关词} can help you \fillin{邀请带来的好处}.\\
再者，普遍认为在\fillin{题目相关词}中\fillin{要做的事}有助于你\fillin{邀请带来的好处}。

3.\ In addition, by \fillin{要做的事} during \fillin{题目相关词}, you will be able to \fillin{邀请带来的好处}, which is of great benefit to you.\\
另外，在\fillin{题目相关词}期间\fillin{要做的事}，将使你能够\fillin{邀请带来的好处}，这对你大有裨益。

4.\ What is more, taking this chance to \fillin{要做的事} in \fillin{题目相关词} will significantly enhance \fillin{邀请相关能力或品质}.\\
更重要的是，借此机会在\fillin{题目相关词}中\fillin{要做的事}，将显著提升\fillin{邀请相关能力或品质}。

5.\ Besides, \fillin{要做的事} in the context of \fillin{题目相关词} will provide you with first-hand experience to \fillin{邀请带来的好处}.\\
此外，在\fillin{题目相关词}这一情境下\fillin{要做的事}，还能为你提供\fillin{邀请带来的好处}的一手体验。

6.\ Furthermore, engaging in \fillin{题目相关词} and \fillin{要做的事} will create a rare platform for you to \fillin{邀请带来的好处}.\\
此外，参与\fillin{题目相关词}并\fillin{要做的事}，将为你搭建一个难得的平台，使你得以上\fillin{邀请带来的好处}。

7.\ Moreover, through \fillin{题目相关词}, the chance to \fillin{要做的事} will help you build a stronger sense of \fillin{邀请相关短语}.\\
再者，通过\fillin{题目相关词}，\fillin{要做的事}的机会将帮助你建立更强的\fillin{邀请相关短语}。
\end{eg}

\subsubsection{第 4 句功能句（7 条）}

\begin{eg}
1.\ Last but not least, your presence will make \fillin{题目相关词} more meaningful and enable you to \fillin{好处是什么}.\\
最后，你的到来将使\fillin{题目相关词}更加有意义，并使你得以\fillin{好处是什么}。

2.\ Last but not least, accepting this invitation to \fillin{题目相关词} will allow you to \fillin{好处是什么} and leave you with \fillin{邀请收获}.\\
最后，接受此次关于\fillin{题目相关词}的邀请，会让你\fillin{好处是什么}，并收获\fillin{邀请收获}。

3.\ Last but not least, by taking part in \fillin{题目相关词}, you will open up new possibilities to \fillin{好处是什么}.\\
最后，通过参与\fillin{题目相关词}，你将为\fillin{好处是什么}打开新的可能性。

4.\ Last but not least, joining \fillin{题目相关词} serves as a rare chance for you to \fillin{好处是什么}, which should not be missed.\\
最后，参与\fillin{题目相关词}是你\fillin{好处是什么}的难得机会，实在不应错过。

5.\ Last but not least, with your active participation in \fillin{题目相关词}, the goal of \fillin{好处是什么} will be much easier to achieve.\\
最后，只要你积极参与\fillin{题目相关词}，\fillin{好处是什么}的目标将更容易实现。

6.\ Taken together, your involvement in \fillin{题目相关词} will not only \fillin{好处之一} but also \fillin{好处之二}.\\
总的来看，你参与\fillin{题目相关词}，不仅会\fillin{好处之一}，还会\fillin{好处之二}。

7.\ Handled in this way, \fillin{题目相关词} will become a wonderful occasion for you to \fillin{好处是什么}.\\
如此一来，\fillin{题目相关词}将成为你\fillin{好处是什么}的绝佳契机。
\end{eg}


% =====================================================================
\subsection{道歉类}
% =====================================================================

\begin{tpl}{道歉类 · 第二段骨架}
\textbf{第 2 句}：承认过失行为，并说明它给对方造成了什么影响。\par
\textbf{第 3 句}：解释原因、承担责任——不推给客观因素，落点在「我本应……」上。\par
\textbf{第 4 句}：给出具体补救措施，说明今后如何避免。\par
（道歉类最容易失分的地方是「只道歉不补救」，第 4 句不能省。）
\end{tpl}

\subsubsection{第 2 句功能句（7 条）}

\begin{eg}
1.\ At this point, I must admit that \fillin{过失行为} concerning \fillin{题目相关词} has caused \fillin{产生的影响}.\\
在此，我必须承认，在\fillin{题目相关词}方面的\fillin{过失行为}已经造成了\fillin{产生的影响}。

2.\ More specifically, my improper handling of \fillin{题目相关词}, namely \fillin{过失行为}, has brought you \fillin{产生的影响}.\\
更具体地说，我对\fillin{题目相关词}处理不当，也就是\fillin{过失行为}，给你带来了\fillin{产生的影响}。

3.\ In fact, the way I dealt with \fillin{题目相关词} by \fillin{过失行为} fell far short of your reasonable expectations.\\
事实上，我通过\fillin{过失行为}来处理\fillin{题目相关词}的方式，远远达不到你合理的期待。

4.\ To be honest, my decision to \fillin{过失行为} in relation to \fillin{题目相关词} turned out to be inconsiderate and unprofessional.\\
说实话，我在\fillin{题目相关词}上作出\fillin{过失行为}的决定，确实欠考虑，也不够专业。

5.\ As a result of \fillin{过失行为} with regard to \fillin{题目相关词}, you had to face \fillin{产生的影响}, which should not have happened.\\
由于我在\fillin{题目相关词}方面\fillin{过失行为}，你不得不面对\fillin{产生的影响}，这本不该发生。

6.\ When looking back, it is clear that \fillin{过失行为} about \fillin{题目相关词} was a serious mistake on my part.\\
回头来看，我在\fillin{题目相关词}上\fillin{过失行为}，无疑是我这边的严重失误。

7.\ In connection with \fillin{题目相关词}, my act of \fillin{过失行为} directly led to \fillin{产生的影响}, which I deeply regret.\\
在\fillin{题目相关词}这一问题上，我\fillin{过失行为}的做法直接导致了\fillin{产生的影响}，对此我深感遗憾。
\end{eg}

\subsubsection{第 3 句功能句（7 条）}

\begin{eg}
1.\ Although \fillin{客观原因} partly explains why this happened, the main responsibility still lies with me.\\
尽管\fillin{客观原因}在一定程度上解释了事情发生的原因，但主要责任仍然在我。

2.\ Admittedly, factors such as \fillin{客观原因} played a role, yet I should have foreseen the risk in \fillin{题目相关词}.\\
诚然，\fillin{客观原因}等因素确实起到了一定作用，但我本应预见到\fillin{题目相关词}中的风险。

3.\ Instead of simply blaming \fillin{客观原因}, I must acknowledge my lack of \fillin{主观原因，如 proper planning} in dealing with \fillin{题目相关词}.\\
与其把问题简单归咎于\fillin{客观原因}，我更应该承认自己在处理\fillin{题目相关词}时缺乏\fillin{主观原因，如 proper planning}。

4.\ In retrospect, if I had \fillin{本应做到的事} when handling \fillin{题目相关词}, the situation would not have become so serious.\\
事后反思，如果我在处理\fillin{题目相关词}时能够\fillin{本应做到的事}，情况就不会发展到如此严重的地步。

5.\ What matters more is that my failure to \fillin{本应做到的事} about \fillin{题目相关词} clearly shows my negligence.\\
更为重要的是，我没能在\fillin{题目相关词}上及时\fillin{本应做到的事}，这清楚地反映出我的疏忽。

6.\ Beyond any doubt, \fillin{主观原因，如 my carelessness} in relation to \fillin{题目相关词} is the real reason behind the problem.\\
毫无疑问，在\fillin{题目相关词}上的\fillin{主观原因，如 my carelessness}才是这一问题的真正根源。

7.\ From this experience, I have learned a lesson I must take seriously: \fillin{教训，用名词短语，如 the importance of being more cautious with ...}.\\
通过这次经历，我意识到\fillin{教训，如 I must be more cautious with ...}这一教训必须认真吸取。
\end{eg}

\subsubsection{第 4 句功能句（7 条）}

\begin{eg}
1.\ Last but not least, \fillin{补救措施} will be carried out immediately so as to reduce the negative impact of \fillin{题目相关词}.\\
最后，我会立刻落实\fillin{补救措施}，以减轻\fillin{题目相关词}带来的负面影响。

2.\ Last but not least, I have already planned to \fillin{补救措施}, which I hope can partly make up for the loss caused by \fillin{题目相关词}.\\
最后，我已着手计划\fillin{补救措施}，希望能够在一定程度上弥补\fillin{题目相关词}所造成的损失。

3.\ Last but not least, concrete steps such as \fillin{补救措施} will be taken, ensuring that similar problems with \fillin{题目相关词} will not occur again.\\
最后，我将采取\fillin{补救措施}等具体举措，确保今后在\fillin{题目相关词}上不再发生类似问题。

4.\ Last but not least, by \fillin{补救措施}, I will try my best to restore \fillin{需要挽回的方面，如 your trust} in the context of \fillin{题目相关词}.\\
最后，通过\fillin{补救措施}，我会尽最大努力在\fillin{题目相关词}方面挽回\fillin{需要挽回的方面，如 your trust}。

5.\ Last but not least, \fillin{补救措施} has been put on my priority list, which will help avoid repeating the same mistake in \fillin{题目相关词}.\\
最后，我已将\fillin{补救措施}列为首要任务，这将有助于避免今后在\fillin{题目相关词}上重蹈覆辙。

6.\ Taken together, \fillin{补救措施} and my renewed attention to \fillin{题目相关词} are meant to prevent such an incident from happening again.\\
总的来说，通过\fillin{补救措施}以及我对\fillin{题目相关词}重新投入的重视，目的是防止类似事件再次发生。

7.\ From now on, \fillin{补救措施} will be strictly implemented whenever \fillin{题目相关词} is involved, so that the same problem can be thoroughly avoided.\\
从现在起，每当涉及\fillin{题目相关词}时，我都会严格执行\fillin{补救措施}，以彻底避免同样的问题。
\end{eg}


% =====================================================================
\subsection{祝贺类}
% =====================================================================

\begin{tpl}{祝贺类 · 第二段骨架}
\textbf{第 2 句}：具体夸成绩——把对方的成就说清楚，并指出它印证了什么品质。\par
\textbf{第 3 句}：归因——这份成功来自努力、习惯、坚持或自律。\par
\textbf{第 4 句}：展望未来——这一成就将成为新起点。\par
（祝贺类的关键是「夸得具体」，只写 Congratulations 会显得敷衍。）
\end{tpl}

\subsubsection{第 2 句功能句（7 条）}

\begin{eg}
1.\ In particular, your outstanding performance in \fillin{题目相关词} fully proves that you have \fillin{优秀品质或能力}.\\
具体而言，你在\fillin{题目相关词}上的出色表现，充分证明了你具备\fillin{优秀品质或能力}。

2.\ More specifically, the fact that you have achieved \fillin{取得的成就} in \fillin{题目相关词} clearly shows your \fillin{品质或能力}.\\
更具体地说，你在\fillin{题目相关词}方面取得\fillin{取得的成就}，清楚地展现了你的\fillin{品质或能力}。

3.\ As far as \fillin{题目相关词} is concerned, your success in \fillin{取得的成就} is truly impressive and inspiring.\\
就\fillin{题目相关词}而言，你在\fillin{取得的成就}上的成功，确实令人钦佩，也颇具鼓舞意义。

4.\ In relation to \fillin{题目相关词}, your ability to \fillin{具体表现} strongly confirms your \fillin{品质或能力}.\\
在\fillin{题目相关词}方面，你能够\fillin{具体表现}，这有力地印证了你的\fillin{品质或能力}。

5.\ With such a remarkable outcome in \fillin{题目相关词}, it is evident that you have made persistent efforts in \fillin{付出或过程}.\\
在\fillin{题目相关词}上取得如此优异的结果，足以说明你在\fillin{付出或过程}方面付出了长期不懈的努力。

6.\ Seen from the perspective of \fillin{题目相关词}, your success vividly reflects \fillin{品质或努力过程}, which deserves full recognition.\\
从\fillin{题目相关词}的角度来看，你的成功生动体现了\fillin{品质或努力过程}，理应得到充分肯定。

7.\ At this point, it is fair to say that your achievement in \fillin{题目相关词} has set a good example for \fillin{受众对象}.\\
此时，可以毫不夸张地说，你在\fillin{题目相关词}上的成就已经为\fillin{受众对象}树立了良好榜样。
\end{eg}

\subsubsection{第 3 句功能句（7 条）}

\begin{eg}
1.\ To a large extent, this success comes from your willingness to \fillin{关键努力}, especially when dealing with \fillin{题目相关词}.\\
在很大程度上，这一成功源于你愿意在\fillin{题目相关词}上\fillin{关键努力}。

2.\ Moreover, it is your habit of \fillin{良好习惯} in \fillin{题目相关词} that lays a solid foundation for such an achievement.\\
再者，正是你在\fillin{题目相关词}中形成的\fillin{良好习惯}，为这一成就奠定了坚实基础。

3.\ In addition, your persistence in \fillin{长期坚持的行为} regarding \fillin{题目相关词} has turned difficulties into opportunities.\\
另外，你在\fillin{题目相关词}上坚持\fillin{长期坚持的行为}，把一次次困难转化为机遇。

4.\ What is more, your readiness to \fillin{积极行为} whenever \fillin{题目相关词} is involved has greatly contributed to this outcome.\\
更重要的是，每当涉及\fillin{题目相关词}时，你总是愿意\fillin{积极行为}，这对取得这一结果起到了关键作用。

5.\ Besides, the way you \fillin{具体做法} in \fillin{题目相关词} demonstrates a mature attitude towards \fillin{相关价值}.\\
此外，你在\fillin{题目相关词}中\fillin{具体做法}的方式，也展现出你对\fillin{相关价值}的成熟态度。

6.\ Furthermore, your ability to balance \fillin{题目相关词} with \fillin{其他任务或责任} shows impressive self-discipline.\\
此外，你能在\fillin{题目相关词}与\fillin{其他任务或责任}之间取得平衡，体现出非同一般的自律能力。

7.\ More importantly, through \fillin{题目相关词}, you have gradually learned to \fillin{学习到的能力或品质}, which makes this success even more meaningful.\\
更为重要的是，通过\fillin{题目相关词}，你逐渐学会了\fillin{学习到的能力或品质}，这也使得这次成功更加弥足珍贵。
\end{eg}

\subsubsection{第 4 句功能句（7 条）}

\begin{eg}
1.\ Last but not least, this achievement in \fillin{题目相关词} will surely encourage you to \fillin{未来行动}, helping you reach \fillin{下一阶段目标}.\\
最后，\fillin{题目相关词}上的这一成就必将激励你\fillin{未来行动}，助你迈向\fillin{下一阶段目标}。

2.\ Last but not least, with this success as a new starting point in \fillin{题目相关词}, you are well positioned to \fillin{未来发展方向}.\\
最后，以这次在\fillin{题目相关词}上的成功为新起点，你已具备良好条件去\fillin{未来发展方向}。

3.\ Last but not least, the experience gained from \fillin{题目相关词} will pave the way for you to \fillin{未来收获或成就}.\\
最后，你在\fillin{题目相关词}中积累的经验，将为今后实现\fillin{未来收获或成就}铺平道路。

4.\ Taken together, this milestone in \fillin{题目相关词} not only rewards your past efforts but also opens up broader space for \fillin{未来成长}.\\
总的来看，\fillin{题目相关词}上的这一里程碑，不仅是对你过去努力的回报，也为\fillin{未来成长}打开了更广阔的空间。

5.\ From now on, with the confidence built through \fillin{题目相关词}, it will be easier for you to \fillin{未来挑战或目标}.\\
从现在起，凭借你在\fillin{题目相关词}中树立的自信，你将更从容地应对\fillin{未来挑战或目标}。

6.\ Viewed in a long-term perspective, this success in \fillin{题目相关词} will become a lasting source of motivation for you to \fillin{长期目标}.\\
从长远来看，你在\fillin{题目相关词}上的这次成功，将成为你持续追求\fillin{长期目标}的重要动力源泉。

7.\ Under such a promising start in \fillin{题目相关词}, it is reasonable to expect that you will soon \fillin{更高层次的成就}.\\
在\fillin{题目相关词}如此良好的开端之下，我们有理由期待你不久后能\fillin{更高层次的成就}。
\end{eg}


% =====================================================================
\subsection{介绍类}
% =====================================================================

\begin{tpl}{介绍类 · 第二段骨架}
\textbf{第 2 句}：给出定义或定位——「这是什么 / 它本质上是什么」。\par
\textbf{第 3 句}：展开特点与作用——包含哪些方面、带来什么好处。\par
\textbf{第 4 句}：总结价值——为什么它值得选择。\par
（介绍类含「我的日常生活」式小题，第 2 句的第 8～10 条专为此类准备。）
\end{tpl}

\subsubsection{第 2 句功能句（10 条）}

\begin{eg}
1.\ To be more specific, \fillin{题目相关词} refers to \fillin{介绍对象或概念}, which is characterized by \fillin{核心特点}.\\
更具体而言，\fillin{题目相关词}是指\fillin{介绍对象或概念}，其突出特点在于\fillin{核心特点}。

2.\ In essence, \fillin{题目相关词} can be understood as \fillin{概括性定义}, with a special focus on \fillin{重点方面}.\\
从本质上看，\fillin{题目相关词}可以理解为\fillin{概括性定义}，尤其强调\fillin{重点方面}。

3.\ As a whole, \fillin{题目相关词} serves as \fillin{功能或作用}, providing \fillin{受众对象} with \fillin{具体益处}.\\
总体而言，\fillin{题目相关词}起到\fillin{功能或作用}的作用，为\fillin{受众对象}提供\fillin{具体益处}。

4.\ In terms of its basic nature, \fillin{题目相关词} combines \fillin{要素1} with \fillin{要素2}, forming a unique pattern.\\
就其基本属性而言，\fillin{题目相关词}融合了\fillin{要素1}与\fillin{要素2}，形成了独特格局。

5.\ From a practical perspective, \fillin{题目相关词} mainly involves \fillin{主要内容}, aiming to \fillin{总体目标}.\\
从实践角度看，\fillin{题目相关词}主要涵盖\fillin{主要内容}，其根本目标在于\fillin{总体目标}。

6.\ Viewed broadly, \fillin{题目相关词} functions as \fillin{角色定位}, playing an important part in \fillin{应用场景}.\\
从更广泛的角度看，\fillin{题目相关词}扮演\fillin{角色定位}的角色，在\fillin{应用场景}中发挥着重要作用。

7.\ In its simplest form, \fillin{题目相关词} can be described as \fillin{简要描述}, which distinguishes it from \fillin{对比对象}.\\
用最简洁的方式来说，\fillin{题目相关词}可以描述为\fillin{简要描述}，这使它有别于\fillin{对比对象}。

8.\ More specifically, in my daily life I usually \fillin{日常习惯1} and \fillin{日常习惯2}, which helps me maintain \fillin{生活相关优点，如 a balanced and healthy lifestyle}.\\
更具体而言，在日常生活中，我通常会\fillin{日常习惯1}并且\fillin{日常习惯2}，这有助于我保持\fillin{生活相关优点，如 a balanced and healthy lifestyle}。

9.\ As for my daily routine, I tend to \fillin{作息习惯，如 get up early} and \fillin{学习或锻炼习惯，如 set aside time for regular exercise}, which enables me to stay \fillin{状态，如 energetic and focused}.\\
至于我的日常作息，我通常会\fillin{作息习惯，如 get up early}，并\fillin{学习或锻炼习惯，如 set aside time for regular exercise}，这使我始终保持\fillin{状态，如 energetic and focused}。

10.\ In my everyday life, I not only \fillin{学习习惯1} but also \fillin{兴趣习惯2}, and this combination shapes \fillin{生活方式优点，如 a positive and self-disciplined lifestyle}.\\
在日常生活中，我既会\fillin{学习习惯1}，也会\fillin{兴趣习惯2}，而这种组合塑造了\fillin{生活方式优点，如 a positive and self-disciplined lifestyle}。
\end{eg}

\subsubsection{第 3 句功能句（7 条）}

\begin{eg}
1.\ To begin with, it covers \fillin{方面1}, enabling \fillin{受众对象} to \fillin{具体收获}.\\
首先，它涵盖\fillin{方面1}，使\fillin{受众对象}能够\fillin{具体收获}。

2.\ Moreover, its emphasis on \fillin{方面2} helps to \fillin{具体作用}, which is especially valuable for \fillin{特定群体}.\\
此外，它对\fillin{方面2}的重视，有助于\fillin{具体作用}，对\fillin{特定群体}尤为重要。

3.\ In addition, by integrating \fillin{内容1} with \fillin{内容2}, \fillin{题目相关词} provides a comprehensive way to \fillin{整体效果}.\\
另外，通过将\fillin{内容1}与\fillin{内容2}结合起来，\fillin{题目相关词}为\fillin{整体效果}提供了一种较为完整的途径。

4.\ What is more, the design of \fillin{题目相关词} highlights \fillin{亮点设计}, making it easier to \fillin{使用或理解方面的好处}.\\
更重要的是，\fillin{题目相关词}的设计突出\fillin{亮点设计}，从而更便于\fillin{使用或理解方面的好处}。

5.\ Besides, its flexible form allows it to be adapted to \fillin{应用场景1} and \fillin{应用场景2} without much difficulty.\\
此外，其形式较为灵活，可以较为容易地适用于\fillin{应用场景1}和\fillin{应用场景2}。

6.\ Furthermore, the clear structure of \fillin{题目相关词} makes it convenient for users to \fillin{使用动作}, improving \fillin{提升点}.\\
再者，\fillin{题目相关词}结构清晰，便于使用者\fillin{使用动作}，从而提升\fillin{提升点}。

7.\ On top of that, \fillin{题目相关词} offers practical guidance on \fillin{具体内容}, which can effectively support \fillin{支持对象}.\\
并且，\fillin{题目相关词}还就\fillin{具体内容}提供了实践性指引，可有效支撑\fillin{支持对象}。
\end{eg}

\subsubsection{第 4 句功能句（7 条）}

\begin{eg}
1.\ Taken together, these features make \fillin{题目相关词} an ideal choice for \fillin{目标人群} who wish to \fillin{目标需求}.\\
总的来看，这些特征使\fillin{题目相关词}成为\fillin{目标人群}实现\fillin{目标需求}的理想选择。

2.\ In this sense, \fillin{题目相关词} not only clarifies \fillin{帮助理解的内容} but also strengthens \fillin{延伸价值}.\\
从这个意义上讲，\fillin{题目相关词}既有助于厘清\fillin{帮助理解的内容}，也强化了\fillin{延伸价值}。

3.\ As a result, \fillin{题目相关词} serves as a reliable tool for \fillin{使用目的}, bringing long-term benefits to \fillin{受益者}.\\
因此，\fillin{题目相关词}成为实现\fillin{使用目的}的可靠工具，为\fillin{受益者}带来长远收益。

4.\ In summary, with its focus on \fillin{核心优势}, \fillin{题目相关词} can effectively support \fillin{具体任务或目标}.\\
总之，凭借其在\fillin{核心优势}方面的突出表现，\fillin{题目相关词}能够有力支撑\fillin{具体任务或目标}。

5.\ Under such a description, it is clear that \fillin{题目相关词} can play a constructive role in \fillin{应用方向}.\\
在上述介绍之下，不难看出，\fillin{题目相关词}能够在\fillin{应用方向}中发挥积极作用。

6.\ Viewed in this light, \fillin{题目相关词} is far more than \fillin{表面印象}, but a meaningful approach to \fillin{深层价值}.\\
由此可见，\fillin{题目相关词}远不止于\fillin{表面印象}，更是一种实现\fillin{深层价值}的重要途径。

7.\ Altogether, \fillin{题目相关词} stands out for \fillin{突出优点}, making it highly valuable in \fillin{适用领域}.\\
综上，\fillin{题目相关词}因\fillin{突出优点}而脱颖而出，在\fillin{适用领域}中具有较高价值。
\end{eg}


% =====================================================================
\subsection{通知类}
% =====================================================================

\begin{tpl}{通知类 · 第二段骨架}
\textbf{第 2 句}：交代基本要素——时间 / 地点 / 参与对象 / 主题 / 主办单位。\par
\textbf{第 3 句}：交代主要内容——有哪些环节、以什么形式开展。\par
\textbf{第 4 句}：交代注意事项——需要携带什么、遵守什么要求。\par
（通知是唯一允许写得规整并列的文体：\kw{To begin with / As for / Moreover / Last but not least} 直接铺开即可，不必强求非谓语结构。中间段以\textbf{事件本体}为主，目的、意义、影响不要写太多。）
\end{tpl}

\subsubsection{第 2 句功能句（6 条）}

\begin{eg}
1.\ First, \fillin{时间信息} has been set as the official time for \fillin{题目相关词}, and all \fillin{参加对象} are required to attend on time.\\
首先，\fillin{时间信息}被确定为\fillin{题目相关词}的正式时间，届时要求全体\fillin{参加对象}准时参加。

2.\ To begin with, \fillin{题目相关词} will take place at \fillin{地点信息}, providing a fixed venue for all \fillin{参加对象}.\\
首先，\fillin{题目相关词}将在\fillin{地点信息}举行，为全体\fillin{参加对象}提供固定场所。

3.\ In terms of schedule, \fillin{题目相关词} is arranged from \fillin{起始时间} to \fillin{结束时间}, during which \fillin{主要环节} will be carried out.\\
在时间安排上，\fillin{题目相关词}定于\fillin{起始时间}至\fillin{结束时间}进行，期间将开展\fillin{主要环节}。

4.\ As for the participants, \fillin{题目相关词} mainly targets \fillin{参加对象}, who are kindly asked to follow the arrangements below.\\
至于参与人员，\fillin{题目相关词}主要面向\fillin{参加对象}，请相关人员注意以下安排。

5.\ With regard to the theme, \fillin{题目相关词} will focus on \fillin{活动主题或内容}, aiming to address \fillin{目标问题}.\\
在主题方面，\fillin{题目相关词}将围绕\fillin{活动主题或内容}展开，旨在回应\fillin{目标问题}。

6.\ From the organizational point of view, \fillin{题目相关词} is jointly held by \fillin{主办单位} and \fillin{协办单位}, ensuring smooth coordination.\\
从组织层面来看，\fillin{题目相关词}由\fillin{主办单位}与\fillin{协办单位}联合举办，以保障整体协调有序。
\end{eg}

\subsubsection{第 3 句功能句（6 条）}

\begin{eg}
1.\ Furthermore, the main contents of \fillin{题目相关词} include \fillin{内容1}, \fillin{内容2} and \fillin{内容3}, forming a clear structure.\\
此外，\fillin{题目相关词}的主要内容包括\fillin{内容1}、\fillin{内容2}和\fillin{内容3}，整体结构清晰明了。

2.\ Moreover, \fillin{题目相关词} will be conducted in the form of \fillin{形式1}, followed by \fillin{形式2}, so as to enhance \fillin{预期效果}.\\
再者，\fillin{题目相关词}将采用\fillin{形式1}的方式开展，并配合\fillin{形式2}，以增强\fillin{预期效果}。

3.\ In addition, key parts such as \fillin{重点环节1} and \fillin{重点环节2} will be highlighted to ensure substantial gains for \fillin{参加对象}.\\
另外，活动将重点突出\fillin{重点环节1}和\fillin{重点环节2}，以确保\fillin{参加对象}切实有所收获。

4.\ What is more, \fillin{题目相关词} will integrate \fillin{要素1} with \fillin{要素2}, offering participants a comprehensive experience.\\
更重要的是，\fillin{题目相关词}将融合\fillin{要素1}与\fillin{要素2}，为参与者提供较为全面的体验。

5.\ Besides, during \fillin{题目相关词}, \fillin{具体安排，如 group discussion or Q and A} will be arranged to encourage active participation.\\
此外，在\fillin{题目相关词}过程中将安排\fillin{具体安排，如 group discussion or Q and A}，以鼓励大家积极参与。

6.\ At the same time, sufficient time will be reserved for \fillin{环节，如 free communication}, so that \fillin{参加对象} can \fillin{具体收获}.\\
同时，将预留充足时间用于\fillin{环节，如 free communication}，以便\fillin{参加对象}能够\fillin{具体收获}。
\end{eg}

\subsubsection{第 4 句功能句（6 条）}

\begin{eg}
1.\ Last but not least, all \fillin{参加对象} are required to \fillin{注意事项或要求}, so that \fillin{题目相关词} can proceed in an orderly way.\\
最后，要求全体\fillin{参加对象}\fillin{注意事项或要求}，以确保\fillin{题目相关词}有序开展。

2.\ Last but not least, participants should \fillin{携带物品或准备事项}, which is necessary for the smooth progress of \fillin{题目相关词}.\\
最后，请参与人员\fillin{携带物品或准备事项}，这是确保\fillin{题目相关词}顺利进行的必要条件。

3.\ Last but not least, please pay special attention to \fillin{纪律或安全要求}, so as to maintain a safe and tidy environment for \fillin{题目相关词}.\\
最后，请特别注意\fillin{纪律或安全要求}，以维持\fillin{题目相关词}的安全、有序环境。

4.\ Taken together, the above arrangements call for your active cooperation, without which \fillin{题目相关词} cannot achieve the expected results.\\
总的来看，以上安排需要大家积极配合，否则\fillin{题目相关词}难以达到预期效果。

5.\ From now on, any changes related to \fillin{题目相关词} will be updated through \fillin{通知渠道}, and all \fillin{参加对象} should keep informed in time.\\
自即日起，凡与\fillin{题目相关词}有关的变动将通过\fillin{通知渠道}发布，请全体\fillin{参加对象}及时关注。

6.\ Under these requirements, \fillin{题目相关词} is expected to run efficiently and create real value for \fillin{参加对象}.\\
在上述要求下，\fillin{题目相关词}有望高效运行，并为\fillin{参加对象}带来实质性收获。
\end{eg}


% =====================================================================
\subsection{事件类第二段：六要素表达（时间 / 地点 / 身份 / 特征 / 起因 / 经过 / 结果）}
% =====================================================================

\begin{note}
来自御风。思路是借小说六要素（时间、地点、人物、起因、经过、结果）穷举一个事件的所有描述角度，再把它们与各信件的专用表达拼起来，就能撑起第二段。上面的六类功能句解决「说什么」，这里的六要素解决「怎么把细节写具体」。使用时随手取用，不必全部用上。
\end{note}

\subsubsection{（1）时间（注意介词搭配）}

\begin{longtable}{L{9.4cm}L{6.6cm}}
\toprule
\textbf{英文} & \textbf{中文} \\
\midrule
\endhead
\bottomrule
\endlastfoot
the meeting will be held at 9 a.m. this Friday. & 会议将于本周五上午 9 点举行。 \\
is scheduled on May 12. & 定于 5 月 12 日（举行）。 \\
during the coming summer vacation & 在即将到来的暑假期间 \\
at the beginning of the new term & 在新学期开始时 \\
by the end of this month & 在本月底之前 \\
over the past few weeks & 在过去几周里 \\
at the appointed time & 在约定时间 \\
deadline & 截止日期 \\
in due course & 在适当的时候（补充） \\
at your earliest convenience & 在你方便的时候尽早（补充） \\
from 10 a.m. to 4 p.m. next Saturday & 下周六上午 10 点到下午 4 点（补充，2022 英二范文） \\
before July 1st & 在 7 月 1 日之前（补充，招募截止） \\
\end{longtable}

\subsubsection{（2）地点}

\begin{longtable}{L{9.4cm}L{6.6cm}}
\toprule
\textbf{英文} & \textbf{中文} \\
\midrule
\endhead
\bottomrule
\endlastfoot
at the school gate & 在学校大门口 \\
in front of the main building & 在主楼前 \\
in the school auditorium & 在学校礼堂 \\
in the lecture hall & 在报告厅／阶梯教室 \\
in the conference room & 在会议室 \\
at the student activity center & 在学生活动中心 \\
in the exhibition hall & 在展厅 \\
at the community service center & 在社区服务中心 \\
at the railway station / airport & 在火车站／机场 \\
in a nearby park & 在附近的公园 \\
in the central square & 在中心广场（补充，2022 英二美食节范文） \\
in the online meeting room & 在线上会议室（补充，2021 英二线上会议） \\
\end{longtable}

\subsubsection{（3）身份（人物）}

\begin{longtable}{L{9.4cm}L{6.6cm}}
\toprule
\textbf{英文} & \textbf{中文} \\
\midrule
\endhead
\bottomrule
\endlastfoot
undergraduate students & 本科生 \\
postgraduate students & 研究生 \\
freshmen & 大一新生 \\
senior students & 高年级学生 \\
international students & 留学生 \\
exchange students & 交换生 \\
applicants & 申请者 \\
candidates & 候选人；考生；应聘者 \\
participants & 参与者 \\
contestants / competitors & 参赛者 \\
volunteers & 志愿者 \\
organizers & 组织者 \\
committee members & 委员会成员 \\
local residents & 当地居民 \\
distinguished guests & 尊敬的嘉宾 \\
\end{longtable}

\subsubsection{（4）特征（形容人物与活动的品质词）}

\begin{note}
这些形容词用在第二段里指人——描述参与者、志愿者、候选人、嘉宾的品质；也可用来夸人（祝贺类、推荐类的第 2 句）。
\end{note}

\begin{longtable}{L{9.4cm}L{6.6cm}}
\toprule
\textbf{英文} & \textbf{中文} \\
\midrule
\endhead
\bottomrule
\endlastfoot
distinguished & 杰出的；卓越的；德高望重的 \\
responsible & 有责任心的 \\
reliable & 可靠的，可信赖的 \\
competent & 有能力的，能胜任的 \\
qualified & 合格的，胜任的 \\
experienced & 经验丰富的 \\
dedicated & 投入的，有奉献精神的 \\
committed & 尽心尽力的，坚定投入的 \\
diligent & 勤奋的 \\
conscientious & 认真负责的，一丝不苟的 \\
self-disciplined & 自律的 \\
motivated & 积极主动的，有动力的 \\
enthusiastic & 热情的 \\
proactive & 主动积极的 \\
cooperative & 有合作精神的 \\
considerate & 体贴的，考虑周到的 \\
patient & 有耐心的 \\
approachable & 平易近人的，容易接近的 \\
adaptable & 适应能力强的 \\
organized & 有条理的 \\
insightful & 有见地的，有洞察力的 \\
\end{longtable}

\subsubsection{（5）起因（交代活动为什么办、信为什么写）}

\begin{longtable}{L{9.4cm}L{6.6cm}}
\toprule
\textbf{英文} & \textbf{中文} \\
\midrule
\endhead
\bottomrule
\endlastfoot
against the background of... & 在……的背景下 \\
in light of... & 鉴于；考虑到 \\
in view of... & 鉴于；考虑到 \\
given the importance of... & 鉴于……的重要性 \\
with the aim of... & 旨在…… \\
as part of our effort to... & 作为我们努力……的一部分 \\
in an effort to... & 为了努力…… \\
in response to... & 为了回应…… \\
inspired by... & 受到……启发 \\
out of concern for... & 出于对……的关心 \\
due to the growing need for... & 由于对……日益增长的需求 \\
with growing attention paid to... & 随着人们越来越关注…… \\
as more importance is attached to... & 随着人们更加重视…… \\
in order to raise awareness of... & 为了增强对……的意识 \\
to create an opportunity for... & 为了给……创造机会 \\
\end{longtable}

\subsubsection{（6）经过（活动怎么推进、准备怎么展开）}

\begin{longtable}{L{9.4cm}L{6.6cm}}
\toprule
\textbf{英文} & \textbf{中文} \\
\midrule
\endhead
\bottomrule
\endlastfoot
in the course of the activity & 在活动过程中 \\
throughout the whole process & 在整个过程中 \\
during the preparation stage & 在准备阶段 \\
at the initial stage & 在初始阶段 \\
as the activity progressed & 随着活动的推进 \\
after careful consideration & 经过认真考虑后 \\
after several rounds of discussion & 经过几轮讨论后 \\
following your suggestion & 按照你的建议；在听取你的建议后 \\
under your guidance & 在你的指导下 \\
with the support of the organizers & 在组织者的支持下 \\
with the joint efforts of all participants & 在所有参与者的共同努力下 \\
as previously arranged & 按照此前安排 \\
as scheduled & 按计划；如期 \\
meanwhile & 与此同时 \\
simultaneously & 同时 \\
subsequently & 随后；其后 \\
shortly afterwards & 不久之后 \\
in the following days & 在接下来的几天里 \\
in an orderly manner & 有序地 \\
despite some unexpected difficulties & 尽管遇到一些意外困难 \\
be arranged to & 被安排去（后接动词） \\
followed by & ……之后 \\
\end{longtable}


% =====================================================================
\subsection{各类信件「结果 / 效果」句}
% =====================================================================

\begin{note}
来自御风。第二段（或第三段）写到最后总要落到一个「结果 / 效果」上，这张表按信件类型给出了固定的落点。答题时\textbf{逐点回应、句句有回应}：题目提到的每一个要点都要有对应句子，不重不漏不跑题。
\end{note}

\begin{longtable}{L{1.8cm}L{9.5cm}L{4.4cm}}
\toprule
\textbf{信件类型} & \textbf{结果 / 效果表达（英文）} & \textbf{中文} \\
\midrule
\endfirsthead
\toprule
\textbf{信件类型} & \textbf{结果 / 效果表达（英文）} & \textbf{中文} \\
\midrule
\endhead
\bottomrule
\endlastfoot

感谢信 & benefit me greatly & 使我受益匪浅 \\
感谢信 & leave a deep impression on me & 给我留下深刻印象 \\
感谢信 & enable me to make further progress & 使我能够取得更大进步 \\
感谢信 & provide practical guidance for me & 为我提供切实的指导 \\
\midrule
邀请信 & add great value to the event & 为活动增添重要价值 \\
邀请信 & make the activity more meaningful & 使活动更有意义 \\
邀请信 & benefit all participants & 使所有参与者受益 \\
邀请信 & promote mutual understanding & 促进相互理解 \\
\midrule
建议信 & be taken into consideration & 被纳入考虑 \\
建议信 & lead to a better outcome & 带来更好的结果 \\
建议信 & make the arrangement more effective & 使安排更有效 \\
建议信 & improve communication & 改善沟通 \\
\midrule
投诉信 & cause considerable inconvenience & 造成相当大的不便 \\
投诉信 & fail to meet my expectations & 未能达到我的期望 \\
投诉信 & affect my original plan & 影响我原定的计划 \\
投诉信 & be dealt with properly & 得到妥善处理 \\
\midrule
道歉信 & cause inconvenience to you & 给你带来不便 \\
道歉信 & affect the original arrangement & 影响原有的安排 \\
道歉信 & make up for the inconvenience & 弥补这一不便 \\
道歉信 & avoid similar problems in the future & 避免今后出现类似问题 \\
\midrule
申请信 & gain valuable experience & 获得宝贵经验 \\
申请信 & make further progress & 取得更大进步 \\
申请信 & contribute to the program & 为该项目做出贡献 \\
申请信 & live up to your expectations & 不辜负你的期望 \\
\midrule
祝贺信 & win wide recognition & 赢得广泛认可 \\
祝贺信 & achieve remarkable progress & 取得显著进步 \\
祝贺信 & deserve such an honor & 配得上这份荣誉 \\
祝贺信 & make greater achievements in the future & 今后取得更大成就 \\
\midrule
推荐信 & be fully qualified for the position & 完全胜任该职位 \\
推荐信 & make a valuable contribution & 做出宝贵贡献 \\
推荐信 & perform well in the program & 在该项目中表现出色 \\
推荐信 & meet your requirements & 符合你的要求 \\
\midrule
通知、活动安排 & achieve the expected result & 达到预期效果 \\
通知、活动安排 & be carried out smoothly & 顺利开展 \\
通知、活动安排 & receive positive feedback & 获得积极反馈 \\
通知、活动安排 & enrich students' campus life & 丰富学生的校园生活 \\
\midrule
求助、咨询信 & provide me with practical guidance & 为我提供切实指导 \\
求助、咨询信 & help me make a better decision & 帮助我做出更好的决定 \\
求助、咨询信 & clear up my confusion & 消除我的困惑 \\
求助、咨询信 & be of great help to my preparation & 对我的准备大有帮助 \\

\end{longtable}


% =====================================================================
\subsection{第二段整段示范}
% =====================================================================

\begin{tpl}{事件类示范 · 2022 英语一（邀请教授组队参加创新大赛）}
The contest, revolving around the theme of inspiring students' creativity, will be held at my university next month. Serving as a high-level platform to encourage campus innovation, it has attracted extraordinary teams from all over the world to sign up. And through engaging in this game, competitors not only gain a chance to immerse themselves in the atmosphere of collaboration and competition, but also enhance their practical skills in addressing real-world issues. Therefore, given your extensive experience in contests of this kind, we would like to extend an invitation to you.\\[2pt]
本次大赛以激发学生创造力为主题，将于下月在我校举行。作为鼓励校园创新的高水平平台，它已吸引来自世界各地的优秀团队报名。参与这场赛事，参赛者不仅能沉浸于协作与竞争的氛围之中，还能提升自己解决现实问题的实践能力。因此，鉴于您在这类赛事上的丰富经验，我们首先想到向您发出邀请。\\[3pt]
\textbf{逻辑骨架}：主题 + 时间（回应题目）$\rightarrow$ 平台介绍（脑补，不与题目冲突）$\rightarrow$ 效果与好处（效果）$\rightarrow$ 夸夸 + 邀请（回应题目）。
\end{tpl}

\begin{tpl}{非事件类示范 · 2023 英语二（在艺术展与机器人展之间做建议）}
The main reason is that the robot show will offer you a valuable chance to experience the latest development of science and technology. By watching different kinds of robots in action, you can not only broaden your knowledge, but also develop a deeper interest in innovation. Besides, since you have always been curious about modern technology, I believe this show will be more suitable for you than the art exhibition.\\[2pt]
主要原因是，机器人展将为你提供一个体验科技最新发展的宝贵机会。通过观看各类机器人的现场演示，你不仅能拓宽知识面，还能培养对创新更深的兴趣。此外，既然你一直对现代科技充满好奇，我认为这个展比艺术展更适合你。\\[3pt]
\textbf{逻辑骨架}：直接给理由 $\rightarrow$ 展开好处（not only ... but also）$\rightarrow$ 结合对方特点做对比收束。非事件类（建议信／请求信／道歉信／介绍信）的第二段写法与大作文第二段相同，但不必写那么复杂，可以带主观色彩。
\end{tpl}

\begin{note}
\textbf{上表中被合并或改写的条目说明}：

（1）原书建议类第 2 句的第 1、2 条只差一个动词（putting / cultivating），语义完全一致，本册保留为两条并注明可互换。

（2）原书 OCR 把 \textit{specific / recommendation / comprehensive / Furthermore / engaging / missed / first-hand} 等词识别成 \textit{spccific / rccommcndation / comprchensive / Furthcrmore / cngaging / inissed / first hand}，已全部还原。

（3）原书邀请类第 3 句第 6 条的 \textit{Furthcrmore} 有拼写错误，已更正为 \textit{Furthermore}；第 7 条的 \textit{participants} 在通知类第 4 条中 OCR 成 \textit{Participants}（首字母误大写），已还原。

（4）原书建议类第 4 句第 7 条只有 \textit{Only by ...} 一句、没有 Last but not least 前缀，属起首式变体，保留原样。

（5）御风的六要素表原为「时间 / 地点 / 人物 / 起因 / 经过 / 结果」六组，其中「特征」是原书单独列出的一组形容词，本册保留为独立一组，共七组，无词条删减。
\end{note}

% =====================================================================
% 80_letter_p3.tex
% 小作文（Part A 应用文）第三段功能句库
% 来源：《英语作文模板》（御风）——「书信类型 $\rightarrow$ 第一句 / 第二句」大表
%       （30 种信件类型）、通知类第三段 12 句、真题范文第三段；
%       《mmg 作文模板》三、通用小作文——「第三段功能句」21 条。
% =====================================================================

\section{小作文第三段功能句库}

\begin{note}
第三段只写 1～2 句，占全文 15～20 词：

\textbf{第 1 句「总结 + 期望」}——把前面说的事收个尾，说一句「我希望……」。

\textbf{第 2 句「期待回应 / 结果」}——用 \kw{Looking forward to ...} 收束，表明等待对方的回复、确认、理解或解决方案。

第三段是\textbf{「格式分 / 得体分」的主要来源}：格式对不对、语气得不得体，阅卷人往往就在这最后两句上定印象。号召式句子照样要贴合题目——给老师写信用 \kw{Looking forward to your kind understanding}，给朋友写信用 \kw{Looking forward to seeing you soon}，不能互换。
\end{note}

\begin{tpl}{第三段公式}
\textbf{第 1 句}：\kw{I / We sincerely hope that ...}（或 \kw{Once again, I would like to ...}）——总结 + 期望。\par
\textbf{第 2 句}：\kw{Looking forward to ...}——期待回应 / 结果。\par
（落款另起一行：\kw{Yours sincerely,} + \kw{Li Ming}；题目不要求写地址，就不要写。）
\end{tpl}

\rubric{第三段对应评分档位里的「\textbf{格式和语气恰当}」。称呼—正文—落款三件套齐全、\kw{Looking forward to} 之后的动词形式正确（looking forward to \textbf{doing}）、落款名字符合题目要求，这几点守住，格式分就不会丢。}

\zone{三个小雷区：（1）\kw{look forward to} 后面误接动词原形——必须是动名词 \kw{receiving / hearing / seeing}；（2）在第三段里又补一条新的写作要点，导致段落重心错乱；（3）落款署名写了真名或自创名字，题目要求用 Li Ming 就必须用 Li Ming。}


% =====================================================================
\subsection{第 1 句：总结与期望}
% =====================================================================

\begin{note}
下表完整收入御风「书信类型 $\rightarrow$ 第一句：总结 + 期望」的 30 种类型，无删减。中文为逐条补译。找不到完全对应的类型时，按\textbf{收信人身份}就近选：给老师 / 机构用正式版，给朋友用生活化版。
\end{note}

\begin{longtable}{L{2.2cm}L{9.0cm}L{4.7cm}}
\toprule
\textbf{信件类型} & \textbf{第一句：总结 + 期望（英文）} & \textbf{中文} \\
\midrule
\endfirsthead
\toprule
\textbf{信件类型} & \textbf{第一句：总结 + 期望（英文）} & \textbf{中文} \\
\midrule
\endhead
\bottomrule
\endlastfoot

邀请信 & On behalf of the organizing committee, we sincerely anticipate your attendance at this meaningful event. & 我谨代表组委会，诚挚期待您出席这场意义非凡的活动。 \\
邀请嘉宾信 & On behalf of the contest committee, we sincerely anticipate your attendance and valuable guidance. & 我谨代表赛事委员会，诚挚期待您的出席与宝贵指导。 \\
邀请团队信 & On behalf of the organizing committee, we sincerely anticipate your attendance with your talented team. & 我谨代表组委会，诚挚期待您携优秀团队出席。 \\
申请信 & I sincerely hope that my qualifications and enthusiasm will meet your requirements. & 我真诚希望我的资质与热情能够符合您的要求。 \\
求职信 & I sincerely hope that my experience and abilities will make me a suitable candidate for this position. & 我真诚希望我的经验与能力能使我成为该职位的合适人选。 \\
推荐信 & I sincerely believe that he/she is fully qualified for this opportunity and will perform well if accepted. & 我坚信他／她完全能胜任这次机会，若能获选定会表现出色。 \\
感谢信 & Once again, I would like to express my heartfelt gratitude for your generous help and support. & 再次对您慷慨的帮助与支持表达我由衷的感谢。 \\
道歉信 & Once again, I sincerely apologize for the inconvenience I have caused and hope you can accept my apology. & 再次为我造成的不便致以诚挚歉意，希望您能接受我的道歉。 \\
投诉信 & I sincerely hope that this matter can be dealt with properly and without further delay. & 我真诚希望此事能得到妥善处理，不再拖延。 \\
建议信 & I sincerely hope that my suggestions will be taken into careful consideration. & 我真诚希望我的建议能被认真考虑。 \\
咨询信 & I sincerely hope that you could provide me with more detailed information about the matter. & 我真诚希望您能就此事为我提供更详细的信息。 \\
求助信 & I sincerely hope that you could offer me some practical guidance on this matter. & 我真诚希望您能就此事给我一些切实的指导。 \\
祝贺信 & Once again, please accept my warmest congratulations on your remarkable achievement. & 再次请接受我对你所取得的卓越成就最热烈的祝贺。 \\
慰问信 & I sincerely hope that you will recover soon and return to your normal life with confidence. & 我真诚希望你早日康复，满怀信心地回归正常生活。 \\
建议信：校园生活 & I sincerely hope that these suggestions will help improve students' campus experience. & 我真诚希望这些建议有助于改善学生的校园体验。 \\
建议信：学习规划 & I sincerely hope that these suggestions will be helpful to your study and future preparation. & 我真诚希望这些建议对你的学习和今后的准备有所帮助。 \\
感谢信：老师指导 & Once again, I would like to express my sincere gratitude for your patient guidance. & 再次对您耐心的指导表达我诚挚的感谢。 \\
感谢信：朋友帮助 & Once again, I truly appreciate your timely help and warm support. & 再次衷心感谢你及时的帮助与温暖的支持。 \\
活动通知、告知信 & We sincerely hope that this activity will provide all participants with a valuable and enjoyable experience. & 我们真诚希望本次活动能为所有参与者带来宝贵而愉快的体验。 \\
志愿者招募信 & We sincerely hope that more responsible and enthusiastic students will join us in this meaningful activity. & 我们真诚希望更多有责任心、有热情的同学加入这项有意义的活动。 \\
失物招领信 & I sincerely hope that the owner can contact me soon and get the lost item back. & 我真诚希望失主能尽快联系我，取回遗失物品。 \\
寻物启事信 & I sincerely hope that anyone who has found the item can offer me timely help. & 我真诚希望拾到该物品的人能及时给我帮助。 \\
辞职信 & I sincerely hope that my resignation will not cause you too much inconvenience. & 我真诚希望我的辞职不会给您的安排带来太多不便。 \\
拒绝信 & I sincerely hope that my decision will not cause you any inconvenience. & 我真诚希望我的决定不会给您的安排带来任何不便。 \\
介绍信 & I sincerely hope that this introduction will help you gain a clearer understanding of the activity. & 我真诚希望这份介绍能帮助你更清楚地了解本次活动。 \\
询问活动安排 & I sincerely hope that you could help me clarify the relevant arrangements. & 我真诚希望您能帮我明确相关的安排。 \\
文化交流信 & We sincerely hope that this activity will serve as a bridge for mutual understanding and cultural exchange. & 我们真诚希望本次活动能成为相互理解与文化交流的桥梁。 \\
环保活动信 & We sincerely hope that this activity will raise students' awareness of environmental protection. & 我们真诚希望本次活动能提升学生的环保意识。 \\
读书活动信 & We sincerely hope that this activity will encourage more students to develop a lasting interest in reading. & 我们真诚希望本次活动能鼓励更多学生培养持久的阅读兴趣。 \\
比赛相关信 & We sincerely hope that the contest will offer all participants a valuable chance to show their abilities. & 我们真诚希望本次比赛能为所有参赛者提供展示能力的宝贵机会。 \\

\end{longtable}

\begin{note}
\textbf{中文说明合并如下}：这 30 句的句式高度统一，只有三个变体——\kw{I sincerely hope that ...}（个人身份，最常用）、\kw{We sincerely hope that ...}（代表集体，用于通知、活动、文化／环保／读书／比赛类）、\kw{Once again, I would like to ... / Once again, please accept ...}（感谢、道歉、祝贺、慰问类，用「再次」加强情感）。背句子时只要记\textbf{「主语 + sincerely hope that + 期望内容」}这一个骨架，替换中间的名词即可。
\end{note}


% =====================================================================
\subsection{第 2 句：期待回应}
% =====================================================================

\begin{note}
同样完整收入御风原表的 30 条第二句。这一列的定义是「期待回应 / 结果」，全部以 \kw{Looking forward to} 开头——注意它后面\textbf{一律接动名词}。中文为逐条补译。
\end{note}

\begin{longtable}{L{2.2cm}L{9.0cm}L{4.7cm}}
\toprule
\textbf{信件类型} & \textbf{第二句：期待回应 / 结果（英文）} & \textbf{中文} \\
\midrule
\endfirsthead
\toprule
\textbf{信件类型} & \textbf{第二句：期待回应 / 结果（英文）} & \textbf{中文} \\
\midrule
\endhead
\bottomrule
\endlastfoot

邀请信 & Looking forward to receiving your favorable reply at your earliest convenience. & 期待在你方便的时候尽早收到你的肯定答复。 \\
邀请嘉宾信 & Looking forward to receiving your confirmation at your earliest convenience. & 期待在你方便的时候尽早收到你的确认。 \\
邀请团队信 & Looking forward to receiving your early reply. & 期待早日收到你的回复。 \\
申请信 & Looking forward to your favorable consideration. & 期待得到你的肯定与考虑。 \\
求职信 & Looking forward to an opportunity for an interview. & 期待获得一次面试机会。 \\
推荐信 & Looking forward to your favorable consideration of my recommendation. & 期待你对我的推荐给予肯定考虑。 \\
感谢信 & Looking forward to having the opportunity to repay your kindness in the future. & 期待今后有机会报答你的善意。 \\
道歉信 & Looking forward to your kind understanding. & 期待你的谅解。 \\
投诉信 & Looking forward to receiving a satisfactory solution as soon as possible. & 期待尽快得到一个令人满意的解决方案。 \\
建议信 & Looking forward to seeing positive changes in the near future. & 期待在不久的将来看到积极的变化。 \\
咨询信 & Looking forward to receiving your detailed reply at your earliest convenience. & 期待在你方便的时候尽早收到你的详细回复。 \\
求助信 & Looking forward to receiving your helpful advice. & 期待收到你有所帮助的建议。 \\
祝贺信 & Looking forward to hearing more good news from you in the future. & 期待今后听到你更多的好消息。 \\
慰问信 & Looking forward to hearing about your full recovery soon. & 期待早日听到你完全康复的消息。 \\
建议信：校园生活 & Looking forward to seeing a better campus environment in the near future. & 期待在不久的将来看到一个更好的校园环境。 \\
建议信：学习规划 & Looking forward to hearing about your progress. & 期待听到你的进展。 \\
感谢信：老师指导 & Looking forward to learning more from you in the future. & 期待今后能继续向你学习。 \\
感谢信：朋友帮助 & Looking forward to meeting you again soon. & 期待不久后与你再见。 \\
活动通知、告知信 & Looking forward to your active participation. & 期待大家积极参与。 \\
志愿者招募信 & Looking forward to receiving your application. & 期待收到你的申请。 \\
失物招领信 & Looking forward to receiving further information from you. & 期待收到你的进一步消息。 \\
寻物启事信 & Looking forward to receiving any useful information. & 期待收到任何有用的信息。 \\
辞职信 & Looking forward to your kind understanding and approval. & 期待你的理解与批准。 \\
拒绝信 & Looking forward to your understanding. & 期待你的理解。 \\
介绍信 & Looking forward to your further interest and participation. & 期待你进一步关注并参与。 \\
询问活动安排 & Looking forward to receiving the detailed schedule. & 期待收到详细的日程安排。 \\
文化交流信 & Looking forward to your participation and support. & 期待你的参与与支持。 \\
环保活动信 & Looking forward to your active involvement. & 期待你积极参与。 \\
读书活动信 & Looking forward to your participation and valuable suggestions. & 期待你的参与和宝贵建议。 \\
比赛相关信 & Looking forward to your confirmation and support. & 期待你的确认与支持。 \\

\end{longtable}

\begin{note}
\textbf{中文说明合并如下}：30 条第二句只有四种落点——\kw{receiving ...}（期待收到回复、确认、申请、方案、日程）、\kw{hearing ...}（期待听到消息、进展、好消息、康复）、\kw{seeing ...}（期待看到变化、更好的环境）、\kw{your + 名词}（期待对方的理解、参与、支持、确认与配合）。选句时按\textbf{你希望对方做出什么动作}来定，不要一律套 receiving your reply。
\end{note}

\begin{tpl}{真题范文第三段两例}
\textbf{2022 英语一（邀请教授组队参赛）}\\
On behalf of the contest committee, we sincerely anticipate your attendance with your talented team. Looking forward to receiving your early reply.\\[2pt]
我谨代表赛事委员会，诚挚期待您携优秀团队出席。期待早日收到您的回复。\\[4pt]
\textbf{2012 英语二（投诉电子词典）}\\
Deeply disappointed though I am, I still believe that a responsible seller should attach due importance to the legitimate concerns of its customers. I would therefore appreciate it if your customer service center could look into this matter without delay and provide me with a reasonable and satisfactory solution.\\[2pt]
尽管我深感失望，但我仍然相信，一个负责任的商家应当重视顾客的正当诉求。因此，如果贵客服中心能尽快调查此事并给我一个合理、令人满意的解决方案，我将不胜感激。\\[3pt]
\textbf{提示}：第二例的第三段偏长（近 50 词），是「第三段也可以写得充分」的示范；考场上若第二段已经写满，第三段压缩回 2 句即可。
\end{tpl}


% =====================================================================
\subsection{通用万能第三段}
% =====================================================================

\begin{note}
考场上来不及按类型挑句子时，用下面 6 组万能收尾。每组都是「第 1 句 + 第 2 句」的完整组合，直接照抄，只把末尾的 \fillin{事由} 换成题目词。A、B、C 三组覆盖 90\% 的书信与通知。
\end{note}

\begin{eg}
\textbf{A 组 · 通用（书信首选，最百搭）}\par
Thank you for your time and consideration. I look forward to your reply at your earliest convenience.\\
感谢你抽出时间阅读与考虑，期待你尽早回复。\par

\textbf{B 组 · 通用（通知、告示首选）}\par
Thank you for your time and attention, and I look forward to your positive response.\\
感谢大家的时间与关注，期待大家给予积极回应。\par

\textbf{C 组 · 预先致谢（偏正式，适合请求、申请、投诉）}\par
Thank you in advance for your cooperation. Your support will be highly valued and sincerely appreciated.\\
预先感谢你的配合，你的支持将受到高度重视并深表感谢。\par

\textbf{D 组 · 期望效果（适合建议、邀请、通知）}\par
I sincerely hope that you will make full use of this opportunity. Once again, thank you for your consideration and support in this matter.\\
我真诚希望你能充分利用这次机会。再次感谢你对此事的关心与支持。\par

\textbf{E 组 · 共同愿景（适合活动、公益、集体类）}\par
With joint efforts, this matter is sure to bring about positive changes. Thank you again for your support.\\
相信在共同努力下，此事必将带来积极变化。再次感谢你的支持。\par

\textbf{F 组 · 极简版（字数吃紧时用，全文 18 词）}\par
Your kind consideration would be highly appreciated. Looking forward to your reply.\\
如蒙考虑，不胜感激；期待你的回复。
\end{eg}


% =====================================================================
\subsection{通知类第三段（御风 12 句）}
% =====================================================================

\begin{note}
通知、告示写 2～3 句收尾。主语按发通知的身份调整：学生会、组委会、学院发通知用 \kw{We / All participants / All students}；个人发通知改用 \kw{I / You are kindly requested to}。
\end{note}

\begin{eg}
1.\ Your attention and cooperation are sincerely appreciated.\\
感谢大家的关注与配合。

2.\ We sincerely appreciate your understanding and support.\\
我们真诚感谢大家的理解与支持。

3.\ Your cooperation with the above arrangements will be highly appreciated.\\
如能配合上述安排，我们将不胜感激。

4.\ All participants are expected to follow the above arrangements accordingly.\\
所有参与者应按上述安排执行。

5.\ Please make sure to arrive on time and observe the relevant requirements.\\
请务必准时到达，并遵守相关要求。

6.\ We look forward to your active participation.\\
期待大家积极参与。

7.\ All students are warmly encouraged to take part in this event.\\
热忱鼓励全体同学参加本次活动。

8.\ Should you have any questions, please feel free to contact the Students' Union.\\
如有疑问，请随时联系学生会。

9.\ We sincerely hope that this event will bring you a rewarding and memorable experience.\\
我们真诚希望本次活动能给大家带来充实而难忘的体验。

10.\ We believe that your participation will add great value to this event.\\
我们相信大家的参与将为本次活动增添重要价值。

11.\ We look forward to seeing you enjoy the event and benefit from it.\\
我们期待看到大家享受本次活动并从中受益。

12.\ May this event serve as a meaningful opportunity for learning, communication and personal growth.\\
愿本次活动成为大家学习、交流和个人成长的宝贵机会。
\end{eg}

\tip{\textbf{挑选建议（只选 2 句）}：第三段第 1 句选 \kw{3}（最正式）或 \kw{5}（含具体要求时最实）；第 2 句选 \kw{6}（最稳）或 \kw{12}（要显得有文采时用）。\kw{8}（留有疑问请联系）适合招募类、报名类通知，放在倒数第二句。\kw{4} 与 \kw{5} 是「要求句」，只有通知里真的规定了要求时才用，不要硬套。}

\tip{通知的\textbf{标题} Notice 居中大写，\textbf{时间}写在标题下方第一行最右侧（如 June 20th, 2022），\textbf{落款}右对齐写 The Students' Union。这三处格式比句子本身更容易被扣分。}


% =====================================================================
\subsection{第三段补充功能句（mmg 21 句）}
% =====================================================================

\begin{note}
原书「第三段功能句」21 条，本册全部收入。它们与上面御风的表示例用途一致，但更偏\textbf{通知、告示、集体类}，且大量使用 \fillin{此处填写题干事由} 占位（期待动作、时间节点、对象、总体目的），填空后即可直接用。第 1～5 条适合当第三段第 1 句（总结 + 期望），第 6～21 条适合当第 2 句（期待回应 / 结果）。
\end{note}

\begin{eg}
1.\ I would appreciate it very much if you could \fillin{期待动作} by \fillin{时间节点}.\\
如果你能在\fillin{时间节点}前\fillin{期待动作}，我将不胜感激。

2.\ It would be greatly appreciated if all \fillin{对象} could \fillin{期待动作} in accordance with the above arrangements.\\
若全体\fillin{对象}能够按照以上安排\fillin{期待动作}，将不胜感激。

3.\ You are kindly requested to \fillin{期待动作} so that \fillin{题目相关词} can be implemented more effectively.\\
请你（各位）\fillin{期待动作}，以便\fillin{题目相关词}能够得到更为有效的落实。

4.\ Please make sure that you \fillin{期待动作} before \fillin{时间节点}, ensuring that \fillin{题目相关词} proceeds smoothly.\\
请务必在\fillin{时间节点}前\fillin{期待动作}，以确保\fillin{题目相关词}顺利开展。

5.\ All \fillin{对象} are expected to \fillin{期待动作}, which is essential for the success of \fillin{题目相关词}.\\
希望全体\fillin{对象}\fillin{期待动作}，这一点对\fillin{题目相关词}的顺利开展至关重要。

6.\ Thank you for your time and attention, and I look forward to your positive response.\\
感谢你的时间与关注，期待你给予积极回应。

7.\ Thank you in advance for your cooperation with \fillin{题目相关词}.\\
预先感谢你在\fillin{题目相关词}方面的配合。

8.\ Many thanks for your support of \fillin{题目相关词}, and your understanding is highly appreciated.\\
衷心感谢你对\fillin{题目相关词}的支持，也非常感谢你的理解。

9.\ Your attention and cooperation will be highly valued and sincerely appreciated.\\
你的重视与配合将受到高度重视，并深表感谢。

10.\ Once again, thank you for your consideration and support in this matter.\\
再次感谢你对此事的关心与支持。

11.\ In summary, the above arrangements about \fillin{题目相关词} are made to \fillin{总体目的}.\\
总之，以上关于\fillin{题目相关词}的安排，旨在\fillin{总体目的}。

12.\ To conclude, all the measures mentioned above are designed to \fillin{总体目的}, and your participation is of great importance.\\
总而言之，上述各项措施均是为\fillin{总体目的}而设，你的参与至关重要。

13.\ Taken together, these efforts are expected to promote \fillin{题目相关词} and bring real benefits to \fillin{受益对象}.\\
综合来看，这些举措有望推进\fillin{题目相关词}，并为\fillin{受益对象}带来实实在在的收获。

14.\ Under the above arrangements, \fillin{题目相关词} is believed to play a more constructive role in \fillin{应用场景}.\\
在上述安排下，\fillin{题目相关词}有望在\fillin{应用场景}中发挥更加积极的作用。

15.\ I sincerely hope that you will make full use of this opportunity to \fillin{期待成效}.\\
我真诚希望你能充分利用这次机会\fillin{期待成效}。

16.\ It is my sincere hope that all \fillin{对象} can grow through \fillin{题目相关词} and achieve \fillin{期望结果}.\\
我衷心希望全体\fillin{对象}能够通过\fillin{题目相关词}得到成长，并实现\fillin{期望结果}。

17.\ With joint efforts, \fillin{题目相关词} is sure to bring about \fillin{积极变化}.\\
相信在共同努力下，\fillin{题目相关词}必将带来\fillin{积极变化}。

18.\ I firmly believe that your active involvement in \fillin{题目相关词} will lead to \fillin{积极结果}.\\
我坚信，你积极参与\fillin{题目相关词}，必然会促成\fillin{积极结果}。

19.\ Please read this notice carefully and pass it on to those concerned.\\
请认真阅读本通知，并转告相关人员。

20.\ All \fillin{对象} should keep a close eye on further notices related to \fillin{题目相关词}.\\
请全体\fillin{对象}密切关注后续关于\fillin{题目相关词}的相关通知。

21.\ If there is any adjustment to \fillin{题目相关词}, additional information will be released in due course.\\
如若\fillin{题目相关词}有任何调整，届时将另行通知。
\end{eg}

\begin{note}
\textbf{把三份来源合成一套的使用顺序}：

（1）先看题目是\textbf{通知}还是\textbf{书信}。通知 → 从「通知类第三段 12 句」里挑 2 句；书信 → 用「第 1 句 + 第 2 句」大表按类型挑。

（2）表里没有对应类型时用「通用万能第三段」的 A～F 组。

（3）字数还有富余，就从「mmg 21 句」里再补一句（第 1～5 条），或用第 6～21 条替换掉上面选的句子，避免多篇作文里第三段用句重复。

\textbf{合并说明}：御风通知第三段的第 1 句（Your attention and cooperation are sincerely appreciated.）与 mmg 第 9 条（Your attention and cooperation will be highly valued and sincerely appreciated.）语义高度重合，本册两条都保留，前者更简洁、后者更饱满，可任选。
\end{note}

% ============================================================
% 第 13 章 通用弹药库
% 来源：御风《英语作文模板》附录一（逻辑词）、十种句式、利弊短语、
%       Q2 趋势与程度词、Q3 数字基本知识；mmg《作文模板》话题词库
% ============================================================

\section{通用弹药库}

\begin{note}
这一章不给段落框架，只给\textbf{可以随时替换进任何框架里的词和句}。按御风的说法，一篇作文临场真正要决定的只有四样东西：\kw{衔接词}、\kw{论证词}、\kw{主题词}、\kw{主题短语}。衔接词和论证词是全题库通用的——因为你只需要写一篇文章；主题词则要靠“上位概念”（见本章第 5 小节）去覆盖。
\end{note}

\subsection{逻辑衔接词（15 类，每类精记 2--5 个）}

\zone{衔接词用错比不用更糟。御风原话：“有很多词看起来好像是可以用的，但是其实它的使用条件非常苛刻，一定要查清楚。”所以这里虽然是“弹药库”，也只背\textbf{你自己 100\% 会用}的那 3--5 个/类；表里带【注】的条目是原作者自己标出的使用限制，以及本次整理新发现的限制，\textbf{看看就好，不要上考场}。}

\par\addvspace{6pt}\noindent\textbf{1. 先后关系（Sequence / Chronological Order）}\par\addvspace{2pt}
\begin{longtable}{L{5.4cm}L{10.6cm}}
\toprule
\textbf{衔接词} & \textbf{英语二示范（短语或半句）} \\
\midrule
\endhead
at this time 此时 & At this time, the figure stood at 35\%. \\
first / first of all 首先 & First of all, the chart shows a clear upward trend in online learning. \\
second 其次 & Second, the proportion of rural users rose to 42\%. \\
next 接下来 & Next, the number climbed to 18 million. \\
previously 先前 & Previously, fewer than one in five families owned a car. \\
simultaneously 同时 & Simultaneously, the number of private cars also grew. \\
eventually 最终（结果） & Eventually, the figure leveled off at about 60\%. \\
last but not least 最后但同样重要的是 & Last but not least, personal effort should not be ignored. \\
to begin with / to start with 首先 & To begin with, the data reveal a widening gap between the two groups. \\
to end with 最后 & To end with, the two categories moved in opposite directions. \\
finally 最后 & Finally, the number returned to its 2015 level. \\
since then 从那时起 & Since then, the figure has risen year after year. \\
afterwards 之后 & Afterwards, the growth slowed down noticeably. \\
following this 在此之后 & Following this, spending on education kept rising. \\
preceding this 在此之前 & Preceding this, the figure had remained stable for a decade. \\
prior to this 在此之前 & Prior to this, few students chose to study abroad. \\
at last 最终 & At last, the two lines met at 40\%. \\
\bottomrule
\end{longtable}

\par\addvspace{6pt}\noindent\textbf{2. 因果关系（Cause and Effect）}\par\addvspace{2pt}
\begin{longtable}{L{5.4cm}L{10.6cm}}
\toprule
\textbf{衔接词} & \textbf{英语二示范（短语或半句）} \\
\midrule
\endhead
\textbf{表原因} & \\
because 因为 & The figure rose because average incomes increased. \\
because of this 由于这一点 & Because of this, more families could afford cultural activities.【注】原书列在“表原因”，但它实际跟在结果句后，属“承上”用法 \\
since 既然 / 因为 & Since online shopping is convenient, its share keeps growing. \\
as 由于 & As the population ages, demand for medical care rises. \\
for 因为（书面，放主句后） & The trend is understandable, for living standards have improved. \\
owing to 由于（多接消极或中性） & Owing to the fierce competition, young people face great pressure. \\
due to 由于 & The rise is largely due to government support. \\
for the reason that 原因是 & The change happened for the reason that policies were adjusted. \\
in view of 鉴于 & In view of these facts, the trend is easy to explain. \\
for such a reason 出于这一原因 & For such a reason, the two groups behaved differently. \\
as a result of 由于 & As a result of urbanization, the urban population doubled. \\
\midrule
\textbf{表结果} & \\
therefore 因此 & Therefore, the gap has narrowed. \\
consequently 因此（多接消极结果） & Consequently, traffic congestion became worse. \\
as a result 结果 & As a result, the number of private cars doubled. \\
thus 因而 & Thus, online learning has become part of daily life. \\
hence 因此（书面） & Hence, the figure has kept climbing since 2015. \\
so 所以 & So, more graduates choose flexible employment. \\
so that 以便 / 结果 & The government invested more so that rural schools could catch up. \\
in consequence 结果 & In consequence, the two curves moved in opposite directions. \\
as a consequence 因此 & As a consequence, cultural spending rose sharply. \\
accordingly 相应地 & Accordingly, the number of teachers also increased. \\
inevitably 不可避免地 & Inevitably, the gap between urban and rural areas widened. \\
under these conditions 在这种情况下 & Under these conditions, the trend is likely to continue. \\
\bottomrule
\end{longtable}

\par\addvspace{6pt}\noindent\textbf{3. 转折关系（Contrast / Concession）}\par\addvspace{2pt}
\begin{longtable}{L{5.4cm}L{10.6cm}}
\toprule
\textbf{衔接词} & \textbf{英语二示范（短语或半句）} \\
\midrule
\endhead
but 但是 & The figure rose, but the growth was uneven. \\
even so 即便如此 & Even so, the overall trend remained upward. \\
however 然而 & However, the increase was far from even. \\
though / even though 尽管 & Even though the policy was introduced, the change came slowly. \\
despite that / in spite of that 尽管如此 & In spite of that, the two groups still differed greatly. \\
regardless of 不管 / 不顾 & Regardless of income, people spend more on health. \\
independent of 不受……影响【注】 & The trend is independent of regional differences.【注】这是形容词短语，不是转折连词，考试按 but / however 用 \\
reckless of 不顾【注】 & Reckless of the cost, some consumers buy on impulse.【注】原文列入转折类，属归类宽松；非常用表达，建议不用 \\
while\ldots\ yet\ldots 虽然……但是…… & While the number rose, yet its share declined. \\
unless 除非 & Unless effective measures are taken, the gap will widen. \\
\bottomrule
\end{longtable}

\par\addvspace{6pt}\noindent\textbf{4. 并列关系（Addition / Parallel）}\par\addvspace{2pt}
\begin{longtable}{L{5.4cm}L{10.6cm}}
\toprule
\textbf{衔接词} & \textbf{英语二示范（短语或半句）} \\
\midrule
\endhead
and 和 & The chart shows the changes and the differences. \\
also 也 & The number of users also increased. \\
too 也（句末） & Cultural spending grew, too. \\
as well as 以及 & Online payment as well as live commerce has spread quickly. \\
either\ldots, or\ldots 或者……或者…… & Either the government or individuals should take action. \\
both\ldots\ and\ldots 既……又…… & Both quality education and lifelong learning matter. \\
for one thing\ldots, for another\ldots 一方面……另一方面…… & For one thing, incomes rose; for another, prices fell. \\
meanwhile 同时 & Meanwhile, the number of private cars kept growing. \\
at the same time 同时 & At the same time, the share of rural users increased. \\
\bottomrule
\end{longtable}

\par\addvspace{6pt}\noindent\textbf{5. 递进关系（Progression / Further Addition）}\par\addvspace{2pt}
\begin{longtable}{L{5.4cm}L{10.6cm}}
\toprule
\textbf{衔接词} & \textbf{英语二示范（短语或半句）} \\
\midrule
\endhead
furthermore 此外 & Furthermore, the trend is visible in every age group. \\
moreover 而且 & Moreover, the gap has widened since 2015. \\
further 进一步 & Further, online services have become more affordable. \\
in this way 这样 / 通过这种方式 & In this way, the two groups can share the same resources. \\
still 仍然 & Still, the figure remained below the 2019 level. \\
not only\ldots\ but also\ldots 不仅……而且…… & Not only did incomes rise, but spending also increased. \\
not\ldots\ but\ldots 不是……而是…… & The change is driven not by policy but by technology. \\
in addition (to) / additionally 此外 & In addition, the number of participants doubled. \\
more importantly 更重要的是 & More importantly, the trend reflects a change in attitudes.【注】原文为 much more interesting，并自注“通常作句子成分，较少单独作连接词”，故替换为 more importantly \\
more specifically 更具体地说 & More specifically, the share of people aged over 60 rose fastest. \\
next 其次 & Next, we should consider the cost of the change. \\
besides 此外 & Besides, the data cover only urban residents. \\
as far as\ldots\ is concerned 就……而言 & As far as education is concerned, the figures tell a different story. \\
in other words 换言之 & In other words, the gap has become wider.【注】更偏向重申，见第 14 类 \\
on the one hand\ldots; on the other hand\ldots 一方面……另一方面…… & On the one hand, it saves time; on the other hand, it reduces face-to-face contact. \\
even 甚至 & Even in small cities, online payment is now common. \\
as a popular saying goes\ldots 俗话说…… & As a popular saying goes, “Rome was not built in a day.” \\
in order to do it\ldots 为了做到这一点…… & In order to do it, both the government and individuals must act.【注】偏向目的关系，见第 13 类 \\
to make things better (worse) 更好（更糟）的是 & To make things worse, the gap is still widening. \\
accordingly 相应地 & Accordingly, the number of nursing homes also grew. \\
\bottomrule
\end{longtable}

\par\addvspace{6pt}\noindent\textbf{6. 比较关系 / 相同点（Comparison / Similarity）}\par\addvspace{2pt}
\begin{longtable}{L{5.4cm}L{10.6cm}}
\toprule
\textbf{衔接词} & \textbf{英语二示范（短语或半句）} \\
\midrule
\endhead
similarly 同样地 & Similarly, the number of rural users also rose. \\
in like manner 同样地（书面） & In like manner, the two cities followed the same path. \\
in comparison with / when compared with / compared with 与……相比 & Compared with 2015, the figure nearly doubled. \\
when in fact\ldots 而事实上…… & The chart looks positive, when in fact the growth has stalled. \\
like\ldots 像……一样 & Like developed countries, China now faces an aging population. \\
likewise 同样 & Likewise, spending on education grew faster than income. \\
similarly important 同样重要的是 & Similarly important is the change in people's attitudes. \\
apart from (doing)\ldots 除……之外 & Apart from saving time, online payment also reduces costs. \\
\ldots rather than\ldots 而不是…… & The rise is due to policy rather than to luck.【注】通常表对比或替代，也可用于比较 \\
by doing so 通过这样做 & By doing so, the government eased the pressure on families. \\
both\ldots\ and\ldots 既……又…… & Both men and women spent more on culture. \\
in the same way 同样地 & In the same way, the two categories grew at similar rates. \\
not only\ldots\ but (also)\ldots 不仅……而且…… & Not only did sales rise, but profits also improved. \\
\bottomrule
\end{longtable}

\par\addvspace{6pt}\noindent\textbf{7. 对照关系 / 不同点（Contrast / Difference）}\par\addvspace{2pt}
\begin{longtable}{L{5.4cm}L{10.6cm}}
\toprule
\textbf{衔接词} & \textbf{英语二示范（短语或半句）} \\
\midrule
\endhead
yet / still 然而 & Yet the gap between the two groups remained wide. \\
for all of that 尽管如此 & For all of that, the overall trend is encouraging. \\
notwithstanding 尽管（极书面） & Notwithstanding these efforts, the problem persists. \\
rather 相反 / 更确切地说 & It is not a failure; rather, it is a delay. \\
neither\ldots\ nor\ldots 既不……也不…… & Neither the young nor the old were satisfied with the change. \\
although / though 虽然 & Although the number rose, the share fell. \\
but / however 但是 & But the picture is not entirely positive. \\
something is just the other way around 情况恰恰相反 & For culture, the situation is just the other way around. \\
conversely 相反地 & Conversely, spending on fast food declined. \\
unlike / different from this 与……不同 & Unlike the previous decade, growth has slowed. \\
opposed to / as opposed to 与……相对 / 而不是 & They chose online courses as opposed to traditional ones. \\
in contrast 相比之下 & In contrast, rural users made up only 27\%. \\
by way of 以……方式 & By way of comparison, the figure was twice as high in 2010.【注】通常意为“通过……方式”，表对比不如 in contrast 常用 \\
on the contrary 恰恰相反 & On the contrary, the number kept falling. \\
nevertheless 尽管如此 & Nevertheless, the difference remains striking. \\
contrary to 与……相反 & Contrary to expectations, the share did not rise. \\
whereas / while 然而 / 而 & Reading stayed stable, whereas gardening declined. \\
\bottomrule
\end{longtable}

\par\addvspace{6pt}\noindent\textbf{8. 举例关系（Exemplification）}\par\addvspace{2pt}
\begin{longtable}{L{5.4cm}L{10.6cm}}
\toprule
\textbf{衔接词} & \textbf{英语二示范（短语或半句）} \\
\midrule
\endhead
for example / for instance 例如 & For instance, the number of online learners doubled. \\
in this case 在这种情况下 & In this case, the figure should be read with caution. \\
namely 即 & Only one group, namely students, chose that option. \\
as you know 如你所知（口语） & As you know, exams test much more than memory. \\
like / such as 例如 & Cultural activities such as reading and travelling grew. \\
a case in point is\ldots 一个典型的例子是…… & A case in point is the boom in live commerce. \\
in particular 尤其 & Health spending, in particular, rose sharply. \\
including\ldots 包括…… & Several activities, including reading, remained stable. \\
to put it simply / roughly stated 简而言之 / 粗略地说 & To put it simply, the chart tells a story of growth.【注】原文为 put it simply / stated roughly，均为不完整结构，此处补全 \\
as an illustration 作为例证 & As an illustration, consider the rise of online payment. \\
A good example would be\ldots 一个很好的例子是…… & A good example would be the spread of mobile payment. \\
To detail this, I would like to\ldots 为详细说明这一点，我想…… & To detail this, I would like to mention the rural share. \\
It is interesting to note that\ldots 值得注意的是…… & It is interesting to note that the gap has narrowed. \\
as proof 作为证据 & As proof, the figure rose by 12 percentage points. \\
take the case of / take\ldots\ as an example 以……为例 & Take the case of online learning; it grew fastest during the pandemic. \\
as for / as regards 至于 & As for the elderly, their leisure time nearly doubled. \\
according to 根据 & According to the chart, the figure stood at 55\%. \\
on this occasion 这一次 & On this occasion, the two groups behaved quite differently. \\
\bottomrule
\end{longtable}

\par\addvspace{6pt}\noindent\textbf{9. 强调关系（Emphasis）}\par\addvspace{2pt}
\begin{longtable}{L{5.4cm}L{10.6cm}}
\toprule
\textbf{衔接词} & \textbf{英语二示范（短语或半句）} \\
\midrule
\endhead
in fact / in truth 事实上 & In fact, the gap has widened since 2015. \\
especially / particularly 尤其 & Particularly in rural areas, the change is visible. \\
moreover 而且 & Moreover, the trend has lasted for a decade. \\
naturally 自然 / 当然 & Naturally, income growth drives such spending. \\
what is particular 值得注意的是 & What is particular about the chart is that the oldest group spends the least.【注】原文 what is particular 单独使用不合语法，必须接 is that 从句或补出 about\ldots\ is that \\
not to mention\ldots 更不用说…… & Not to mention the pressure it puts on public services. \\
believe it or not 信不信由你（口语） & Believe it or not, the figure has tripled in ten years. \\
undeniably / undoubtedly / without a doubt 不可否认 / 毫无疑问 & Undoubtedly, technology has changed the way we learn. \\
it is certain / sure that\ldots 可以肯定…… & It is certain that the trend will continue. \\
by definition 从定义上说 & By definition, lifelong learning never ends. \\
definitely 肯定（口语） & The policy has definitely made a difference. \\
obviously 显然 & Obviously, the chart reflects a change in habits. \\
not only\ldots 不仅…… & Not only the young but also the elderly use mobile payment. \\
in any event 无论如何 & In any event, the overall trend is upward. \\
without reservation 毫无保留地 & We can say without reservation that living standards have improved. \\
\bottomrule
\end{longtable}

\par\addvspace{6pt}\noindent\textbf{10. 条件关系（Condition）}\par\addvspace{2pt}
\begin{longtable}{L{5.4cm}L{10.6cm}}
\toprule
\textbf{衔接词} & \textbf{英语二示范（短语或半句）} \\
\midrule
\endhead
if / unless / lest 如果 / 除非 / 以免 & If the trend continues, the gap will widen. \\
provided that 只要（书面） & Provided that incomes keep rising, spending will follow. \\
if it is the case 如果是这样的话 & If it is the case, the policy should be adjusted. \\
in this sense 在这个意义上 & In this sense, the chart reflects social progress. \\
once\ldots 一旦…… & Once people get used to online payment, they rarely go back. \\
if possible / if necessary 如果可能 / 必要时 & If necessary, the government should offer subsidies. \\
if so / if not 如果是这样 / 如果不是 & If so, the trend is likely to continue. \\
if not all / if anything 即使不是全部 / 要说有的话 & If anything, the gap has grown wider. \\
\bottomrule
\end{longtable}

\par\addvspace{6pt}\noindent\textbf{11. 归纳总结（Summary / Conclusion）}\par\addvspace{2pt}
\begin{longtable}{L{5.4cm}L{10.6cm}}
\toprule
\textbf{衔接词} & \textbf{英语二示范（短语或半句）} \\
\midrule
\endhead
in other words 换言之 & In other words, the two groups are moving apart. \\
to put it in a nutshell 简而言之 & To put it in a nutshell, the chart shows steady progress. \\
in sum / in short / in brief / in summary 总之 / 简言之 & In short, technology has reshaped daily life. \\
therefore / hence 因此 & Therefore, the trend deserves our attention. \\
to sum up / to conclude / in conclusion 总之 / 最后 & To sum up, the change is both inevitable and beneficial. \\
the conclusion can be drawn that\ldots 可以得出这样的结论…… & The conclusion can be drawn that living standards have improved. \\
\bottomrule
\end{longtable}

\tip{图表作文最后一段的段首词，御风明确推荐用 \textbf{Overall}，而不要用 In conclusion（见源文件第 725 行）。上面这一类的词留给小作文和正文中的句内总结。}

\par\addvspace{6pt}\noindent\textbf{12. 方位 / 空间关系（Spatial / Location）}\par\addvspace{2pt}
\begin{longtable}{L{5.4cm}L{10.6cm}}
\toprule
\textbf{衔接词} & \textbf{英语二示范（短语或半句）} \\
\midrule
\endhead
beyond 在……之外 & Beyond the city, the figure is much lower. \\
opposite to 在……对面 & Opposite to the hospital stands a new library. \\
adjacent to 邻近 & The park adjacent to the community is always crowded. \\
at the same place / there / here and there 在同一地点 / 那里 / 到处 & Here and there, shared bikes can be seen. \\
over / above / below 在……上方 / 下方 & The figure for group A is far above that for group B. \\
in the middle / between / around 在中间 / 之间 / 周围 & In the middle of the chart is the largest slice. \\
in front of 在……前面 & In front of the building lies a small square. \\
in the distance / farther 在远处 / 更远 & In the distance, new buildings are under construction. \\
at the right / on this side 在右边 / 这一侧 & At the right of the chart is the smallest share. \\
\bottomrule
\end{longtable}

\tip{方位词在图表作文里只在两种场合有用：描述饼图、柱图的相对位置（“最大的一块在中间”），以及小作文的通知、指路、场地介绍。不要为了用方位词而生造句子。}

\par\addvspace{6pt}\noindent\textbf{13. 目的关系（Purpose）}\par\addvspace{2pt}
\begin{longtable}{L{5.4cm}L{10.6cm}}
\toprule
\textbf{衔接词} & \textbf{英语二示范（短语或半句）} \\
\midrule
\endhead
with this object / for this purpose 为此目的 & For this purpose, the government set up special funds. \\
in order that / so that 为了 / 以便 & So that more students could benefit, the policy was extended. \\
in this way 这样 & In this way, resources can be used more efficiently. \\
since 既然 & Since time is limited, we should focus on key data.【注】多表原因，表目的时较少见 \\
on that account 因此 & On that account, the plan was delayed.【注】多表结果或原因 \\
in case 以防 & In case of emergency, the school has a clear plan. \\
with a view to 为了（后接动名词） & With a view to improving public health, the city built more parks. \\
for the same reason 出于同样的原因 & For the same reason, cultural spending also rose.【注】偏向原因 \\
\bottomrule
\end{longtable}

\par\addvspace{6pt}\noindent\textbf{14. 重申关系（Restatement / Clarification）}\par\addvspace{2pt}
\begin{longtable}{L{5.4cm}L{10.6cm}}
\toprule
\textbf{衔接词} & \textbf{英语二示范（短语或半句）} \\
\midrule
\endhead
in other words 换言之 & In other words, the rich are getting richer. \\
that is to say 也就是说 & That is to say, the gap has not narrowed. \\
namely 即 & Two groups, namely students and workers, chose differently. \\
to put it in another way 换一种说法 & To put it in another way, time is the scarcest resource. \\
as I have said 正如我所说 & As I have said, the trend is not accidental. \\
again / once again 再次 & Once again, technology proves to be a double-edged sword. \\
\bottomrule
\end{longtable}

\par\addvspace{6pt}\noindent\textbf{15. 时间关系（Time）}\par\addvspace{2pt}
\begin{longtable}{L{5.4cm}L{10.6cm}}
\toprule
\textbf{衔接词} & \textbf{英语二示范（短语或半句）} \\
\midrule
\endhead
at once / immediately 立刻 & The figure fell immediately after the policy. \\
at length 最终 / 详尽地（书面） & At length, the two lines met. \\
in the meantime / meanwhile / at the same time 同时 & In the meantime, the rural share also rose. \\
in the end / at last / finally 最终 & In the end, the figure reached 72\%. \\
then / soon / later / next point 然后 / 不久 / 稍后 & Soon after, the number began to climb. \\
not long after / a few minutes later 不久之后 & Not long after, the trend reversed. \\
at present / from this time on 目前 / 从此 & At present, the figure stands at 55\%. \\
all of a sudden 突然 & All of a sudden, the line dropped sharply. \\
from time to time / occasionally 有时 / 偶尔 & Occasionally, the figure dipped below 20\%. \\
since then / formerly / some time ago 从那时起 / 以前 & Since then, the number has doubled. \\
when / whenever / as / once 当……时 & When the policy was introduced, the figure jumped. \\
in a moment / shortly / whereupon 立刻 / 不久 / 于是（书面） & Shortly afterwards, the growth slowed. \\
\bottomrule
\end{longtable}

\impr{15 类逻辑词已全部收入（共 160 余条），并对以下条目做了处理：（1）合并了原文中重复出现的衔接词（next、finally、in other words、accordingly 等跨类重复出现，已在首次出现处给出示范，后续保留词条并标注所属关系）；（2）纠正了 put it simply / stated roughly 的不完整结构，改为 to put it simply / roughly stated；（3）把原文自注“使用条件苛刻”或归类宽松的 independent of、reckless of、by way of、what is particular、since（表目的）、on that account、for the same reason、much more interesting 全部保留词条并在表内以【注】标出限制；（4）为每条衔接词配了可直接用于英语二作文的示范半句，示范句统一使用图表作文语境。}

\subsection{十种高分句式（GRA 加分项）}

\rubric{评分原则明确：“不以生僻词、复杂句或连接词的数量决定分数。”下面十类句式的价值在于\textbf{让准确的意思更准确地说出来}，而不是堆砌。每类都给了“英语二改写版”——因为御风的原始例句大多出自英语一图画题或通用话题，直接搬进图表作文会跑题。}

\par\addvspace{6pt}\noindent\textbf{一、否定句（Negative Sentences）}\par\addvspace{2pt}
\begin{eg}
1. \en{Instead of indulging in playing computer games, children should be taught how to benefit from useful information on the Internet.}\\
应该教孩子们如何从互联网获取有益的信息，而不是沉溺于玩电脑游戏。\\
2. \en{On no account (Under no circumstances) can we ignore the immense value of knowledge.}\\
我们绝不能忽视知识的巨大价值。\\
3. \en{College students take part-time jobs not for more money but for a better understanding of society.}\\
大学生参加兼职工作不是为了赚更多钱，而是为了更好地了解社会。\\
4. \en{One's salary does not depend so much on his educational background as on his ability and contribution to society.}\\
一个人的工资与其说取决于他的教育背景，倒不如说取决于他的能力和对社会的贡献。\\
5. \en{In terms of nutrition, fast food is far from satisfactory.}\\
从营养角度来说，快餐远非令人满意。\\
6. \en{Parents would not expect their children to become useful persons without working hard.}\\
父母们不能指望孩子们不经过刻苦努力就可以成才。\\
\end{eg}
\begin{note}
\textbf{英语二可直接用的改写}\\
1. \en{Instead of spending their spare time on short videos, college students should be encouraged to read more widely.}\\
2. \en{On no account should we turn a blind eye to the digital divide revealed by the chart.}
\end{note}

\par\addvspace{6pt}\noindent\textbf{二、非限制性定语从句（Non-restrictive Relative Clauses）}\par\addvspace{2pt}
\begin{eg}
1. \en{Undoubtedly, practical courses can be applied to reality, which is of vital importance to their future development.}\\
毫无疑问，实用性课程可以用于实际，这对他们未来的发展是非常重要的。\\
2. \en{Children tend to imitate what they have seen and heard on mass media, which is sometimes dangerous and harmful.}\\
孩子们倾向于模仿大众媒体上的所见所闻，这在有时是危险和有害的。\\
3. \en{The majority of students believe that part-time jobs will provide them with more opportunities to develop their interpersonal skills, which may put them in a favorable position in the future job market.}\\
大部分学生相信业余工作会使他们有更多机会发展人际交往能力，这对他们未来找工作非常有利。\\
\end{eg}
\begin{note}
\textbf{英语二可直接用的改写}\\
1. \en{The proportion of online learners rose from 25\% to 58\%, which indicates the growing popularity of digital education.}\\
2. \en{Only about 35\% of the elderly take part in activities requiring much cognitive effort, which suggests that spiritual care may be falling behind.}
\end{note}

\par\addvspace{6pt}\noindent\textbf{三、让步句（Concessive Clauses）}\par\addvspace{2pt}
\begin{eg}
1. \en{This view is widely acknowledged; however, there is little evidence that smoking is beneficial to people's health.}\\
这个观点被广泛认可，然而几乎没有证据表明吸烟对人们健康有利。\\
2. \en{Although (Whilst) the computer has been widely used in class, it cannot replace the role of teachers.}\\
尽管计算机已经广泛用于课堂，但是它不能取代教师的作用。\\
3. \en{Reasonable as the opinion sounds, it cannot bear much analysis.}\\
虽然这个观点听起来有道理，但是它经不住分析。\\
4. \en{In spite of the many conveniences that cars bring to people's life, they can create a series of serious problems.}\\
尽管汽车给人们的生活带来了许多便利，但是也产生了一系列严重的问题。\\
\end{eg}
\begin{note}
\textbf{英语二可直接用的改写}\\
1. \en{Though online learning has grown rapidly, it cannot fully replace classroom teaching.}\\
2. \en{Succinct as the chart is, it conveys abundant insights.}
\end{note}

\par\addvspace{6pt}\noindent\textbf{四、It 引导的句子（It-patterns）}\par\addvspace{2pt}
\begin{eg}
1. \en{It is hard to imagine what our life would look like without computers in modern society.}\\
很难想象如果现代社会没有了计算机，我们的生活会变成什么样子。\\
2. \en{It is conceivable that being physically active does good to health.}\\
可想而知，积极参加体育活动有利于身体健康。\\
3. \en{It is a highly controversial issue whether women should join the armed forces or not.}\\
女性是否应该参军是一个非常有争议的问题。\\
4. \en{The Internet has made it easier for modern people to communicate with each other in a few seconds.}\\
现代人很容易通过互联网在几秒钟内就取得联系。\\
5. \en{It is worth caring about the way a child behaves.}\\
关心孩子们的行为方式是值得的。\\
6. \en{It has been a few decades since the computer came into being.}\\
自从计算机问世以来已经有几十年的时间了。\\
\end{eg}
\begin{note}
\textbf{英语二可直接用的改写}\\
1. \en{It is hard to imagine what our daily life would look like without mobile payment.}\\
2. \en{It is worth noting that the gap between urban and rural users has narrowed only slightly.}
\end{note}

\par\addvspace{6pt}\noindent\textbf{五、假设句（Conditional Sentences）}\par\addvspace{2pt}
\begin{eg}
1. \en{If we destroy old buildings, then we will ruin the traditional culture and heritage.}\\
如果我们推倒老房子，就要破坏传统的文化和遗产。\\
2. \en{You would miss the chance to interact with other students if you used the Internet at home.}\\
假如你在家上网，你就会失去与其他同学交流的机会。\\
3. \en{Once you change your present job, you will be faced with the danger of being unemployed.}\\
一旦你变换了现在的工作，就面临着失业的危险。\\
\end{eg}
\begin{note}
\textbf{英语二可直接用的改写}\\
1. \en{If the current trend continues, the gap between the two groups will widen further.}\\
2. \en{Once online services become cheaper, more elderly people will join in.}
\end{note}

\par\addvspace{6pt}\noindent\textbf{六、倒装句（Inversion）}\par\addvspace{2pt}
\begin{eg}
1. \en{Only in this way can the problem be tackled properly.}\\
只有这样才能妥善地解决这个问题。\\
2. \en{Only when children make arduous efforts can they become successful.}\\
只有当孩子们付出艰苦的努力，他们才能获得成功。\\
3. \en{Not only does studying in school serve academic purposes, but it also teaches students how to handle interpersonal relations.}\\
在校学习不仅为了学术目的，而且还可以学会如何处理人际关系。\\
4. \en{Under no circumstances should youngsters follow negative information on mass media blindly.}\\
青少年决不应当盲从大众媒体上的不良信息。\\
5. \en{So valuable is time that we cannot afford to waste it.}\\
时间是如此珍贵，我们经不起浪费它。\\
6. \en{Nowadays, the most dangerous tendency for youngsters is to indulge in playing PC games.}\\
如今对青少年最为危险的事情是倾向于过多地玩电脑游戏。\\
\end{eg}
\begin{note}
\textbf{英语二可直接用的改写}\\
1. \en{Only in this way can we narrow the gap revealed by the chart.}\\
2. \en{So striking is the difference that it deserves our close attention.}
\end{note}

\par\addvspace{6pt}\noindent\textbf{七、强调句（Emphatic Sentences）}\par\addvspace{2pt}
\begin{eg}
1. \en{It is for the benefit of maintaining the ecological balance that human beings ought to protect the endangered animal species.}\\
正是为了维持生态平衡，人类才应该保护地球上濒危的动物物种。\\
2. \en{It is not until people suffer from fatal diseases such as SARS and AIDS that they become aware of how significant it is for the government to invest more money in medical care.}\\
直到人们患了像非典和艾滋病这样致命的疾病时，才意识到政府把更多的钱投入到医疗上的重要性。\\
\end{eg}
\begin{note}
\textbf{英语二可直接用的改写}\\
1. \en{It is the improvement of living standards that has made such spending possible.}\\
2. \en{It is not until the gap becomes obvious that people begin to pay attention to it.}
\end{note}

\par\addvspace{6pt}\noindent\textbf{八、比较句（Comparative Sentences）}\par\addvspace{2pt}
\begin{eg}
1. \en{Studies show that juvenile delinquency rates are twice as high for youngsters from single-parent families as for those in traditional households.}（同级比较）\\
研究表明，来自单亲家庭的青少年犯罪率是来自传统家庭的两倍。\\
2. \en{Generally speaking, people in modern times enjoy less leisure time than they did previously.}（比较级）\\
一般说来，现代人比过去享有更少的休闲时间。\\
3. \en{Compared with children of average abilities, intellectually gifted children do better in their school work and intelligence tests.}\\
与能力不同的孩子相比，聪明的孩子在学业和智力测验方面表现更加突出。\\
4. \en{Unlike pop stars, other professionals such as doctors and teachers create true value for society.}\\
与明星不同，其他专业人员如医生和教师为社会创造了真正的价值。\\
\end{eg}
\begin{note}
\textbf{英语二可直接用的改写}\\
1. \en{The number of online learners in 2024 was twice as high as that recorded in 2015.}\\
2. \en{People in big cities spend far more on culture than those in small towns do.}
\end{note}

\par\addvspace{6pt}\noindent\textbf{九、插入语（Parenthesis）}\par\addvspace{2pt}
\begin{eg}
1. \en{Computers, most important of all, make worldwide communication possible.}\\
最重要的是，计算机使世界范围内的交流成为可能。\\
2. \en{Providing more parking areas, in the long run, has proved to be a practical way out in many large cities in the world.}\\
从长远的利益来看，提供更多的停车位证明是许多世界大城市切实可行的出路。\\
3. \en{Observing local culture, consciously and unconsciously, can reduce the chance of offending the locals, or it will lead to embarrassment and even conflict.}\\
无论是否意识到，遵循当地的文化可以减少冒犯当地人的机会，否则就会导致尴尬甚至冲突。\\
4. \en{Space exploration, some people believe, will bring some unexpected discoveries.}\\
有人相信空间探索将会带来意想不到的发现。\\
\end{eg}
\begin{note}
\textbf{英语二可直接用的改写}\\
1. \en{Online learning, most important of all, has made education accessible to more people.}\\
2. \en{The gap, in the long run, will narrow as incomes rise.}
\end{note}

\par\addvspace{6pt}\noindent\textbf{十、原因句（Sentences of Cause）}\par\addvspace{2pt}
\begin{eg}
1. \en{Violent films can do harm to young people, because they contain a great deal of negative information.}\\
暴力影片对年轻人是有害的，因为其中包含大量负面信息。\\
2. \en{Human beings are chiefly responsible for wild animal extinction.}\\
野生动物灭绝的主要责任在人类。\\
3. \en{We should attribute (ascribe) medical advances to the animals which make substantial contributions to experiments.}\\
我们应该把医学的发展归因于动物对实验做出的巨大贡献。\\
4. \en{Because of overpopulation, water shortage has become one of the most serious problems.}\\
由于人口过剩，水资源短缺已经成为一个最严重的问题。\\
5. \en{Owing to the fierce competition in today's world, a great many young people find themselves under great pressure.}\\
由于当今世界的激烈竞争，大量的年轻人承受巨大的压力。\\
\end{eg}
\begin{note}
\textbf{英语二可直接用的改写}\\
1. \en{The change is chiefly due to the improvement of living standards.}\\
2. \en{Owing to the rapid spread of the Internet, online services have become part of daily life.}
\end{note}

\impr{十种句式共 43 条例句已全部收入，并在原文基础上做了以下校正：（1）“a better understanding of societies”→“society”、“contribution to the society”→“society”（society 泛指“社会”时不加冠词、不用复数）；（2）“practical courses can be used to the reality”→“can be applied to reality”；（3）“serve academic purpose”→“serve academic purposes”；（4）“do more well in”→“do better in”；“Compared with those with different abilities”→“Compared with children of average abilities”，使比较双方互不包含；“children with intelligence”→“intellectually gifted children”；（5）“because they contains numerous negative information”→“because they contain a great deal of negative information”（information 不可数）；（6）“do substantial contributions to the experiment”→“make substantial contributions to experiments”（贡献搭配 make a contribution）；（7）“they are becoming aware”→“they become aware”（It is not until\ldots that 从句用一般现在时）；（8）原文中的全角分号、逗号（；，）改为英文标点；（9）“In spite of\ldots it can create”的主语由 it 改为 they，指代 cars；（10）“Computers\ldots create wide communication”改为“make worldwide communication possible”。所有中文释义保留原样。}

\subsection{评述“利弊 / 影响”的短语}

\begin{note}
第二段分析利弊时最缺的不是观点，而是动词。下面这些短语的价值在于：\textbf{一个词就能领起一个分句}，且比 good / bad / helpful 精确得多。注意 fuel、power、feed、drive、boost、bolster 都是及物动词，后面直接接宾语。
\end{note}

\par\addvspace{5pt}\noindent\textbf{一、“有利”“促进”之意（11 条）}\par\addvspace{2pt}
\begin{longtable}{L{4.8cm}L{2.8cm}L{7.8cm}}
\toprule
\textbf{短语} & \textbf{中文} & \textbf{英语二场景例句} \\
\midrule
\endhead
hold\ldots\ in good stead & 使……处于有利地位 & Solid language skills will hold you in good stead in the job market. \\
pave the way for\ldots & 为……铺平道路 & Online payment has paved the way for the rapid growth of e-commerce. \\
be remarkably beneficial to\ldots & 对……很有好处 & Regular exercise is remarkably beneficial to both body and mind. \\
serve as a stepping stone for\ldots & 可作为……的垫脚石 & Internships serve as a stepping stone for graduates entering the job market. \\
be instrumental in\ldots & 在……方面起重要作用 & Government support has been instrumental in narrowing the digital divide. \\
fuel\ldots & 助推（多接消极或中性的势头） & The growing demand for private cars fuels traffic congestion. \\
power\ldots & 为……提供动力 & Digital technology powers the growth of the sharing economy. \\
feed\ldots & 助长；滋养 & Endless advertisements feed consumers' desire for new products. \\
drive\ldots & 推动 & Rising incomes have driven the growth of cultural spending. \\
boost\ldots & 促进；提升 & Online platforms boost the efficiency of small businesses. \\
bolster\ldots & 增强；巩固 & Lifelong learning bolsters one's competitiveness in the job market. \\
\bottomrule
\end{longtable}

\par\addvspace{5pt}\noindent\textbf{二、“有害”“阻碍”之意（10 条）}\par\addvspace{2pt}
\begin{longtable}{L{4.8cm}L{2.8cm}L{7.8cm}}
\toprule
\textbf{短语} & \textbf{中文} & \textbf{英语二场景例句} \\
\midrule
\endhead
deliver a blow to\ldots & 打击…… & The pandemic delivered a heavy blow to the tourism industry. \\
cast / throw a cloud over\ldots & 使……蒙上阴影 & Rising costs have cast a cloud over small businesses. \\
take a toll on\ldots & 对……造成严重伤害 & Long working hours take a heavy toll on people's health. \\
exert negative effects on\ldots & 对……有消极影响 & Excessive screen time exerts negative effects on children's eyesight. \\
serve as a barrier to\ldots & 是……的妨碍 & High costs serve as a barrier to lifelong learning. \\
serve as a stumbling block for\ldots & 是……的绊脚石 & A lack of digital skills serves as a stumbling block for older job seekers. \\
get in the way of\ldots & 阻挡……的道路 & Short-term profits often get in the way of long-term development. \\
be detrimental to\ldots & 对……有害 & Overwork is detrimental to both health and efficiency. \\
impede\ldots & 阻碍 & Poor infrastructure impedes the growth of rural e-commerce. \\
hinder\ldots & 妨碍 & The digital divide hinders equal access to education. \\
\bottomrule
\end{longtable}

\begin{note}
\textbf{使用分寸}\\
（1）deliver a blow to、take a toll on、cast a cloud over 语气很重，适合写“危害”类第二段，不要用来描述轻微的不便。\\
（2）fuel 后面多接消极名词（congestion、inflation、anxiety），接积极内容会不自然。\\
（3）impede / hinder 是正式书面动词，比 stop 准确：stop 是“使停止”，impede 是“使变慢、变难”。\\
（4）“stumbling stone” 是错误搭配，正确习语为 \kw{stumbling block}（绊脚石），已在表中更正。
\end{note}

\impr{（1）“有利 / 促进”类原文实为 11 条（编号 1--11），已全部收入，未按任务的“10 条”截断；（2）“serve as a stumbling stone for”更正为“serve as a stumbling block for”；（3）为原文只给中文的 6 个光杆动词（fuel / power / feed / drive / boost / bolster）补出英语二场景例句与使用分寸；（4）原文第 6、7 条的编号断行造成的排版错位（fuel、power 与 drive、boost、bolster 之间的页码插入）已按语义复原。}

\subsection{趋势与程度词汇}

\begin{note}
御风讲得很直接：“有两块是一定可以用得上的，也就是趋势的描述和程度的描述，可以多换换词。”换词的前提是\textbf{分得清轻重}：soar 和 rise 不是同义词，前者自带“猛”的意思；marginally 和 dramatically 更不能互换。下面按“方向 + 力度”分表，每条都给出图表作文里可以直接抄的片段。
\end{note}

\par\addvspace{5pt}\noindent\textbf{一、上升趋势（温和：climb / increase / rise / grow）}\par\addvspace{2pt}
\begin{longtable}{L{2.9cm}L{1.9cm}L{10.9cm}}
\toprule
\textbf{词} & \textbf{中文} & \textbf{图表作文示范片段} \\
\midrule
\endhead
climb & 攀升 & The number of online learners climbed steadily from 40 million to 65 million. \\
increase & 增加 & Investment in renewable energy increased by 15\% over the decade. \\
rise & 上升 & The proportion of rural households with Internet access rose from 30\% to 70\%. \\
grow & 增长 & Cultural consumption grew at an average annual rate of 9\%. \\
\bottomrule
\end{longtable}

\par\addvspace{5pt}\noindent\textbf{二、上升趋势（剧烈：spiral / soar / rocket / surge / shoot up / leap）}\par\addvspace{2pt}
\begin{longtable}{L{2.9cm}L{1.9cm}L{10.9cm}}
\toprule
\textbf{词} & \textbf{中文} & \textbf{图表作文示范片段} \\
\midrule
\endhead
spiral & 急剧攀升（多含失控义） & Housing prices spiralled in the largest cities. \\
soar & 猛增 & The number of new energy vehicles soared from 0.2 million to 3 million. \\
rocket & 飞涨（幅度极大，偏口语） & E-commerce sales rocketed to 1.2 trillion yuan. \\
surge & 激增（图表作文最常用） & Spending on education surged by 12 percentage points. \\
shoot up & 暴涨 & The figure shot up after 2015. \\
leap & 跃升（跨过一个大台阶） & Online sales leapt from 5\% to 25\% of total retail sales. \\
\bottomrule
\end{longtable}

\par\addvspace{5pt}\noindent\textbf{三、下降趋势（温和：dip / slide / fall / decline / drop / decrease / diminish）}\par\addvspace{2pt}
\begin{longtable}{L{2.9cm}L{1.9cm}L{10.9cm}}
\toprule
\textbf{词} & \textbf{中文} & \textbf{图表作文示范片段} \\
\midrule
\endhead
dip & 小幅下降（也可作名词 a dip） & The share dipped slightly to 18\%. \\
slide & 下滑（常含“不知不觉地”） & Attendance slid from 60\% to 45\% over five years. \\
fall & 下降（最常用） & The birth rate fell by 2 per 1,000. \\
decline & 下降（较正式） & The number of smokers declined gradually. \\
drop & 下降（常用 drop to） & The figure dropped to its lowest point in a decade. \\
decrease & 减少（及物、不及物皆可） & Average working hours decreased modestly. \\
diminish & 逐渐减小（多用于抽象事物） & The role of cash diminished year by year. \\
\bottomrule
\end{longtable}

\par\addvspace{5pt}\noindent\textbf{四、下降趋势（剧烈：plumb / plunge / plummet / nosedive / tumble / slump）}\par\addvspace{2pt}
\begin{longtable}{L{2.9cm}L{1.9cm}L{10.9cm}}
\toprule
\textbf{词} & \textbf{中文} & \textbf{图表作文示范片段} \\
\midrule
\endhead
plumb & 跌至（多接低点） & The index plumbed a ten-year low. \\
plunge & 暴跌 & The stock index plunged by 30\% in a single month. \\
plummet & 骤降 & Sales plummeted from 2 million to 0.5 million. \\
nosedive & 急剧下滑 & Profits nosedived in the second quarter. \\
tumble & 大幅下跌 & The price tumbled to a record low. \\
slump & 暴跌；陷入低谷 & Exports slumped by nearly a fifth. \\
\bottomrule
\end{longtable}

\par\addvspace{5pt}\noindent\textbf{五、程度轻微（slightly / modestly / moderately / marginally）}\par\addvspace{2pt}
\begin{longtable}{L{2.9cm}L{1.9cm}L{10.9cm}}
\toprule
\textbf{词} & \textbf{中文} & \textbf{图表作文示范片段} \\
\midrule
\endhead
slightly & 略微 & The figure rose slightly by 3\%. \\
modestly & 小幅地 & Spending grew modestly over the decade. \\
moderately & 适度地 & The number increased moderately. \\
marginally & 微弱地（几乎可忽略） & The share edged up marginally. \\
\bottomrule
\end{longtable}

\par\addvspace{5pt}\noindent\textbf{六、程度显著（10 个）}\par\addvspace{2pt}
\begin{longtable}{L{2.9cm}L{1.9cm}L{10.9cm}}
\toprule
\textbf{词} & \textbf{中文} & \textbf{图表作文示范片段} \\
\midrule
\endhead
considerably & 相当大地 & The number rose considerably after 2018. \\
remarkably & 显著地 & Efficiency improved remarkably. \\
notably & 显著地（也作“尤其”） & Cultural spending was notably higher in cities. \\
noticeably & 明显地 & The gap became noticeably wider. \\
markedly & 明显地（较正式） & The two groups behaved markedly differently. \\
substantially & 大幅度地 & Rural incomes rose substantially. \\
significantly & 显著地（图表题高频） & The share of online shoppers was significantly higher than in 2015. \\
dramatically & 急剧地 & The cost dropped dramatically. \\
radically & 根本性地（强调性质变化） & The way people pay has changed radically. \\
exponentially & 指数级地 & Online data grew exponentially. \\
\bottomrule
\end{longtable}

\par\addvspace{5pt}\noindent\textbf{七、波动、持稳与极值（本次整理新增）}\par\addvspace{2pt}
\begin{longtable}{L{2.9cm}L{1.9cm}L{10.9cm}}
\toprule
\textbf{词} & \textbf{中文} & \textbf{图表作文示范片段} \\
\midrule
\endhead
fluctuate & 波动 & The figure fluctuated between 40\% and 45\%. \\
hover around & 在……附近徘徊 & The figure hovered around 30\% for five years. \\
level off & 趋于平稳（高位之后的持平） & After 2020, the growth leveled off. \\
plateau & 进入平台期 & Demand plateaued at around 60 million. \\
remain stable & 保持稳定 & The proportion remained stable at about 55\%. \\
stabilize & 稳定下来 & Prices stabilized after the policy was introduced. \\
peak & 达到峰值 & It peaked at 90 billion yuan. \\
bottom out & 触底 & The figure bottomed out in 2016 and then recovered. \\
hit a low & 跌至低点 & Sales hit a ten-year low. \\
edge up / edge down & 微升 / 微降 & The share edged up to 21\%. \\
\bottomrule
\end{longtable}

\impr{原文只给了上升、下降、程度轻微、程度显著四组词，但御风自己的模板例句里反复出现“underwent a steady increase”“climbed to a peak”“remained at approximately 55\%”“registered a substantial increase”（源文件第 591--601、692--701 行）。为保证词表能覆盖图表题的全部描述需要，本表补出“波动、持稳与极值”一组，共 10 条，属\textbf{从源材料例句中提炼}，不是外来内容。}

\par\addvspace{5pt}\noindent\textbf{八、极值、聚类与差异（御风 Q1：什么信息最该写）}\par\addvspace{2pt}
\begin{longtable}{L{3.6cm}L{6.6cm}L{5.5cm}}
\toprule
\textbf{功能} & \textbf{可用词} & \textbf{示范} \\
\midrule
\endhead
极值（最重要） & highest, main, major, the most, the lowest & The highest participation rate went to watching TV, at 90.8\%. \\
聚类（相同） & similar, the same, likewise, similarly, also, similar to, at a similar level & Reading and gardening stayed at a similar level, at about 30\%. \\
聚类 / 差异（不同） & in contrast, on the other hand, different from, distinct from, substantial / striking difference between, the gap between & In contrast, activities requiring cognitive effort made up less than 35\%, and the gap between the two groups was striking. \\
\bottomrule
\end{longtable}

\tip{动态图优先写\textbf{趋势与比较}，其次写\textbf{极值}、最新数据、起始数据；静态图没有趋势，优先写\textbf{极值、比较、聚类}。这是“数据写不完”时决定写哪几个数的唯一标准。}

\par\addvspace{5pt}\noindent\textbf{九、倍数与增减幅度}\par\addvspace{2pt}
\begin{longtable}{L{3.6cm}L{6.8cm}L{5.3cm}}
\toprule
\textbf{用法} & \textbf{正确} & \textbf{要避免的写法} \\
\midrule
\endhead
double / triple / quadruple（动词，翻倍 / 增至三倍 / 增至四倍） & The number doubled between 2010 and 2020. Online sales tripled in five years. & The number doubled by two times.（累赘） \\
twice / three times（限定词，两倍 / 三倍） & The figure was twice as high as that recorded in 2015. & The figure was twice higher than that in 2015.（twice 不与比较级连用） \\
twofold / threefold（倍数副词） & Spending rose threefold over the decade. & Spending rose three times.（易被读成“涨了三次”） \\
\ldots times as \ldots\ as（是……的几倍） & A is three times as large as B. & A is three times larger than B.（有歧义：是“三倍”还是“多三倍”） \\
by a factor of（达到……倍）【补】 & Sales increased by a factor of three. & Sales increased by three factors.（错误搭配） \\
\bottomrule
\end{longtable}

\par\addvspace{5pt}\noindent\textbf{十、图表数值万能动词（御风“逐个写法”三件套：关于 + 指标 + 怎么样）}\par\addvspace{2pt}
\begin{longtable}{L{3.2cm}L{6.6cm}L{5.9cm}}
\toprule
\textbf{功能} & \textbf{可用词} & \textbf{示范} \\
\midrule
\endhead
“关于”（引出对象） & concerning, regarding, as for, in the case of, for, with regard to, in terms of & Regarding the proportion of graduates, the figure stood at 35\%, representing the largest share among all categories. \\
“指标”（换掉 percentage） & percentage, share, quantity, amount, proportion, figure, index & In terms of public transport, the percentage of daily commuters remained at approximately 55\%. \\
“怎么样”（换掉 is / has） & account for, make up, constitute, represent, stand at, stay at, reach, hit, register & As for smartphone users, the figure surged to 1.2 billion. With regard to household spending on health care, it reached 8,000 yuan per capita. \\
\bottomrule
\end{longtable}

\impr{第十张表来自御风静态图“逐个写法”的词表（源文件第 663--701 行），原本是描述段素材。它与第 3 章的“关于 / 指标 / 怎么样”三件套互为补充：第 3 章管句式，本表管动词。}

\subsection{高频话题词库（按主题分类）}

\begin{note}
mmg 的词表是“题目相关词”，共 70 条（其中 sense of responsibility 与 responsibility 重复，实为 69 条）。这里按主题重新分组，好处是：\textbf{同一个主题下的词可以互相激活}，考场上想起一个就能想起一串；同时把每条词的中文全部补齐、把 OCR 错误全部改正。带【补】的是本次从御风源文件（第 992--995 行的 filial piety、resist the temptation，以及第 974--979 行的主题方向）中补入的。
\end{note}

\par\addvspace{5pt}\noindent\textbf{A. 科技与数字社会}\par\addvspace{2pt}
\begin{longtable}{L{5.8cm}L{10.2cm}}
\toprule
\textbf{英文} & \textbf{中文} \\
\midrule
\endhead
Mobile Internet & 移动互联网 \\
netizen & 网民 \\
artificial intelligence & 人工智能 \\
electronic product & 电子产品 \\
big data & 大数据 \\
facial recognition & 人脸识别 \\
digital divide & 数字鸿沟 \\
data privacy & 数据隐私 \\
smart city & 智慧城市 \\
technology addiction & 科技依赖（手机 / 网络成瘾） \\
5G technology & 5G 技术 \\
\bottomrule
\end{longtable}

\par\addvspace{5pt}\noindent\textbf{B. 教育与学习}\par\addvspace{2pt}
\begin{longtable}{L{5.8cm}L{10.2cm}}
\toprule
\textbf{英文} & \textbf{中文} \\
\midrule
\endhead
online learning & 在线学习 \\
academic pressure & 学业压力 \\
quality education & 素质教育 \\
exam-oriented education & 应试教育 \\
lifelong learning & 终身学习 \\
peer competition & 同龄人竞争 \\
parental expectation & 家长期望 \\
youth independence & 青年独立性 \\
critical thinking & 批判性思维 \\
study-abroad trend & 留学趋势 \\
\bottomrule
\end{longtable}

\par\addvspace{5pt}\noindent\textbf{C. 就业与职业}\par\addvspace{2pt}
\begin{longtable}{L{5.8cm}L{10.2cm}}
\toprule
\textbf{英文} & \textbf{中文} \\
\midrule
\endhead
seek a job & 找工作 \\
flexible employment & 灵活就业 \\
work-life balance & 工作与生活的平衡 \\
career planning & 职业规划 \\
youth entrepreneurship & 青年创业 \\
\bottomrule
\end{longtable}

\par\addvspace{5pt}\noindent\textbf{D. 社会与人口}\par\addvspace{2pt}
\begin{longtable}{L{5.8cm}L{10.2cm}}
\toprule
\textbf{英文} & \textbf{中文} \\
\midrule
\endhead
citizen & 公民 \\
population aging & 人口老龄化 \\
aging society & 老龄化社会 \\
spoil children & 溺爱孩子 \\
raise children & 抚养孩子 \\
taking care of the elderly parents & 赡养年迈的父母 \\
charity & 慈善 \\
volunteer spirit & 志愿者精神 \\
filial piety【补】 & 孝顺 \\
family bonds【补】 & 亲情 \\
interpersonal relationships【补】 & 人际关系 \\
\bottomrule
\end{longtable}

\par\addvspace{5pt}\noindent\textbf{E. 环境与可持续}\par\addvspace{2pt}
\begin{longtable}{L{5.8cm}L{10.2cm}}
\toprule
\textbf{英文} & \textbf{中文} \\
\midrule
\endhead
environmental awareness & 环保意识 \\
garbage classification & 垃圾分类 \\
public transport & 公共交通 \\
subway & 地铁 \\
biodiversity & 生物多样性 \\
natural reserve & 自然保护区 \\
low-carbon lifestyle & 低碳生活 \\
energy conservation & 节能 \\
food waste & 食物浪费 \\
recycling awareness & 回收意识 \\
climate change & 气候变化 \\
green consumption & 绿色消费 \\
environmental protection & 环境保护 \\
\bottomrule
\end{longtable}

\par\addvspace{5pt}\noindent\textbf{F. 文化与价值}\par\addvspace{2pt}
\begin{longtable}{L{5.8cm}L{10.2cm}}
\toprule
\textbf{英文} & \textbf{中文} \\
\midrule
\endhead
culture and civilization & 文化和文明 \\
moral education & 道德教育 \\
honesty & 诚信 \\
sense of responsibility & 责任感（与 responsibility 合并） \\
devotion & 奉献 \\
national identity & 国家认同 \\
innovation & 创新意识 \\
a positive attitude【补】 & 积极心态 \\
youth values【补】 & 青年价值观 \\
respect for the elderly【补】 & 尊老敬老 \\
\bottomrule
\end{longtable}

\par\addvspace{5pt}\noindent\textbf{G. 健康与生活方式}\par\addvspace{2pt}
\begin{longtable}{L{5.8cm}L{10.2cm}}
\toprule
\textbf{英文} & \textbf{中文} \\
\midrule
\endhead
food security & 食品安全 \\
infectious disease & 传染病 \\
time management & 时间管理 \\
self-discipline & 自律 \\
resilience & 韧性（抗压与恢复能力） \\
resist the temptation【补】 & 抵制诱惑 \\
\bottomrule
\end{longtable}

\par\addvspace{5pt}\noindent\textbf{H. 数字经济、消费与经济}\par\addvspace{2pt}
\begin{longtable}{L{5.8cm}L{10.2cm}}
\toprule
\textbf{英文} & \textbf{中文} \\
\midrule
\endhead
live commerce & 直播带货 \\
online payment & 在线支付 \\
e-commerce & 电子商务 \\
shared economy & 共享经济 \\
digital economy & 数字经济 \\
inflation & 通货膨胀 \\
economic resilience & 经济韧性 \\
economic globalization & 经济全球化 \\
credibility & 信誉 \\
urbanization & 城市化 \\
consumption upgrading【补】 & 消费升级 \\
living standards【补】 & 生活水平 \\
disposable income【补】 & 可支配收入 \\
\bottomrule
\end{longtable}

\impr{（1）OCR 错误已全部更正：environment awareness → \kw{environmental awareness}；lifelonglearning → \kw{lifelong learning}；seek ajob → \kw{seek a job}；（2）sense of responsibility 与 responsibility 同义，合并为一条；（3）taking care of the elderly parents 保留原形，但更地道的写法是 taking care of one's elderly parents / supporting elderly parents；（4）为每个主题补出 3--4 个御风源文件里的方向大词（filial piety、family bonds、interpersonal relationships、a positive attitude、youth values、respect for the elderly、resist the temptation、consumption upgrading、living standards、disposable income），并用【补】标出；（5）以下词跨主题，已在主表收入一次，阅读时可交叉激活：online learning（B，也属 A）、technology addiction（A，也属 G）、work-life balance（C，也属 G）、critical thinking（B，也属 F）、national identity（F，也属 D）、low-carbon lifestyle 与 green consumption（E，也属 H）、youth entrepreneurship（C，也属 H）、urbanization（H，也属 D）、parental expectation 与 peer competition（B，也属 D）。}

\begin{tpl}{「上位概念」表：把题目映射成方向大词}
\begin{longtable}{L{3.4cm}L{5.7cm}L{5.6cm}}
\toprule
\textbf{图表常见主题} & \textbf{方向大词（3--5 个）} & \textbf{可用论证角度} \\
\midrule
\endhead
收入 / 消费 / 城镇化 & economic development, living standards, consumption upgrading, disposable income, urbanization & 收入增长 → 消费结构变化 → 生活质量提升 \\
人口老龄化 & demographic change, social welfare, public services, filial piety, aging society & 老龄化 → 养老负担 → 社会保障与家庭责任 \\
网络 / 手机 / APP & the digital age, information technology, digital literacy, online services, smart life & 技术普及 → 生活便利 → 数字鸿沟 \\
教育 / 读书 / 留学 & human capital, lifelong learning, quality education, personal development, competitiveness & 教育投入 → 人力资本 → 个人与社会的回报 \\
环境 / 垃圾 / 低碳 & sustainable development, ecological awareness, green lifestyle, low-carbon economy, environmental protection & 环保意识 → 行为改变 → 可持续发展 \\
健康 / 运动 / 饮食 & public health, healthy lifestyle, quality of life, medical services, well-being & 健康意识 → 生活方式 → 医疗负担 \\
就业 / 创业 / 兼职 & employment pressure, the job market, career development, entrepreneurship, work-life balance & 就业压力 → 能力提升 → 多元就业 \\
文化 / 传统 / 读书 & cultural confidence, traditional culture, spiritual life, national identity, soft power & 文化需求 → 精神生活 → 文化自信 \\
志愿 / 慈善 / 互助 & social responsibility, civic awareness, public spirit, social harmony & 责任意识 → 社会参与 → 社会和谐 \\
诚信 / 道德 / 规则 & moral standards, social trust, integrity, civic morality & 失信成本 → 社会信任 → 制度与教育 \\
\bottomrule
\end{longtable}
\end{tpl}

\tip{\textbf{为什么背“上位概念”而不是背具体题目词？}御风讲得很清楚：主题划分是不确定的，同一道题可以是科技、人际关系、亲情、网络时代中任何一个方向；而且“泛泛而谈容易跑题，太具体很难压中”。上位概念处在中间层——它比具体题目词（如 facial recognition）通用，又比 good / important 这种空词有内容，正好对应“主题外无效，主题内无敌”的选词标准。\textbf{使用提示}：本表是查表用的，不必背诵。第 6 章另有一张按写作顺序排列的简表（主题 → 上位概念 → 评论方向），写第 2 段时用那张更顺；本表覆盖的主题更全，复习阶段查这份。}

\begin{note}
第 6 章已经给出一张内容相同的「句内功能动词表」（定义 / 归因 / 影响 / 对比假设 / 举例 / 评价六类，并配英语二示范句）。为避免同一张表出现两次，此处不再重复，需要造句时直接翻到第 6 章。
\end{note}



\subsection{数字表达速查（by / at / with / to、比较、倍数、时间）}

\begin{note}
这一小节对应御风“Q3：基础非常差的同学，再补一点基本知识”。它的定位不是“高级”，而是\textbf{不丢分}：图表作文里最容易被判的语法错误，几乎全部集中在这四类数字表达上。每条都给“正确 / 典型错误”对照。
\end{note}

\zone{先分清“百分比”和“百分点”。\kw{rise by 5\%} 指在原有基础上涨了 5\%（相对变化）；\kw{rise by 5 percentage points} 指从一个百分数变成另一个百分数，差 5 个点（绝对变化）。图表里比较两个比例时，必须用 percentage points。另外，数字后面写 \kw{percent}（不加冠词），“百分比”这个名词才写 \kw{percentage}。误：The share rose from 20\% to 25\%, an increase of 5\%.（应为 5 percentage points）}

\par\addvspace{5pt}\noindent\textbf{一、描述数字：By / At / With / To}\par\addvspace{2pt}
\begin{longtable}{L{4.0cm}L{6.2cm}L{5.5cm}}
\toprule
\textbf{规则} & \textbf{正确} & \textbf{典型错误} \\
\midrule
\endhead
By + 幅度（变化了多少） & The proportion gradually declines by 10\% every decade. & The proportion gradually declines 10\% every decade.（漏掉介词，考试中视为语法错误） \\
At + 具体数字（达到某个点） & It peaked at 90 billion yuan. & It peaked 90 billion yuan.（peak 是不及物动词，必须接 at） \\
With + 具体数字（作补充说明；主语不能是 number、figure 这类表数字的名词） & Sales reached 2 billion yuan, with online channels accounting for 40\%. & The number with online channels accounting for 40\% \ldots（主语是 number，不能用 with 引数字） \\
To + 变化后的终点（常与 from 呼应） & The figure rose from 20 million to 35 million. & The figure rose to 20 million from 35 million.（方向颠倒，读者会以为起点是 35 million） \\
\bottomrule
\end{longtable}

\par\addvspace{5pt}\noindent\textbf{二、比较数字}\par\addvspace{2pt}
\begin{longtable}{L{2.4cm}L{5.9cm}L{7.4cm}}
\toprule
\textbf{含义} & \textbf{表达} & \textbf{正确 / 要避免} \\
\midrule
\endhead
多达 & up to / as much as & 正：Up to 60\% of the respondents preferred online courses. 注意：as much as 接不可数或整体量，可数名词复数用 as many as。 \\
仅仅 & only / a mere / no more than & 正：A mere 5\% of the elderly took part in such activities. 误：only a mere 5\%（三者意义重复，不要叠用） \\
超过 & top / exceed / in excess of & 正：The number exceeded 10 million. / Sales topped 2 billion yuan. 误：exceeded than 10 million \\
少于 & less than / below & 正：Fewer than 200 students chose that option. 注意：可数名词用 fewer，less 接不可数；below 多接水平、标准。 \\
多于 & more than / above / over & 正：More than half of the users were under 30. \\
剩余的 & remaining & 正：The remaining 30\% chose to study abroad. 误：The remained 30\%（remaining 是形容词，remained 是动词过去式） \\
\bottomrule
\end{longtable}

\par\addvspace{5pt}\noindent\textbf{三、倍数：三种写法互换}\par\addvspace{2pt}
\begin{longtable}{L{2.1cm}L{3.1cm}L{4.1cm}L{5.9cm}}
\toprule
\textbf{汉语} & \textbf{动词式} & \textbf{限定词式} & \textbf{倍数副词 / 百分比式} \\
\midrule
\endhead
翻倍（增至两倍） & The number doubled. & twice as high as that in 2015 & twofold / a 100\% increase \\
增至三倍 & The number tripled. & three times as high as that in 2015 & threefold / a 200\% increase \\
增至四倍 & The number quadrupled. & four times as high as that in 2015 & fourfold / a 300\% increase \\
\bottomrule
\end{longtable}

\tip{“增长了三倍”在汉语里常指“涨到四倍”，英语必须写清楚：increase by 200\%（净增两倍）还是 increase to three times（涨到三倍）。倍数句式的正误辨析（three times larger than 的歧义、“increased three times”的误解）见上一小节第九表。}

\par\addvspace{5pt}\noindent\textbf{四、时间}\par\addvspace{2pt}
\begin{longtable}{L{2.6cm}L{6.0cm}L{7.0cm}}
\toprule
\textbf{用法} & \textbf{正确} & \textbf{要避免} \\
\midrule
\endhead
期间 & during / over / between \ldots\ and / from \ldots\ to & Over the past decade, the figure rose steadily. 误：Over the past decade, the figure rises steadily.（时间已结束，用过去时） \\
时间节点前 & earlier / by / until & By 2020, the number had reached 80 million. 误：By 2020, the number has reached 80 million.（by + 过去时间用过去完成） \\
时间节点后 & from \ldots\ onwards / in the following \ldots\ / in a row / for consecutive years & From 2015 onwards, the trend reversed. The figure rose for three consecutive years. 误：for three consecutively years（形容词 consecutive 修饰名词 years） \\
区间搭配 & between \ldots\ and \ldots & Between 2010 and 2020, the number doubled. 误：between 2010 to 2020 \\
\bottomrule
\end{longtable}

\par\addvspace{5pt}\noindent\textbf{五、年龄表达（半动态图专用）}\par\addvspace{2pt}
\begin{longtable}{L{3.0cm}L{6.4cm}L{6.2cm}}
\toprule
\textbf{用法} & \textbf{表达} & \textbf{示范} \\
\midrule
\endhead
People aged + 年龄段 & people aged under 24 / between 30 and 39 / over 50 & People aged over 50 spend the most on health care. \\
in + 年龄层 & people in their twenties / in their fifties & People in their thirties are the most likely to change jobs. \\
Age group & the under-24 age group / the 30--39 age group & The under-24 age group accounts for the largest share. \\
复合形容词 & under-24-year-olds & Under-24-year-olds make up nearly half of all users. \\
\bottomrule
\end{longtable}

\begin{eg}
动态视角一：\en{With age, people are more likely to study out of interest.}\\
动态视角二：\en{The proportion of people who study out of interest increases with age.}\\
静态视角：\en{Older people are more likely to study out of interest than younger people.}\\
\end{eg}

\tip{同一组数据，纵着看是趋势（动态视角），横着看是差异（静态视角）。半动态图（按年龄段、按阶段分组）两种都能写，考场上选你词汇量更大的那一种。}

\impr{（1）By / At / With / To 四条规则、六条比较表达、四条倍数写法、三条时间表达、四类年龄表达，全部来自御风 Q3 与“特殊图表”部分，已逐条配正误对照；（2）“By + 过去时间用过去完成”“可数名词用 fewer”“between \ldots\ and 不搭 to”“in a row / for consecutive years 的形容词形式”等，是原书只写了词、没有点明语法的部分，本次补齐；（3）倍数一表与上一小节第九表是“词汇”与“写法”的分工，不重复内容；（4）“年龄表达”与“半动态图三视角”取自御风“一些特殊图表题的处理方法”一节（源文件第 728--742 行）；第 5 章讲半动态图写法时也给出了同一组年龄表达，两处内容一致，任选一处记即可。}

\begin{tpl}{本章背诵优先级（本次整理）}
\begin{enumerate}[leftmargin=1.6em,itemsep=2pt,topsep=0pt,parsep=2pt]
\item \textbf{第一优先}：15 类衔接词，每类只挑 3--5 个自己一定会用的，其余划掉。这是 CC 提分最快、成本最低的一块。
\item \textbf{第二优先}：趋势词与程度词（第 4 小节），必须成对记——soar 配 dramatically，rise 配 slightly；描述段每个句子都要用到。
\item \textbf{第三优先}：利弊短语（第 3 小节）与上位概念表（第 5 小节），第二段的分析就靠它们撑住。
\item \textbf{最后}：十种句式。它服务于准确性，不是加分装饰；一篇作文用 2--3 种即可，用错不如不用。
\end{enumerate}
\end{tpl}

\section{本次整合做了哪些改进}
\markboth{本次整合做了哪些改进}{}

\begin{note}
这一章把两本模板的原文与改后版本放在一起对照，方便你判断“要不要改、为什么改”。\textbf{凡涉及语法和搭配的改动，都是原文确实有误；凡涉及内容和篇幅的改动，依据都写明了出处。}如果你觉得某处不必改，至少先看完理由。
\end{note}

\subsection{事实与数据准确性}

\begin{longtable}{C{0.8cm} L{4.2cm} L{4.4cm} L{4.6cm}}
\toprule
& \textbf{原模板写法} & \textbf{问题} & \textbf{改进后 / 依据}\\
\midrule
\endfirsthead
\toprule
& \textbf{原模板写法} & \textbf{问题} & \textbf{改进后 / 依据}\\
\midrule
\endhead
1 & 源范文写 \en{... witnessed a sharp rise from 8.8\% to 25.4\%, almost tripling in a mere decade.} & 倍数表述本身没有问题（8.8\% $\rightarrow$ 25.4\% 约为 2.9 倍，\en{almost tripling} 是准确的）。但本册把它改成 \en{rose sharply ... almost tripling}，是为了让\textbf{主语统一在 “proportion” 上}，避免 \en{witness a rise} 这种“数据自己见证变化”的拟人化表达在图表作文里反复出现。 & 保留 \en{almost tripling} 的说法；同时提醒：倍数、幅度的表述必须与数据吻合，\en{tripling} 不能用在只涨了一倍的数据上。依据：批改原则第 2 条。\\
\midrule
2 & 模板中大量 \fillin{百分比 1} 占位符 & 学生容易把它当成“装饰”随手填一个数 & 第 1 章增设\textbf{红线三}：所有数字必须来自图表；第 3 章给出“数字总数 3--5 个”的量化约束。\\
\midrule
3 & 未区分“百分比”与“百分点” & 极易出现 \en{rose by 33\%}（应为 33 percentage points）这类错误 & 第 1 章\textbf{红线二}专门讲；第 4 章与第 13 章给出正确/错误对照例句。依据：批改原则第 2 条明确点出“不混淆百分比与百分点”。\\
\midrule
4 & mmg 动态图第一段模板末句：“\en{These trends reflect not only changes in quantity, but also changes in the influence of X on daily life.}” & 该句本身可用于动态图，但“对日常生活的影响变化”超出了图表能证明的范围，属于把推测写成事实；学生若把它挪用到静态图，还会凭空造出“trend” & 本册把这类句子统一限定在动态图使用，并给“影响类”表述加了分寸词（\en{may / suggests}）。依据：批改原则第 2、3 条。\\
\midrule
5 & 半动态图（年龄段）中未提醒“视角不可混用” & 学生容易把“类别差异”写成“随时间变化” & 第 5 章明确给出“动态视角 / 静态视角”二选一，并给出禁止混用的说明。\\
\bottomrule
\end{longtable}

\subsection{语法与搭配纠错}

下面每一条都是两本模板原文中\textbf{确实存在的语言错误}（不含 OCR 识别错误），已在本册全部改正。

\begin{longtable}{C{0.8cm} L{5.6cm} L{3.2cm} L{5.0cm}}
\toprule
& \textbf{原文} & \textbf{问题} & \textbf{改进后}\\
\midrule
\endfirsthead
\toprule
& \textbf{原文} & \textbf{问题} & \textbf{改进后}\\
\midrule
\endhead
1 & \en{While other activities ... are overshadowed in comparison, none of which exceeds 35\%.} & \en{While} 引导的从句后面直接接了 \en{none of which} 这个关系从句，\en{while} 从句因此没有主句可依附（\en{none of which exceeds} 单独使用本身没有问题，问题出在 \en{While} 结构） & \en{... are far less popular, none of them exceeding 35\%.}\\
\midrule
2 & \en{Every little endeavor counts to enhance the quality of life.} & \en{count to do} 不是“有助于” & \en{counts toward enhancing ...}\\
\midrule
3 & \en{take more participation in volunteer work} & 搭配不存在 & \en{take a more active part in voluntary work}\\
\midrule
4 & \en{The participation rate highlights the prevalence of ...} & 主语与动词搭配生硬 & \en{Watching TV and walking stand out as the most common choices ...}\\
\midrule
5 & \en{abundant enlightenments} & \en{enlightenment} 通常不可数 & \en{abundant insights}\\
\midrule
6 & \en{As a prescription, it is imperative ...} & \en{as a prescription} 不地道 & \en{As a remedy, it is imperative ...} 或直接删去\\
\midrule
7 & \en{the burden of raising their parents and minor kids} & \en{raise} 是“抚养”而非“赡养”；\en{minor kids} 不地道 & \en{support both their aging parents and their young children}\\
\midrule
8 & \en{no longer being aggressive and motivated, they feel worn out} & 悬垂结构，\en{being} 的逻辑主语与主句不一致 & \en{... while coping with declining energy}\\
\midrule
9 & \en{On one hand, ...} & 固定搭配是 \en{On the one hand} & \en{On the one hand, ...}\\
\midrule
10 & \en{Too familiar as such scenario may seem} & 倒装条件不成立；\en{scenario} 缺冠词 & \en{Familiar as such a scenario may seem, ...}\\
\midrule
11 & \en{While society rapidly varies} & \en{vary} 表示“各不相同”，不能表示“变化快” & \en{While society changes rapidly, ...}\\
\midrule
12 & \en{Only if personal ambition and societal needs are closely combined, and simultaneously integrating oneself into ...} & \en{Only if} 从句与 \en{integrating} 分词并列，结构不对称 & 拆成两个独立句；或 \en{Only when personal ambition is combined with societal needs can we ...}\\
\midrule
13 & \en{not only does ... suggest achievements in the past, but acts as a promising start} & \en{not only ... but} 前后动词形式不平衡 & 改为两个简单句\\
\midrule
14 & \en{the figures will not stagnate but scale new peaks} & \en{figures ... scale peaks} 搭配不自然 & \en{turn awareness into everyday practice}\\
\midrule
15 & \en{Words fail to express my excitement to inform you about ...} & 结构不成立：\en{express my excitement to inform} 逻辑不通 & \en{Words cannot express how delighted I was to learn that ...} 或直接 \en{I am writing to ...}\\
\midrule
16 & \en{in the light of 题目相关词}（被用来表示“结合”） & \en{in the light of} 意为“考虑到、根据”，不表示“结合” & 统一改为 \en{with a focus on ...}；表示“在……背景下”时用 \en{in the context of ...}\\
\midrule
17 & \en{around 题目相关词}（表“围绕……”） & \en{around} 这样用不成搭配 & 统一改为 \en{with a focus on ...} / \en{in the area of ...}\\
\bottomrule
\end{longtable}

\subsection{与评分规则冲突之处}

这是本次整合中\textbf{最重要}的一类改动：原模板本身没错，但按英语二的评分规则使用会掉分。

\begin{longtable}{C{0.8cm} L{4.6cm} L{4.4cm} L{4.6cm}}
\toprule
& \textbf{原模板做法} & \textbf{与哪条规则冲突} & \textbf{改进方案}\\
\midrule
\endfirsthead
\toprule
& \textbf{原模板做法} & \textbf{与哪条规则冲突} & \textbf{改进方案}\\
\midrule
\endhead
1 & 讲“英语一和英语二的区别只有第一段，第二段第三段通用”，鼓励直接套用经济类/品质类段落 & 批改原则第 3 条：“评论须与图表相关、逻辑合理” & 第 6 章为每一类语料\textbf{补上锚定句}，并要求把 \fillin{题目相关词} 换成图表里的真实对象\\
\midrule
2 & “体育健康类（如果背了环境生态，直接就能用到体育健康）” & 直接套用会出现 \en{green development} 出现在健康题里的语域错配 & 第 6 章给出三类之间的\textbf{逐词替换表}（\en{green development → healthy living} 等）\\
\midrule
3 & 建议大作文 180--220 词 & 批改原则第 4 条：“题目要求约 150 词时，不擅自改成必须不少于 150 词” & 本册大作文目标定为 \textbf{150--175 词}，并在第 2 章给出逐段词数预算\\
\midrule
4 & 建议小作文 110--120 词 & 大纲为“约 100 词” & 本册定为 \textbf{100--120 词}，并说明称呼与落款不计入\\
\midrule
5 & 大量使用 \en{Succinct as it is, the chart conveys abundant enlightenments, and below are my perspectives.} 等“万能过渡句” & 批改原则第 5 条虽说不因常见句式扣分，但档位描述中的“内容切题”要求内容与题目相关；过渡句过滥会强化模板感 & 第 3、4 章将过渡句改为“自写具体句或直接省略”，并说明理由\\
\midrule
6 & 图画类第三段模板以 \en{In conclusion, ...} 开头（图表类模板本身已用 \en{Overall}） & 图表题只有 150 词，用 \en{In conclusion} 收尾显得笨重；而且同一本材料里两套模板的段首词不统一，容易让学生记混 & 第 7 章把图表题第三段的段首词统一为 \en{Overall}，并说明御风原书也明确建议图表题用 \en{Overall}\\
\midrule
7 & 第三段可能残留数字 & 图表题的第三段应回到观点 & 第 7 章把“不出现数字”列为硬规则并写进自查清单\\
\midrule
8 & 模板的议论句全部是“无 hedge 的断言”（\en{... reflects ...}、\en{... stems from ...}），学生照抄时很容易顺势写成 \en{The chart proves that ...} & 批改原则第 3 条：“合理推测不能写成图表已经证明的事实” & 第 6 章统一采用 \en{may / might / seems to / serves as an implication that}，并在章末提醒：\en{reflects / stems from} 表达的是“你的判断”，不是图表已证明的事实\\
\midrule
9 & 推荐“用豆包之类的 AI 帮你换词” & 依赖外部工具，考场上无法使用 & 第 6 章给出\textbf{可直接查表完成的替换方案}（上位概念表 + 替换表），把“换词能力”变成“查表能力”\\
\midrule
10 & 英语一的图画作文内容与英语二内容混排在同一份材料里 & 考场上容易拿错段落 & 本册\textbf{主体全部为英语二内容}；图画与材料题只作附录 B 的最低优先级保险，并明确标注\\
\bottomrule
\end{longtable}

\subsection{篇幅与结构上的重组}

\begin{longtable}{C{0.8cm} L{5.6cm} L{4.2cm} L{4.4cm}}
\toprule
& \textbf{原材料的组织方式} & \textbf{问题} & \textbf{本册的组织方式}\\
\midrule
\endfirsthead
\toprule
& \textbf{原材料的组织方式} & \textbf{问题} & \textbf{本册的组织方式}\\
\midrule
\endhead
1 & mmg 第三段按“经济类 / 正面品质 / 反面品质 / 校园类 / 环境生态类 / 体育健康类 / 文化自信类”各给一整段，共 8 段，四类通用尾段几乎逐字相同 & 大量重复，实际内容只需 1/3 & 第 6 章整合为\textbf{六类语料 + 逐类锚定改造}；第 7 章把 8 个尾段合并为\textbf{一套骨架 + 一张替换表}\\
\midrule
2 & 御风把“图表作文”与“图画作文”并列讲解，范文中英语一、英语二混排 & 英语二考生辨认成本高 & 第 3--8 章全部为英语二图表内容；英语一范文只保留有借鉴价值的段落并明确标注\\
\midrule
3 & 两个材料的“逻辑衔接词”部分合计 160 余条，平铺罗列，未按使用频率排序 & 记不住，也容易用错 & 第 13 章按 15 类完整收录，但每类标注“精记 2--5 个”，并配英语二示范半句\\
\midrule
4 & 两个材料都没有“自查”环节 & 学生写完不知道好不好 & 新增第 3、4、6、7 章的\textbf{分段自查清单}和第 8 章的\textbf{全篇自查清单}\\
\midrule
5 & 两个材料都没有把评分档位与写作动作对应起来 & 学生只知道“要写好”，不知道“好”在哪几条 & 新增第 1 章的\textbf{档位--行动对照表}\\
\midrule
6 & 两个材料虽然分别有“静态第一段”“动态第一段”小节，但都没有告诉学生\textbf{拿到一张图怎么判断属于哪一种} & 判断错就会把“类别差异”写成“时间趋势”，是掉档最快的错误 & 新增第 2 章的\textbf{判断树}，并把静态与动态分别成章（第 3、4 章）\\
\midrule
7 & 没有单独的“数字表达”小节（御风有，但放在 Q\&A 里） & 容易漏 & 第 13 章单列“数字表达速查”\\
\midrule
8 & mmg 缺少“科技与网络”话题语料 & 而这是英语二的高频方向 & 第 6 章\textbf{新增第六类：科技与网络类}完整段落\\
\bottomrule
\end{longtable}

\subsection{本册新增的内容（原材料没有的）}

\begin{longtable}{C{0.9cm} L{5.4cm} L{8.0cm}}
\toprule
& \textbf{新增内容} & \textbf{解决什么问题}\\
\midrule
\endfirsthead
\toprule
& \textbf{新增内容} & \textbf{解决什么问题}\\
\midrule
\endhead
1 & 全书结构图（首页后） & 一页看清两篇文章的段落构成与评分关注点\\
\midrule
2 & 英语二大作文 / 小作文评分档位表（含“这一档在阅卷现场的真正含义”） & 把官方档位翻译成可执行的动作\\
\midrule
3 & 图表作文四条红线 & 把批改原则里最容易踩的坑前置，避免从 13--15 档掉档\\
\midrule
4 & 档位--行动对照表 & 知道自己目标分对应要做的具体事情\\
\midrule
5 & 审题七问 & 90 秒完成审题，避免动笔就跑偏\\
\midrule
6 & 静态/动态判断树（含半动态） & 从源头杜绝“把类别差异写成时间趋势”\\
\midrule
7 & 九句骨架 & 把全篇拆成 9 个功能位，结构与内容分离\\
\midrule
8 & 逐段字数预算 & 把“约 150 词”变成可执行的分段目标\\
\midrule
9 & 每一类语料的“锚定句” & 让模板议论有图表依据，消除模板感\\
\midrule
10 & 图表锚定的六类话题语料（含新增的科技与网络类） & 覆盖英语二全部高频方向\\
\midrule
11 & 上位概念表 & 用最少的词覆盖最多的题，统一语域\\
\midrule
12 & 四套真题的批注版范文（含修改理由） & 看到“模板如何落地成文章”\\
\midrule
13 & 分段自查清单 + 全篇自查清单 & 写完能自己判分、自己纠错\\
\midrule
14 & 本改进清单 & 让每一处改动可追溯\\
\midrule
15 & 附录 B：图画 / 材料题应急四句 & 应对大纲允许但未曾出现的题型，成本极低\\
\bottomrule
\end{longtable}

\clearpage

\appendix
\renewcommand{\thesection}{附录 \Alph{section}}
\section{英语二大作文历年真题速查}
\markboth{附录 A 英语二大作文历年真题速查}{}

\begin{note}
用法：把每一年的“图表类型 + 静态/动态”标出来，然后\textbf{只看题目，不看范文}，在草稿纸上写出第 1 段的三句话和三个数字。写不出来的年份，就是你需要回头复习的年份。
\end{note}

\begin{longtable}{C{1.4cm} C{2.4cm} C{2.0cm} L{9.4cm}}
\toprule
\textbf{年份} & \textbf{图表类型} & \textbf{静/动} & \textbf{主题}\\
\midrule
\endfirsthead
\toprule
\textbf{年份} & \textbf{图表类型} & \textbf{静/动} & \textbf{主题}\\
\midrule
\endhead
2010 & 柱状图 & 动态 & 发展中国家与发达国家手机用户数量的变化\\
\midrule
2011 & 柱状图 & 动态 & 不同品牌汽车市场份额的变化\\
\midrule
2012 & 表格 & 半动态 & 某公司不同年龄段员工工作满意度调查\\
\midrule
2013 & 柱状图 & 静态 & 不同年级学生兼职情况\\
\midrule
2014 & 柱状图 & 动态 & 城乡人口数量的变化\\
\midrule
2015 & 饼状图 & 静态 & 居民春节假期花销比例\\
\midrule
2016 & 饼状图 & 静态 & 大学生旅游目的调查\\
\midrule
2017 & 折线图 & 动态 & 博物馆数量与参观人数的变化\\
\midrule
2018 & 饼状图 & 静态 & 消费者选择餐厅时关注的因素\\
\midrule
2019 & 柱状图 & 静态 & 本科毕业生去向统计\\
\midrule
2020 & 饼状图 & 静态 & 大学生手机阅读目的调查\\
\midrule
2021 & 柱状图 & 静态 & 居民体育锻炼方式调查\\
\midrule
2022 & 柱状图 & 动态 & 我国快递业务量的变动情况\\
\midrule
2023 & 折线图 & 动态 & 我国居民健康素养水平的变化\\
\midrule
2024 & 柱状图 & 静态 & 劳动实践课学生收获调查\\
\midrule
2025 & 柱状图 & 静态 & 社区老年人休闲活动调查\\
\bottomrule
\end{longtable}

\begin{note}
\textbf{三个统计规律：}
① \textbf{柱状图占绝大多数}（9 次），饼状图 4 次，折线图 2 次，表格 1 次。柱状图既可静态也可动态，务必两种写法都练熟。
② \textbf{静态题略多于动态题}（约 9 : 6）。很多同学只练动态，结果遇到 2024、2025 这种纯静态题反而不会写。
③ \textbf{主题集中在“民生改善”上}：消费、教育、健康、文化、老龄、就业。所以第 6 章的“上位概念表”覆盖度很高。
\end{note}

\clearpage
\section{应急：若 Part B 改为图画或材料题}
\markboth{附录 B 应急：若 Part B 改为图画或材料题}{}

\begin{note}
\textbf{先看清楚这段说明再决定要不要背：}
英语二自 2010 年分卷以来，Part B \textbf{从未考过图画或文字材料题}。大纲虽然把“图画”列为可能的提示形式，但命题实践一直稳定在图表上。
近年英语一出现了“漫画 + 图表”的混合题（2024）和纯表格题（2025），说明 IELTS/图表化的趋势在增强，但对英语二而言，\textbf{图表仍是唯一需要认真准备的题型}。
本附录只给\textbf{一段四句话}，用于极端情况下的保险。\textbf{它的优先级最低}——如果你时间紧张，看完第 3--8 章就够了。
\end{note}

\begin{tpl}{图画题第 1 段（四句话，背下来即可）}
\en{This picture centers on} \fillin{题目相关词}\en{, and its simple composition is nevertheless rich in meaning.} \fillin{一句话描述图画内容}\en{. Although the picture involves nothing more than a small detail from everyday life, what it reflects is not merely an isolated phenomenon, but rather a social issue of a certain universality. In other words, it is precisely through this way of illustrating a larger issue through a small detail that the author endows the picture with greater practical significance.}
\end{tpl}

\begin{note}
\textbf{关键提示（来自 mmg 原书）：}
\begin{enumerate}[leftmargin=1.7em,itemsep=1.5pt,topsep=2pt]
\item \fillin{一句话描述图画内容} \textbf{不要翻译这几个字}，要用英文自己写出图画内容（谁、在做什么、说了什么）。
\item 图画的第 \textbf{2、3 段与图表题完全相同}——直接用第 6 章的话题语料和第 7 章的建议功能句即可。这是 mmg 原书的结论，也是符合评分逻辑的（\en{interpret} 与 \en{comment} 的任务没变）。
\item 如果题目给的是\textbf{一段文字材料}，把 \en{this picture} 换成 \en{this passage / this material} 即可，句式不变。
\end{enumerate}
\end{note}

\begin{tpl}{如果图画信息量很小（无文字说明），改用这个开头}
\en{The cartoon presents an encouraging / disconcerting scene. In the picture,} \fillin{一到两句描述图画内容}\en{. Plain as the picture may appear, it points beyond the surface scene and serves as a broader reflection of} \fillin{中心词}\en{.}
\end{tpl}

\begin{note}
\textbf{最后的提醒：}附录 B 存在的唯一理由是“万一”。请把 99\% 的精力放在图表作文上——因为那是 16 年来 100\% 会出现的题型。
\end{note}

\clearpage
\section{一页速查卡（建议打印）}
\markboth{附录 C 一页速查卡}{}

\begin{tpl}{Part B 图表作文 · 九句骨架}
\textbf{第 1 段}\\
\en{1.} \en{The} \fillin{bar chart / pie chart / line graph / table} \en{illustrates} \fillin{指标} \en{of} \fillin{对象} \en{among} \fillin{范围}\en{.}\\
\en{2.} \en{Overall,} \fillin{极值/趋势}\en{, accounting for / rising from X to Y}\en{.}\\
\en{3.} \en{In contrast,} \fillin{对比项}\en{, none of them exceeding} \fillin{数值}\en{.}\\[3pt]
\textbf{第 2 段}\\
\en{4.} \fillin{锚定句：指向图表核心特征}\\
\en{5.} \en{For one thing, what first comes into view, without any doubt, is} \fillin{分点一}\en{.}\\
\en{6.} \en{Thanks to} \fillin{原因}\en{,} \fillin{解释}\en{.}\\
\en{7.} \en{For another, there still remain potential concerns indicated in the chart.} \fillin{现象} \en{serves as an implication that} \fillin{分点二}\en{.}\\[3pt]
\textbf{第 3 段}\\
\en{8.} \en{Overall, the chart reflects both} \fillin{正面} \en{and} \fillin{隐忧}\en{.}\\
\en{9.} \fillin{判断} \en{+} \fillin{具体建议（执行者明确）}\en{.}
\end{tpl}

\begin{tpl}{Part A 小作文 · 三段公式}
\en{1.} \en{Dear} \fillin{称呼}\en{,}\\
\en{2.} \en{I hope this message finds you well. / It is a great pleasure for me to write to you.}\\
\en{3.} \en{I am writing to} \fillin{invite you to / express my sincere gratitude for / offer some suggestions on / lodge a complaint about} \fillin{事由}\en{.}\\
\en{4.} \fillin{第二段：3--5 句，逐点回应题干要点（时间/地点/内容/意义）}\\
\en{5.} \fillin{第三段：1--2 句，期待回应}\\
\en{6.} \en{Yours sincerely, / Li Ming}
\end{tpl}

\begin{tpl}{动笔前 60 秒的五个问题}
1. 图表类型是什么？\quad 2. 静态还是动态？\quad 3. 时态？\quad 4. 最突出的特征是哪个？\quad 5. 第 2 段我要说哪两点？
\end{tpl}

\begin{tpl}{交卷前 120 秒的六查}
\cb\ 静态/动态有没有混写？\quad \cb\ 数字是否都来自图表？\quad \cb\ 百分比和百分点有没有混？\\
\cb\ 第 3 段有没有残留数字？\quad \cb\ 小作文要点是否逐条回应？\quad \cb\ 称呼/结束语/署名是否完整？
\end{tpl}

\vspace{6pt}
\begin{center}
\rule{0.6\linewidth}{0.6pt}\\[6pt]
{\small\color{mid} —— 全册完 ——}
\end{center}

\end{document}
